% ============================================================
%  300,000 Streets --- Technical Implementation Book
%  Part 1: Chapters 1--4
%  Companion to the final project report.
%  Audience: developer who wants to recreate the project from scratch.
% ============================================================
\documentclass[11pt,a4paper,openany]{book}

% -- Encoding & fonts -----------------------------------------
\usepackage[T1]{fontenc}
\usepackage[utf8]{inputenc}
\usepackage{lmodern}
\usepackage{microtype}

% -- Page geometry ---------------------------------------------
\usepackage[
  a4paper,
  top=25mm, bottom=25mm,
  left=28mm, right=28mm,
  headheight=14pt
]{geometry}

% -- Colour ----------------------------------------------------
\usepackage[dvipsnames,table]{xcolor}
\definecolor{regengreen}{RGB}{46,139,87}
\definecolor{navyblue}{RGB}{10,35,66}
\definecolor{codeblue}{RGB}{0,0,180}
\definecolor{codegray}{RGB}{128,128,128}
\definecolor{codebg}{RGB}{248,248,248}
\definecolor{amber}{RGB}{232,168,56}

% -- Headers & footers -----------------------------------------
\usepackage{fancyhdr}
\pagestyle{fancy}
\fancyhf{}
\fancyhead[LE]{\small\textcolor{navyblue}{\nouppercase{\leftmark}}}
\fancyhead[RO]{\small\textcolor{navyblue}{\nouppercase{\rightmark}}}
\fancyfoot[C]{\small\thepage}
\fancyfoot[R]{\small\textcolor{codegray}{300,000 Streets --- Technical Book}}
\renewcommand{\headrulewidth}{0.4pt}
\renewcommand{\footrulewidth}{0.2pt}

% -- Tables ----------------------------------------------------
\usepackage{booktabs}
\usepackage{array}
\usepackage{longtable}
\usepackage{tabularx}
\newcolumntype{L}[1]{>{\raggedright\arraybackslash}p{#1}}
\newcolumntype{C}[1]{>{\centering\arraybackslash}p{#1}}

% -- Lists -----------------------------------------------------
\usepackage{enumitem}
\setlist{itemsep=3pt,topsep=4pt}

% -- Maths -----------------------------------------------------
\usepackage{amsmath,amssymb}

% -- Code listings ---------------------------------------------
\usepackage{listings}
\lstdefinestyle{python}{
  language=Python,
  basicstyle=\ttfamily\footnotesize,
  backgroundcolor=\color{codebg},
  keywordstyle=\color{codeblue}\bfseries,
  commentstyle=\color{codegray}\itshape,
  stringstyle=\color{regengreen},
  numberstyle=\tiny\color{codegray},
  numbers=left,
  numbersep=8pt,
  stepnumber=1,
  breaklines=true,
  breakatwhitespace=false,
  frame=single,
  rulecolor=\color{black!20},
  showstringspaces=false,
  tabsize=4,
  captionpos=b,
  morekeywords={True,False,None,as,with,yield,lambda,assert},
}
\lstdefinestyle{bash}{
  language=bash,
  basicstyle=\ttfamily\footnotesize,
  backgroundcolor=\color{black!5},
  keywordstyle=\color{navyblue}\bfseries,
  commentstyle=\color{codegray}\itshape,
  stringstyle=\color{regengreen},
  frame=single,
  rulecolor=\color{black!20},
  breaklines=true,
  showstringspaces=false,
  captionpos=b,
}
\lstdefinestyle{text}{
  basicstyle=\ttfamily\footnotesize,
  backgroundcolor=\color{black!5},
  frame=single,
  rulecolor=\color{black!20},
  breaklines=true,
  captionpos=b,
}
\lstset{style=python}

% -- Graphics --------------------------------------------------
\usepackage{graphicx}
\usepackage{float}
\usepackage{caption}
\captionsetup{font=small,labelfont=bf}

% -- Hyperlinks ------------------------------------------------
\usepackage[hidelinks,colorlinks=true,
  linkcolor=navyblue,urlcolor=regengreen,citecolor=navyblue]{hyperref}

% -- Misc ------------------------------------------------------
\usepackage{parskip}
\setlength{\parindent}{0pt}
\usepackage{tcolorbox}
\tcbuselibrary{breakable,skins}

% -- Callout boxes --------------------------------------------
\newtcolorbox{keybox}[1]{
  colback=navyblue!8,colframe=navyblue,
  fonttitle=\bfseries\small,title={#1},
  breakable,left=6pt,right=6pt,top=4pt,bottom=4pt
}
\newtcolorbox{tipbox}[1]{
  colback=regengreen!8,colframe=regengreen,
  fonttitle=\bfseries\small,title={#1},
  breakable,left=6pt,right=6pt,top=4pt,bottom=4pt
}
\newtcolorbox{warnbox}[1]{
  colback=amber!15,colframe=amber!70!black,
  fonttitle=\bfseries\small,title={#1},
  breakable,left=6pt,right=6pt,top=4pt,bottom=4pt
}

% -- Chapter style ---------------------------------------------
\usepackage{titlesec}
\titleformat{\chapter}[display]
  {\normalfont\Huge\bfseries\color{navyblue}}
  {\chaptertitlename\ \thechapter}{10pt}{\Huge}
\titlespacing{\chapter}{0pt}{20pt}{30pt}

% ============================================================
\begin{document}

% -- Title page -----------------------------------------------
\begin{titlepage}
\centering
\vspace*{2cm}
{\color{navyblue}\rule{\linewidth}{2pt}}\\[0.6cm]
{\Huge\bfseries\color{navyblue}
  300,000 Streets of Melbourne\\[0.3cm]
  Technical Implementation Book}\\[0.5cm]
{\color{regengreen}\rule{\linewidth}{1pt}}\\[0.5cm]
{\Large\itshape From Raw Data to Interactive Dashboard:\\
A Step-by-Step Developer's Guide}\\[1.5cm]
{\large
  Regen Melbourne $\times$ RMIT University\\[0.3cm]
  City of Darebin --- School Streets Safety Analysis\\[1.5cm]
  \textbf{Parts 1--7 + Appendices: Chapters 1--24}\\[0.3cm]
  From Raw Data to Production Pipeline\\[2cm]
}
{\normalsize
  \begin{tabular}{ll}
    \textbf{Author:}              & Hishikesh Phukan \\
    \textbf{Academic supervisor:} & Lawrence (RMIT University) \\
    \textbf{Industry partner:}    & Regen Melbourne \\
    \textbf{Schools assessed:}    & Reservoir HS, William Ruthven SC, Preston HS \\
    \textbf{LGA:}                 & City of Darebin, Victoria \\
    \textbf{Data period:}         & 2021--2025 \\
  \end{tabular}
}\\[2cm]
{\color{navyblue}\rule{\linewidth}{2pt}}\\[0.3cm]
{\small\color{codegray}
  This book is a living technical companion. It documents every design decision,
  data source, algorithm, and implementation detail so that any developer can
  recreate the pipeline from scratch.
}
\end{titlepage}

% -- Front matter ---------------------------------------------
\frontmatter
\tableofcontents
\clearpage

% -- Country Acknowledgement ----------------------------------
\thispagestyle{empty}
\vspace*{\fill}
\begin{center}
{\color{regengreen}\rule{0.6\linewidth}{1.5pt}}\\[1.2cm]
{\large\bfseries\color{navyblue} Acknowledgement of Country}\\[1cm]
\end{center}
\begin{center}
\parbox{0.82\linewidth}{%
\setlength{\parskip}{0.8em}
\large\itshape
I acknowledge the Wurundjeri Woi-wurrung people of the Kulin Nation as the
Traditional Custodians of the land on which this research was conducted ---
the land we call Melbourne (\textup{Naarm}), and the surrounding suburbs of
Darebin, Reservoir, Preston, and Thornbury, all of which sit within
Wurundjeri Woi-wurrung Country.

\medskip
I pay my deepest respects to their Elders, past and present, and to all
Aboriginal and Torres Strait Islander peoples whose connection to Country
predates the streets we study and endures long after the data is gathered.

\medskip
This land was never ceded.

\smallskip
\begin{flushright}
\normalfont\normalsize\color{navyblue}
--- Hishikesh Phukan
\end{flushright}
}
\end{center}
\vspace*{\fill}
{\color{regengreen}\rule{\linewidth}{0.4pt}}
\clearpage

% -- Acknowledgements -----------------------------------------
\chapter*{Acknowledgements}
\addcontentsline{toc}{chapter}{Acknowledgements}

\noindent
This work is the product of extraordinary support from people and institutions across
two continents.  I am deeply grateful to all of them.

\bigskip
\noindent\textbf{Australia --- a country that gave me a chance.}\\[0.3em]
I came to Australia to pursue a Master's degree and left with skills, knowledge,
and an understanding of what data-driven research can do for real communities.
Australia opened that door.  The public institutions that underpin this project ---
the ABS, EPA Victoria, VicRoads, Crime Statistics Agency, OpenStreetMap contributors,
and the many councils whose open data made this analysis possible --- reflect the kind
of society that invests in shared infrastructure and shared knowledge.  I am grateful
to be a small part of that.

\bigskip
\noindent\textbf{RMIT University.}\\[0.3em]
RMIT gave me the academic home, the coursework, and the environment to grow as a
researcher and engineer.  The structure of COSC2667/COSC2777 pushed me to build
something real, not just something theoretical.  I am grateful to my academic
supervisor Lawrence for the guidance, the critique, and the patience throughout
this project.

\bigskip
\noindent\textbf{Regen Melbourne.}\\[0.3em]
The 300,000 Streets initiative is Regen Melbourne's vision: that every child in
metropolitan Melbourne deserves a safe street to walk to school.  Partnering with
Regen Melbourne gave this project a purpose beyond the grade.  Their commitment to
equity-centred urban planning kept me focused on what the data is actually for:
the students walking along Plenty Road every morning.

\bigskip
\noindent\textbf{My parents.}\\[0.3em]
Nothing I have done would have been possible without my parents, who pushed me hard,
believed in me without condition, and supported me through every difficulty of a
Master's degree on the other side of the world.  They did not let me quit when it was
hard, and they celebrated every small win as if it were the biggest.  This book is
for them.

\bigskip\bigskip
\begin{flushright}
\itshape
Hishikesh Phukan\\
Melbourne, Victoria, Australia\\
2026
\end{flushright}

\clearpage

% -- Preface --------------------------------------------------
\chapter*{Preface}
\addcontentsline{toc}{chapter}{Preface}

This book documents the complete technical implementation of the \textbf{300,000 Streets of Melbourne} pilot study --- a data-driven assessment of pedestrian safety conditions around three secondary schools in the City of Darebin. It is intended as a companion to the formal project report and is written so that any developer with intermediate Python experience can recreate every output from scratch.

\textbf{Who should read this book.}
This book is for:
\begin{itemize}
  \item Data scientists and engineers who want to extend the pipeline to more schools or other local government areas (LGAs);
  \item Urban planners and researchers who want to understand the methodology behind the Healthy Streets scoring engine;
  \item Students who want a worked example of an end-to-end spatial data science project, from field survey design through to a deployed web dashboard.
\end{itemize}

\textbf{What this book covers.}
Chapters 1--4 cover the project context, the Healthy Streets framework, environment setup, and the overall software architecture. Later chapters (to be released as Parts 2, 3, and beyond) will cover data collection, scoring, machine learning, scenario analysis, equity analysis, crash analysis, and dashboard deployment.

\textbf{Conventions.}
\begin{itemize}
  \item \textbf{Code listings} show exact file contents from the project. Line numbers align with the actual source files.
  \item \texttt{Monospace text} in the body refers to file names, function names, variable names, or shell commands.
  \item Callout boxes highlight key concepts (\textcolor{navyblue}{\textbf{blue}}), tips (\textcolor{regengreen}{\textbf{green}}), and warnings (\textcolor{amber!70!black}{\textbf{amber}}).
  \item All paths are relative to the project root unless stated otherwise.
\end{itemize}

\mainmatter

% ============================================================
%  PART I --- FOUNDATIONS
% ============================================================
\part{Foundations}

% ============================================================
\chapter{Introduction and Project Context}
% ============================================================

\section{The 300,000 Streets Initiative}

The \textbf{300,000 Streets of Melbourne} initiative is a collaboration between \textbf{Regen Melbourne} --- a not-for-profit urban sustainability organisation --- and \textbf{RMIT University}. Its goal is ambitious: systematically assess pedestrian and cyclist safety conditions around every school in metropolitan Melbourne and use that evidence base to advocate for infrastructure investment in the most disadvantaged communities.

The name derives from an estimate of the total lane-kilometres of local streets in metropolitan Melbourne. Most of these streets have never been assessed for walkability or child pedestrian safety; they lack footpath audits, crossing inventories, or any measure of how welcoming the street environment is for someone arriving on foot or by bicycle.

The pilot study documented in this book focuses on \textbf{three secondary schools in the City of Darebin}, a northern-inner-ring LGA whose school zone population spans a wide socioeconomic gradient. The pilot is designed to prove out the methodology before scaling to the full metropolitan network.

\section{Study Scope and School Selection}

Three schools were selected to represent the socioeconomic diversity of the Darebin LGA, using the ABS 2021 SEIFA Index of Relative Socio-economic Disadvantage (IRSD):

\begin{table}[H]
\centering
\caption{Schools assessed in the pilot study}
\begin{tabular}{L{3.8cm}L{5cm}C{2cm}C{2cm}}
\toprule
\textbf{School} & \textbf{Address} & \textbf{SEIFA Decile} & \textbf{Suburb} \\
\midrule
Reservoir High School & 855 Plenty Rd, Reservoir VIC 3073 & D4 & Reservoir \\
William Ruthven Secondary College & 60 Merrilands Rd, Reservoir VIC 3073 & D6 & Reservoir \\
Preston High School & 2--16 Cooma St, Preston VIC 3072 & D7 & Preston \\
\bottomrule
\end{tabular}
\end{table}

SEIFA decile 1 indicates the most socioeconomically disadvantaged, decile 10 the least. All three schools fall in deciles 4--7, covering the lower-middle to middle range of disadvantage. A key research hypothesis is that disadvantage correlates with worse pedestrian safety --- and the pilot provides early evidence that this is indeed the case (\(r = 0.84\)).

\section{The Seven Deliverables}

The project was contracted around seven concrete deliverables. Every technical component in this book serves at least one of them:

\begin{description}[leftmargin=1.5cm,style=nextline]
  \item[D1 --- Field Assessment Report]
    A written report documenting observations at each school gate, including hazard descriptions, photos, and qualitative commentary. This feeds the raw \texttt{school\_data.csv} dataset.
  \item[D2 --- Healthy Streets Scores]
    Quantitative scores for all ten Healthy Streets indicators (HS1--HS10) at each school, plus an overall severity classification. Output: \texttt{outputs/hs\_scores.csv}.
  \item[D3 --- Visualisation Suite]
    Three static charts --- a radar chart, a per-indicator bar chart, and a per-school breakdown chart. Output: \texttt{outputs/chart1\_safety\_scores.png}, \texttt{chart2\_hazard\_severity.png}, \texttt{chart3\_score\_breakdown.png}.
  \item[D4 --- Equity Analysis]
    An analysis of whether SEIFA disadvantage predicts lower Healthy Streets scores. Output: \texttt{outputs/chart\_equity\_seifa.png}.
  \item[D5 --- Crash Trend Analysis]
    A deep-dive into VicRoads crash data (2021--2025) showing pedestrian and cyclist crash trends, school-hours breakdowns, and time-of-day distributions. Output: \texttt{outputs/chart\_crash\_trends.png}.
  \item[D6 --- Machine Learning Pipeline]
    A Ridge regression model trained on open data (OSM, AQI, crime, crash statistics) to predict HS indicator scores at new schools without field surveys. Output: \texttt{outputs/hs\_predictor.pkl} and associated charts.
  \item[D7 --- Interactive Dashboard]
    A GitHub Pages static web dashboard with scenario explorer, hero stats, and equity/crash charts. Deployed via \texttt{docs/} directory.
\end{description}

\section{How to Use This Book}

This book is \textbf{sequential}. Each chapter builds on the previous one. We recommend working through it in order, running each script as it is introduced and inspecting the outputs before moving on.

A rough reading and execution guide:

\begin{table}[H]
\centering
\caption{Reading and execution guide}
\begin{tabular}{C{1.2cm}L{4.5cm}L{6cm}}
\toprule
\textbf{Part} & \textbf{Chapters} & \textbf{Time estimate} \\
\midrule
1 & Ch.~1--4: Foundations & 2--3 hours (reading + environment setup) \\
2 & Ch.~5--8: Data Collection & 4--6 hours (field survey + API calls can be slow) \\
3 & Ch.~9--12: HS Scoring Engine & 2--3 hours \\
4 & Ch.~13--15: Spatial Feature Engineering & 2--3 hours \\
5 & Ch.~16--18: Machine Learning & 1--2 hours \\
6 & Ch.~19--21: Scenario \& Equity Analysis & 2 hours \\
7 & Ch.~22--24: Visualisation \& Dashboard & 3--4 hours \\
\bottomrule
\end{tabular}
\end{table}

\begin{keybox}{Before you begin}
You will need: Python 3.10 or higher, an internet connection for Overpass API calls to OpenStreetMap, access to the VicRoads open data portal, and approximately 2 GB of free disk space for intermediate outputs. Chapter~3 covers environment setup in full detail.
\end{keybox}

\section{Project History and Timeline}

The project ran over twelve weeks during Semester 1, 2025. The pipeline was built in four sprints:

\begin{enumerate}
  \item \textbf{Sprint 1 (Weeks 1--3):} Project scoping, school selection, field survey protocol design, and initial data collection.
  \item \textbf{Sprint 2 (Weeks 4--6):} HS scoring engine, static visualisations, and crash data integration.
  \item \textbf{Sprint 3 (Weeks 7--9):} OSM spatial feature extraction, feature engineering, and Ridge regression ML model.
  \item \textbf{Sprint 4 (Weeks 10--12):} Scenario engine, equity analysis, interactive dashboard, and report writing. Sprint 4 absorbed significant unplanned scope from upstream delays in OSM data quality issues.
\end{enumerate}

The twelve-week, sequentially dependent pipeline left no buffer for upstream delays --- a key limitation noted in the report. Despite this, all seven deliverables were completed on time.

\section{A Note on Reproducibility}

The OSM data used in this project is live: the Overpass API returns whatever is currently mapped at the time of the query. Our spatial features were extracted in April 2025. If you re-run \texttt{spatial\_features.py} today, OSM may have changed, leading to slightly different feature values and therefore slightly different ML predictions. The VicRoads crash dataset is updated quarterly; similarly, running the crash analysis scripts at a later date will produce updated trend charts.

\begin{warnbox}{Non-determinism}
ML results and some chart values will differ slightly depending on when you pull OSM data. For exact replication of the published results, use the CSV files in \texttt{outputs/} that are committed to the repository rather than re-running the data collection scripts.
\end{warnbox}


% ============================================================
\chapter{The Healthy Streets Framework}
% ============================================================

\section{Origins and Purpose}

The \textbf{Healthy Streets} framework was developed by \textbf{Dr Lucy Saunders} for Transport for London (TfL) in 2015. It provides a structured, evidence-based approach to assessing how well a street supports the physical and mental health of people who walk, cycle, or use public transport. The framework defines ten indicators --- each scored from 0 to 10 --- that together capture whether a street environment is genuinely hospitable to active travel.

The original TfL framework was designed for general urban streets. This project adapts it specifically for the \textbf{school gate context}: a 400-metre radius around the main pedestrian entrance of a secondary school, assessed during the school drop-off and pick-up periods (7:30--9:00 am and 2:30--4:30 pm on weekdays).

\begin{keybox}{Core principle of the Healthy Streets framework}
A ``healthy street'' is one that enables and encourages people to walk, cycle, and use public transport as part of their everyday lives --- not just for recreation, but for the journeys they need to make. The framework asks: would a person of any age, ability, or background choose to walk or cycle on this street?
\end{keybox}

\section{The Ten Indicators}

Each of the ten indicators addresses a distinct dimension of street quality. Together they cover physical infrastructure, environmental conditions, and the subjective experience of being on the street.

\begin{longtable}{C{1.2cm}L{3.5cm}L{8cm}}
\caption{Healthy Streets indicators HS1--HS10}\label{tab:hs_indicators}\\
\toprule
\textbf{Code} & \textbf{Name} & \textbf{Description} \\
\midrule
\endfirsthead
\toprule
\textbf{Code} & \textbf{Name} & \textbf{Description} \\
\midrule
\endhead
\bottomrule
\endfoot
HS1 & Pedestrians from all walks of life &
  Does the street infrastructure enable people of all ages, abilities, and mobility needs to walk safely? This covers footpath presence, width, continuity, condition, kerb ramps, and obstacles. \\[4pt]
HS2 & Easy to cross &
  Can a pedestrian cross the road safely and conveniently at or near the school gate? This covers the presence of a signalised or marked crossing, its distance from the gate, visibility, condition, and whether there are pedestrian countdown timers. \\[4pt]
HS3 & Shade and shelter &
  Is there shade from trees, or shelter from rain, for people waiting near the school gate? Covers tree canopy density within 100 m, bus shelters within 200 m, and green space percentage within 400 m. \\[4pt]
HS4 & Places to stop and rest &
  Are there benches or other seating opportunities within 200 m of the gate? Covers bench count (OSM) and park count within 400 m. \\[4pt]
HS5 & Not too noisy &
  Is the noise environment comfortable for pedestrians? Covers speed limit, traffic volume, heavy vehicle presence, arterial road percentage, and presence of traffic calming measures. \\[4pt]
HS6 & People choose to walk, cycle, or use PT &
  Does the street encourage active and sustainable travel? Covers cycling infrastructure quality (using the Cycling Infrastructure Score, CIS), public transport stop count within 400 m, and presence of a school zone. \\[4pt]
HS7 & People feel safe &
  Do people feel safe from crime and threatening behaviour? Covers suburb-level crime rate (incidents per 100,000 population) from Crime Statistics Victoria. \\[4pt]
HS8 & Things to see and do &
  Is there visual interest and activity along the street? Covers amenity count (cafes, shops, community facilities) and street furniture within 400 m. \\[4pt]
HS9 & People feel relaxed &
  Is the overall street environment calming rather than stressful? Derived from a composite of traffic stress indicators: speed limit, traffic volume, parking conflicts, and lighting quality. \\[4pt]
HS10 & Clean air &
  Is air quality acceptable for sustained outdoor activity? Covers EPA AirWatch annual average PM2.5 concentration for the monitoring station nearest the school. \\
\end{longtable}

\section{Scoring Each Indicator: 0 to 10}

Each indicator is scored on a continuous scale from \textbf{0} (worst) to \textbf{10} (best). The scoring functions are defined in \texttt{src/scoring/} (one Python file per indicator). The following principles apply across all indicators:

\begin{enumerate}
  \item \textbf{Component scores are weighted and summed.} Each indicator draws on 2--8 measurable components. Each component is scored 0--1 (or 0--10 in some cases), then combined with a defined weight.
  \item \textbf{Scores are clipped to [0, 10].} Intermediate calculations can exceed this range; the final score is always clipped.
  \item \textbf{Missing data is handled gracefully.} If a required data field is absent (e.g.\ OSM has no bench tags near a school), the component is scored at 0 or at the dataset mean, depending on the indicator.
  \item \textbf{Spatial data augments field data.} Indicators HS3, HS4, HS6, and HS8 combine field observations with OSM-derived spatial features, because these are difficult to observe completely during a single site visit.
\end{enumerate}

\begin{tipbox}{Which indicators need field surveys?}
HS1, HS2, HS5, HS7, and HS9 rely primarily on field observations (footpath width, crossing visibility, traffic volume, etc.). HS3, HS4, HS6, and HS8 rely primarily on OSM data and can therefore be computed for any school without a site visit, making them key candidates for the ML triage pipeline. HS10 comes from EPA monitoring data and requires no field work at all.
\end{tipbox}

\section{Severity Classification}

After computing all ten indicator scores, the pipeline assigns each school a \textbf{severity classification} --- a single label that prioritises schools for intervention:

\begin{table}[H]
\centering
\caption{Severity classification rules}
\begin{tabular}{L{2.5cm}L{9cm}}
\toprule
\textbf{Severity} & \textbf{Rule} \\
\midrule
\textbf{Major} & HS2 $< 4.0$ \textbf{OR} HS1 $< 4.0$ \textbf{OR} HS5 $< 3.0$ \\
\textbf{Moderate} & Two or more indicators $< 6.0$ \textbf{OR} overall HS score $< 5.0$ \\
\textbf{Minor} & All other cases \\
\bottomrule
\end{tabular}
\end{table}

The thresholds are set deliberately conservative: a school is classified as \textbf{Major} if any single critical indicator --- crossing safety (HS2), footpath accessibility (HS1), or noise (HS5) --- falls below 4.0. This reflects the view that a missing or very poor crossing near a school gate is a stand-alone safety crisis regardless of how well the school performs on other indicators.

In our pilot, Preston HS was classified as \textbf{Major} (HS2 = 0.4, no signalised crossing within 400 m of the gate). Both Reservoir HS and William Ruthven SC were classified as \textbf{Moderate}.

\section{Adapting the Framework for Secondary Schools}

The original TfL framework was designed for general urban contexts, not specifically for schools. Several adaptations were made for this project:

\begin{description}
  \item[School zone recognition (HS6).]
    The presence of an active school zone (as per VicRoads school zone data) is treated as a positive indicator for HS6, because it signals that infrastructure investment has already been made for active travel near the school.
  \item[Gate-based spatial analysis.]
    All buffer radii (200 m, 400 m, 800 m) are measured from the main pedestrian gate of each school, not from the centroid of the school parcel.
  \item[School-hours crash filtering (D5).]
    The crash analysis (Chapter 21 in the full book) filters crashes by time of day: 7:30--9:00 am and 2:30--4:30 pm, Monday--Friday, to focus on the highest-risk periods for students.
  \item[SEIFA equity layer (D4).]
    The equity analysis overlays SEIFA disadvantage on top of HS scores to test whether disadvantaged school catchments also have the worst pedestrian safety --- a finding with direct policy implications.
\end{description}

\section{The Cycling Infrastructure Score (CIS)}

HS6 uses a custom ordinal scale called the \textbf{Cycling Infrastructure Score (CIS)} to encode the type of cycling infrastructure present near each school. The CIS is derived from the Level of Traffic Stress (LTS) framework (Mekuria, Furth \& Nixon, 2012) and from VicRoads Traffic Engineering Manual Volume 3, Part 218.

\begin{table}[H]
\centering
\caption{Cycling Infrastructure Score (CIS) mapping}
\begin{tabular}{L{7cm}C{2.5cm}}
\toprule
\textbf{Infrastructure type observed} & \textbf{CIS score} \\
\midrule
Separated bike lane (physical barrier) & 9.0 \\
Shared path or greenway & 8.0 \\
Painted bike lane (on-road, no barrier) & 4.5 \\
Advisory lane or shared road & 2.0 \\
No cycling infrastructure & 1.0 \\
\bottomrule
\end{tabular}
\end{table}

This mapping is defined in \texttt{config.py} as \texttt{CIS\_MAP} and is imported by the HS6 scoring function.

\section{Why Ten Indicators?}

The ten-indicator structure allows the framework to be \textit{disaggregated}: a single overall score hides the fact that a school may score well on air quality but very badly on crossing safety. By examining all ten indicators, decision-makers can pinpoint exactly which infrastructure interventions will have the greatest impact on the overall score.

In our ML pipeline (Chapters 16--18 of the full book), each indicator is predicted \textit{independently}, allowing the scenario engine to show how a single intervention --- such as installing a traffic signal --- affects HS2 and HS9 specifically, rather than just changing the aggregate score.


% ============================================================
\chapter{Development Environment Setup}
% ============================================================

\section{Prerequisites}

Before you can run any code in this project, you need the following installed on your system:

\begin{itemize}
  \item \textbf{Python 3.10 or higher.} The project uses f-strings with \texttt{=} (Python 3.8+), \texttt{pathlib.Path} idioms, and type annotations that require at least 3.10.
  \item \textbf{Git.} To clone the repository and track changes.
  \item \textbf{Conda or pip with virtualenv.} We strongly recommend using a Conda environment to avoid dependency conflicts, particularly between \texttt{geopandas}, \texttt{osmnx}, and \texttt{shapely}.
  \item \textbf{Internet access.} The \texttt{spatial\_features.py} script makes live calls to the OpenStreetMap Overpass API. Each school takes 1--2 minutes. Without internet access, you must use the pre-computed CSVs.
  \item \textbf{$\approx$2 GB free disk space.} The intermediate outputs (CSVs, PNGs, HTML files, the GeoPackage) accumulate to around 1.5 GB.
\end{itemize}

\section{Cloning the Repository}

\begin{lstlisting}[style=bash, caption={Clone the project repository}]
# Clone via HTTPS (no SSH key required)
git clone https://github.com/s4111078-sidd/300-000-School-Streets.git

# Navigate into the project root
cd 300-000-School-Streets

# Verify the directory structure
ls -la
\end{lstlisting}

You should see, among other files: \texttt{main.py}, \texttt{config.py}, \texttt{spatial\_features.py}, \texttt{ml\_model.py}, \texttt{feature\_engineering.py}, \texttt{scenario\_analyzer.py}, \texttt{equity\_analysis.py}, \texttt{crash\_trend\_analysis.py}, and a \texttt{src/} directory.

\section{Creating the Conda Environment}

We strongly recommend Conda because it handles the binary dependencies of \texttt{geopandas} (GDAL, GEOS, PROJ) without requiring you to install system-level libraries manually.

\begin{lstlisting}[style=bash, caption={Create and activate the Conda environment}]
# Create a new environment named "streets" with Python 3.11
conda create -n streets python=3.11 -y

# Activate it
conda activate streets

# Install geospatial dependencies first via conda-forge
# (these have pre-compiled binaries --- much easier than pip)
conda install -c conda-forge \
  geopandas \
  osmnx \
  shapely \
  fiona \
  pyproj \
  networkx \
  -y
\end{lstlisting}

\begin{warnbox}{Why conda-forge for geopandas?}
Installing \texttt{geopandas} via \texttt{pip} on Windows or macOS often fails because it requires compiled C libraries (GDAL, GEOS, PROJ). The \texttt{conda-forge} channel provides pre-built binaries for all platforms. Always install \texttt{geopandas} and \texttt{osmnx} via conda before installing the remaining packages via pip.
\end{warnbox}

\section{Installing Python Dependencies}

Once the geospatial stack is installed via Conda, install the remaining dependencies using \texttt{pip}:

\begin{lstlisting}[style=bash, caption={Install remaining Python packages via pip}]
# Install all remaining dependencies
pip install \
  pandas==2.1.4 \
  numpy==1.26.4 \
  matplotlib==3.8.2 \
  seaborn==0.13.2 \
  scikit-learn==1.4.0 \
  scipy==1.12.0 \
  folium==0.15.1 \
  jinja2==3.1.3

# Verify key packages installed correctly
python -c "import geopandas, osmnx, sklearn, folium; print('All OK')"
\end{lstlisting}

\begin{tipbox}{Using pip only (no Conda)}
If you cannot use Conda (e.g.\ on a server), try: \texttt{pip install geopandas} after first running \texttt{pip install GDAL}. On Ubuntu/Debian, run \texttt{sudo apt-get install gdal-bin libgdal-dev python3-gdal} first. On macOS with Homebrew: \texttt{brew install gdal}.
\end{tipbox}

\section{Full Requirements Reference}

The following table lists all Python packages used across the project, their purpose, and the modules that import them:

\begin{longtable}{L{3.5cm}L{2cm}L{7cm}}
\caption{Python package requirements}\label{tab:requirements}\\
\toprule
\textbf{Package} & \textbf{Version} & \textbf{Used by} \\
\midrule
\endfirsthead
\toprule
\textbf{Package} & \textbf{Version} & \textbf{Used by} \\
\midrule
\endhead
\bottomrule
\endfoot
\texttt{geopandas} & 0.14+ & \texttt{spatial\_features.py}, interactive map \\
\texttt{osmnx} & 1.9+ & \texttt{spatial\_features.py} (OSM API calls) \\
\texttt{shapely} & 2.0+ & Geometry operations (buffers, intersections) \\
\texttt{pyproj} & 3.6+ & CRS transformations (WGS84 $\to$ GDA2020) \\
\texttt{networkx} & 3.2+ & Graph analysis (osmnx dependency) \\
\texttt{pandas} & 2.1+ & All scripts (data manipulation) \\
\texttt{numpy} & 1.26+ & All scripts (numerical operations) \\
\texttt{matplotlib} & 3.8+ & \texttt{main.py}, all chart scripts \\
\texttt{seaborn} & 0.13+ & \texttt{ml\_model.py} (correlation heatmap) \\
\texttt{scikit-learn} & 1.4+ & \texttt{ml\_model.py} (Ridge, Pipeline, LOO-CV) \\
\texttt{scipy} & 1.12+ & KDE heatmap, statistical tests \\
\texttt{folium} & 0.15+ & Interactive map (Leaflet.js wrapper) \\
\texttt{jinja2} & 3.1+ & Dashboard HTML templating \\
\texttt{fiona} & 1.9+ & GeoPackage read/write (osmnx dependency) \\
\end{longtable}

\section{Verifying the Installation}

Run the following verification script to confirm that all critical imports work before touching any project code:

\begin{lstlisting}[language=Python, caption={verify\_install.py --- run this first to confirm your environment}]
"""
verify_install.py
Run this before any project script to confirm the environment is correct.
"""
import sys

checks = []

def check(name, fn):
    try:
        fn()
        checks.append((name, True, None))
    except Exception as e:
        checks.append((name, False, str(e)))

check("Python >= 3.10",    lambda: assert_version())
check("pandas",            lambda: __import__("pandas"))
check("numpy",             lambda: __import__("numpy"))
check("matplotlib",        lambda: __import__("matplotlib"))
check("seaborn",           lambda: __import__("seaborn"))
check("sklearn",           lambda: __import__("sklearn"))
check("geopandas",         lambda: __import__("geopandas"))
check("osmnx",             lambda: __import__("osmnx"))
check("shapely",           lambda: __import__("shapely"))
check("pyproj",            lambda: __import__("pyproj"))
check("folium",            lambda: __import__("folium"))

def assert_version():
    v = sys.version_info
    assert v.major == 3 and v.minor >= 10, f"Need Python 3.10+, got {v}"

print(f"\n{'Package':<20} {'Status'}")
print("-" * 35)
all_ok = True
for name, ok, err in checks:
    status = "OK" if ok else f"FAIL: {err}"
    print(f"  {name:<18} {status}")
    if not ok:
        all_ok = False

print()
if all_ok:
    print("  All checks passed. Environment is ready.")
else:
    print("  Some checks failed. See errors above.")
    sys.exit(1)
\end{lstlisting}

Run it with:
\begin{lstlisting}[style=bash]
python verify_install.py
\end{lstlisting}

\section{Setting Up the Outputs Directory}

All scripts write their outputs to the \texttt{outputs/} directory in the project root. This directory is created automatically by each script if it does not exist, but you can create it manually:

\begin{lstlisting}[style=bash]
mkdir -p outputs
\end{lstlisting}

The Git repository is configured to track the \texttt{outputs/} directory but to ignore large binary files. Check \texttt{.gitignore} if you are adding new output types.

\section{Coordinate Reference Systems}

This is one of the most important setup concepts in the entire project. We use two coordinate reference systems (CRS):

\begin{description}
  \item[\texttt{EPSG:4326} --- WGS84 (geographic).]
    Latitude and longitude in decimal degrees. This is the default CRS of GPS receivers, OSM data, and most web APIs. All school gate coordinates in \texttt{config.py} are stored in WGS84.
  \item[\texttt{EPSG:7855} --- GDA2020 / MGA Zone 55 (projected, metric).]
    A Cartesian coordinate system in metres, covering eastern Victoria. This is the Australian standard for precise local-scale spatial analysis. All buffer distances (200 m, 400 m, 800 m) are computed in this CRS.
\end{description}

\begin{warnbox}{CRS mismatch is a silent killer}
If you compute a 400 m buffer in WGS84 degrees rather than EPSG:7855 metres, you will get a buffer of the wrong size --- and it will \textit{not} raise an error. The project uses pyproj's \texttt{Transformer} to explicitly reproject all geometries before every spatial operation. Never skip this step when adapting the code for a new school or LGA.
\end{warnbox}

The \texttt{Transformer} is created once at module load time in \texttt{spatial\_features.py}:

\begin{lstlisting}[language=Python, caption={CRS transformation setup in spatial\_features.py}]
from pyproj import Transformer

CRS_METRIC = 'EPSG:7855'   # GDA2020 / MGA zone 55
CRS_GEO    = 'EPSG:4326'   # WGS84

_transformer = Transformer.from_crs(CRS_GEO, CRS_METRIC, always_xy=True)

def gate_to_projected_point(lat, lon):
    """Convert WGS84 lat/lon to a Shapely Point in EPSG:7855."""
    x, y = _transformer.transform(lon, lat)
    return Point(x, y)
\end{lstlisting}

The \texttt{always\_xy=True} parameter is critical: without it, pyproj may swap latitude and longitude, producing incorrect coordinates.

\section{Running the Pipeline for the First Time}

Once your environment is set up and verified, run the scripts in this order:

\begin{lstlisting}[style=bash, caption={Full pipeline execution sequence}]
# Step 1: Extract OSM spatial features (takes ~5 minutes)
python spatial_features.py

# Step 2: Extract environmental features (AQI + crime)
python environmental_features.py

# Step 3: Download and filter crash data (VicRoads)
python crash_analysis.py

# Step 4: Run the main scoring pipeline
python main.py

# Step 5: Build the ML feature matrix
python feature_engineering.py

# Step 6: Train the ML model (LOO-CV)
python ml_model.py

# Step 7: Equity and SEIFA analysis
python equity_analysis.py

# Step 8: Crash trend analysis charts
python crash_trend_analysis.py

# Step 9: Run a scenario (optional)
python scenario_analyzer.py --school "Preston HS" \
    --interventions pedestrian_crossing speed_reduction
\end{lstlisting}

At the end of Step 4, check that \texttt{outputs/hs\_scores.csv} exists and contains three rows (one per school) with all ten HS indicator columns populated.

\begin{tipbox}{Using pre-computed outputs}
If you want to skip the slow OSM API calls (Steps 1--3), the repository includes pre-computed CSVs in \texttt{outputs/}. You can start directly from Step 4 (\texttt{main.py}) using those cached inputs.
\end{tipbox}

% ============================================================
\chapter{Repository Structure and Architecture}
% ============================================================

\section{Directory Tree}

The project root contains the following structure. Understanding where each file lives is essential before reading the code chapters.

\begin{lstlisting}[style=text, caption={Project directory structure}]
300-000-School-Streets/
|
|-- config.py                  # Shared constants and lookup tables
|-- main.py                    # Master pipeline (Steps 1-6)
|-- spatial_features.py        # OSM feature extraction
|-- feature_engineering.py     # Build ML feature matrix
|-- ml_model.py                # Ridge regression + LOO-CV
|-- scenario_analyzer.py       # What-if scenario CLI tool
|-- equity_analysis.py         # SEIFA equity charts
|-- crash_trend_analysis.py    # VicRoads crash trend charts
|-- environmental_features.py  # AQI + crime data collection
|-- crash_analysis.py          # Raw crash data download + filter
|-- school_data.csv            # Field observation data (raw input)
|-- demographics_darebin.csv   # ABS suburb-level demographics
|
|-- src/                       # Modular source packages
|   |-- data/
|   |   |-- loader.py          # CSV loading + column renaming
|   |-- scoring/
|   |   |-- __init__.py        # Exports all HS compute functions
|   |   |-- hs1.py ... hs10.py # One file per HS indicator
|   |   |-- severity.py        # Severity classification
|   |-- recommendations/
|   |   |-- engine.py          # Rule-based recommendations
|   |-- visualisation/
|   |   |-- heatmap.py         # KDE heatmap (static + interactive)
|   |   |-- interactive_map.py # Leaflet.js interactive map
|   |-- scenarios/
|       |-- engine.py          # Delta-method scenario runner
|       |-- interventions.py   # Intervention definitions
|
|-- outputs/                   # All generated outputs (git-tracked)
|   |-- hs_scores.csv          # School-level HS1-HS10 scores
|   |-- spatial_features.csv   # OSM feature matrix
|   |-- ml_school_features.csv # Combined ML feature matrix
|   |-- ml_predictions.csv     # LOO-CV predictions
|   |-- hs_predictor.pkl       # Trained Ridge model
|   |-- seifa_darebin.csv      # SEIFA scores per school
|   |-- chart*.png             # All static charts
|   |-- map_interactive.html   # Interactive hazard map
|   |-- map_heatmap.html       # KDE heatmap
|   |-- networks.gpkg          # GeoPackage: OSM network layers
|
|-- docs/                      # GitHub Pages dashboard
|   |-- index.html             # Dashboard entry point
|   |-- data/
|       |-- data.json          # Aggregated data for charts
|
|-- report/
|   |-- report.tex             # LaTeX report source
|   |-- book.tex               # This book (LaTeX source)
|
|-- build_presentation.py      # Presentation slide builder
|-- .gitignore
|-- README.md
\end{lstlisting}

\section{The config.py Module}

\texttt{config.py} is the single source of truth for every constant, path, and lookup table in the project. \textit{Every} other module imports from it --- never hardcodes its own copies of school names, gate coordinates, or indicator lists. This ensures consistency across the pipeline.

Here is the complete, annotated source:

\begin{lstlisting}[language=Python, caption={config.py --- full annotated source}]
"""
Configuration -- paths, constants, and lookup tables shared across all
modules. All modules should import constants from here rather than
defining their own.
"""
from pathlib import Path

PROJECT_ROOT = Path(__file__).parent

# -- Paths -------------------------------------------------------------
# All paths are derived from PROJECT_ROOT so the code works regardless
# of where the user clones the repository.
CSV_FILE    = PROJECT_ROOT / 'school_data.csv'
DEMO_FILE   = PROJECT_ROOT / 'demographics_darebin.csv'
OUT_DIR     = PROJECT_ROOT / 'outputs'
SPATIAL_CSV = OUT_DIR / 'spatial_features.csv'
ENV_CSV     = OUT_DIR / 'environmental_features.csv'
HS_CSV      = OUT_DIR / 'hs_scores.csv'

# -- Healthy Streets indicators -----------------------------------------
# 10 indicators (Lucy Saunders / TfL framework, adapted for school
# streets). Each tuple is (code, display_label). The label uses \n
# for multi-line chart labels in matplotlib polar charts.
HS_INDICATORS = [
    ('HS1',  'Pedestrians from\nall walks of life'),
    ('HS2',  'Easy to\ncross'),
    ('HS3',  'Shade and\nshelter'),
    ('HS4',  'Places to stop\nand rest'),
    ('HS5',  'Not too\nnoisy'),
    ('HS6',  'People choose to\nwalk / cycle / PT'),
    ('HS7',  'People\nfeel safe'),
    ('HS8',  'Things to\nsee and do'),
    ('HS9',  'People\nfeel relaxed'),
    ('HS10', 'Clean\nair'),
]
HS_CODES  = [h[0] for h in HS_INDICATORS]   # ['HS1', ..., 'HS10']
HS_LABELS = [h[1] for h in HS_INDICATORS]   # display labels

# -- School gates ------------------------------------------------------
# Main gate coordinates for all assessed schools, keyed by full name.
# These are used when we need the address for display purposes.
SCHOOL_GATES = {
    'Reservoir High School': {
        'lat': -37.7224, 'lon': 145.0294,
        'addr': '855 Plenty Rd, Reservoir VIC 3073',
    },
    'William Ruthven Secondary College': {
        'lat': -37.69654, 'lon': 145.00299,
        'addr': '60 Merrilands Rd, Reservoir VIC 3073',
    },
    'Preston High School': {
        'lat': -37.7417, 'lon': 145.0071,
        'addr': '2-16 Cooma St, Preston VIC 3072',
    },
}

# Full name -> short display name (chart labels, CSV matching).
# Charts use short names to avoid label overflow.
SCHOOL_SHORT_NAMES = {
    'Reservoir High School':             'Reservoir HS',
    'William Ruthven Secondary College': 'William Ruthven SC',
    'Preston High School':               'Preston HS',
}

# Short name -> gate coords.
# Used by crash_analysis and spatial_features which key their gate
# dicts by short name rather than full name.
SCHOOL_GATES_BY_SHORT = {
    'Reservoir HS':       {'lat': -37.7224,  'lon': 145.0294},
    'William Ruthven SC': {'lat': -37.69654, 'lon': 145.00299},
    'Preston HS':         {'lat': -37.7417,  'lon': 145.0071},
}

# -- Cycling Infrastructure Score mapping ------------------------------
# Reference: LTS (Mekuria, Furth & Nixon 2012);
#            VicRoads TEM Vol. 3 Part 218.
CIS_MAP = {
    'Yes -- separated bike lane':          9.0,
    'Yes -- shared path or greenway':      8.0,
    'Yes -- painted bike lane (on-road)':  4.5,
    'Yes -- advisory lane / shared road':  2.0,
    'No cycling infrastructure':           1.0,
}

# -- Column rename mapping ---------------------------------------------
# Maps raw CSV header strings to internal short names. The raw headers
# are the full questions from the Google Forms field survey. The short
# names are what every other module uses internally.
COLUMN_RENAME = {
    'School name'                      : 'School',
    'Overall hazard severity ...'      : 'Severity',
    'Footpath Accessibility Score ...' : 'FAS',
    'Crossing Safety Score ...'        : 'CSS',
    # ... (see full file for all 35 mappings)
}
\end{lstlisting}

\begin{keybox}{Design principle: single source of truth}
If you need to add a new school to the pipeline, you only edit \texttt{config.py}: add an entry to \texttt{SCHOOL\_GATES}, \texttt{SCHOOL\_SHORT\_NAMES}, and \texttt{SCHOOL\_GATES\_BY\_SHORT}. Every other module will automatically pick up the new school on the next run.
\end{keybox}

\section{Data Flow Architecture}

Figure~\ref{fig:dataflow} describes the end-to-end data flow. Data moves from left (raw inputs) to right (published outputs) through five transformation stages.

\begin{figure}[H]
\centering
\begin{tcolorbox}[colback=white,colframe=navyblue,title=\textbf{Data Flow Diagram},
  fonttitle=\bfseries, width=\textwidth]
\begin{lstlisting}[style=text, frame=none, numbers=none]
RAW INPUTS                 PROCESSING                  OUTPUTS
==========                 ==========                  =======

school_data.csv  --------> src/data/loader.py
                |                |
                |                v
                |          src/scoring/           --> hs_scores.csv
                |          (HS1 - HS10)           --> chart1,2,3 PNGs
                |
                |
OSM Overpass -->|--> spatial_features.py         --> spatial_features.csv
                |           |
VicRoads ------>|--> crash_analysis.py           --> crash_data_*.csv
                |           |
EPA AirWatch -->|--> environmental_features.py   --> environmental_features.csv
Crime Stats ---->           |
                            |
                            v
                      feature_engineering.py      --> ml_school_features.csv
                            |
                            v
                        ml_model.py               --> hs_predictor.pkl
                        (Ridge LOO-CV)            --> chart_hs_*.png
                            |
                            v
                    scenario_analyzer.py          --> scenario_*.png
                    equity_analysis.py            --> chart_equity_seifa.png
                    crash_trend_analysis.py       --> chart_crash_trends.png
                            |
                            v
                      Dashboard (docs/)           --> GitHub Pages
\end{lstlisting}
\end{tcolorbox}
\caption{End-to-end data flow: from field survey and open APIs to published dashboard}
\label{fig:dataflow}
\end{figure}

\section{The Six-Step Main Pipeline}

The \texttt{main.py} script orchestrates the six core steps of the Healthy Streets scoring pipeline. Understanding this script is essential because it ties together all the modules in \texttt{src/}:

\begin{lstlisting}[language=Python, caption={main.py --- pipeline steps overview (abbreviated)}]
# STEP 1 --- Load field survey data
df = load_school_data(str(CSV_FILE))
# df has one row per observation point (multiple per school)

# STEP 2 --- Apply HS scoring (one function per indicator)
for idx, row in df.iterrows():
    sp = _spatial(row['School'])    # OSM features for this school
    ev = _env(row['School'])        # AQI + crime for this school
    df.loc[idx, 'HS1']  = compute_hs1(row)
    df.loc[idx, 'HS2']  = compute_hs2(row)
    df.loc[idx, 'HS3']  = compute_hs3(row, sp)
    # ... HS4 through HS10
df['HS_overall'] = df[HS_CODES].mean(axis=1, skipna=True).round(1)

# STEP 3 --- Generate charts (radar, bar, breakdown)
# See Chapter 22 of the full book for chart code.

# STEP 4 --- Generate rule-based recommendations
rec_df = run_recommendations(df)
rec_df.to_csv('outputs/recommendations.csv', index=False)

# STEP 5 --- Build interactive Leaflet.js map
build_interactive_map(df, rec_df, OUT_DIR_STR)

# STEP 6 --- Build KDE heatmap
build_kde_heatmap(df, rec_df, OUT_DIR_STR)
\end{lstlisting}

\begin{keybox}{Per-row vs per-school scoring}
\texttt{main.py} computes HS scores at the \textit{observation level} (one row per survey point). The school-level summary is then computed by grouping by school and taking the mean across all observation points:
\begin{lstlisting}[language=Python, frame=none, numbers=none, backgroundcolor=\color{navyblue!5}]
school_hs = df.groupby(['School', 'School_short'])[HS_CODES + ['HS_overall']] \
              .mean().round(2).reset_index()
school_hs.to_csv(str(HS_CSV), index=False)
\end{lstlisting}
This means that if a school has multiple observation points (e.g.\ multiple street segments around the gate), their HS scores are averaged. Preston HS had 4 observation points; the published HS scores are the mean of all four.
\end{keybox}

\section{The src/ Package Structure}

The \texttt{src/} directory contains the modular components of the pipeline. Each subdirectory is a Python package (contains an \texttt{\_\_init\_\_.py}):

\subsection{src/data/loader.py}

Loads \texttt{school\_data.csv}, strips whitespace from column headers, applies the \texttt{COLUMN\_RENAME} mapping from \texttt{config.py}, and adds a \texttt{School\_short} column using \texttt{SCHOOL\_SHORT\_NAMES}.

\begin{lstlisting}[language=Python, caption={Key logic in src/data/loader.py}]
def load_school_data(path: str) -> pd.DataFrame:
    df = pd.read_csv(path)
    df.columns = df.columns.str.strip()        # remove accidental whitespace
    df = df.rename(columns=COLUMN_RENAME)      # full name -> short name
    df['School_short'] = df['School'].map(SCHOOL_SHORT_NAMES).fillna(df['School'])
    return df
\end{lstlisting}

\subsection{src/scoring/}

Contains one Python file per HS indicator: \texttt{hs1.py} through \texttt{hs10.py}, plus \texttt{severity.py}. Each file exposes a single function:
\begin{itemize}
  \item \texttt{compute\_hs1(row: pd.Series) -> float}
  \item \texttt{compute\_hs2(row: pd.Series) -> float}
  \item \texttt{compute\_hs3(row: pd.Series, sp: dict) -> float} (sp = spatial features)
  \item \ldots and so on.
\end{itemize}

The \texttt{\_\_init\_\_.py} re-exports all ten functions so that \texttt{main.py} only needs:
\begin{lstlisting}[language=Python, numbers=none]
from src.scoring import (
    compute_hs1, compute_hs2, ..., compute_hs10, compute_severity
)
\end{lstlisting}

\subsection{src/recommendations/engine.py}

Implements rule-based recommendations triggered by indicator scores falling below thresholds. For each school-street combination, it generates up to five recommendations (one per failing indicator), each with a type, priority, estimated cost level, and suggested timeframe.

\subsection{src/visualisation/}

Two modules:
\begin{itemize}
  \item \texttt{heatmap.py} --- generates a KDE (Kernel Density Estimation) heatmap of hazard locations as both a static PNG (for QGIS import) and an interactive Folium HTML map.
  \item \texttt{interactive\_map.py} --- generates the Leaflet.js interactive hazard map with colour-coded markers, popups, and a severity filter control.
\end{itemize}

\subsection{src/scenarios/}

Two modules:
\begin{itemize}
  \item \texttt{interventions.py} --- defines the catalogue of available interventions (e.g.\ pedestrian crossing, speed reduction, bike lane) as Python dictionaries with feature delta values, cost, and timeframe.
  \item \texttt{engine.py} --- implements the delta-method scenario runner. Given a school and a list of intervention keys, it loads the trained Ridge model, applies the feature changes, and computes predicted HS score changes.
\end{itemize}

\section{Naming Conventions}

Consistent naming is enforced across all modules:

\begin{table}[H]
\centering
\caption{Naming conventions used throughout the codebase}
\begin{tabular}{L{3.5cm}L{4cm}L{5cm}}
\toprule
\textbf{Convention} & \textbf{Example} & \textbf{Applies to} \\
\midrule
\texttt{school\_short} & \texttt{Reservoir HS} & CSV column, dict key, chart label \\
\texttt{School} & \texttt{Reservoir High School} & Raw CSV school name column \\
\texttt{HS\_CODES} & \texttt{['HS1', ..., 'HS10']} & DataFrame column names \\
\texttt{HS\_overall} & score 0--10 & Overall mean of HS1--HS10 \\
\texttt{Sev\_clean} & \texttt{Major / Moderate / Minor} & Severity classification label \\
\texttt{outputs/} & all generated files & Every script writes here \\
\bottomrule
\end{tabular}
\end{table}

\begin{warnbox}{CSV column naming is critical}
The \texttt{COLUMN\_RENAME} mapping in \texttt{config.py} maps verbose Google Forms column headers to short internal names. If you re-export the field survey form and the column headers change (e.g.\ if a question is edited), you must update \texttt{COLUMN\_RENAME} accordingly. Otherwise, \texttt{loader.py} will silently leave columns unmapped, and the scoring functions will receive \texttt{NaN} inputs and produce incorrect scores with no error.
\end{warnbox}

\section{Git Workflow and Branch Strategy}

The project uses a feature-branch workflow:

\begin{itemize}
  \item \texttt{main} --- production branch. All final outputs were generated from this branch.
  \item \texttt{rish\_development} --- development branch for Hishikesh Phukan's contributions (ML pipeline, scenario engine, equity analysis).
  \item Feature branches were merged via pull requests on GitHub after code review.
\end{itemize}

The GitHub Pages dashboard is deployed from the \texttt{docs/} directory on the \texttt{main} branch. Every push to \texttt{main} that modifies \texttt{docs/} triggers an automatic rebuild of the dashboard at the project's GitHub Pages URL.

\section{Summary: What to Expect After a Full Pipeline Run}

After running all scripts in the order shown in Section~3.8, the \texttt{outputs/} directory should contain the following files:

\begin{table}[H]
\centering
\caption{Expected outputs after a full pipeline run}
\begin{tabular}{L{5.5cm}L{2cm}L{5cm}}
\toprule
\textbf{File} & \textbf{Size} & \textbf{Generated by} \\
\midrule
\texttt{spatial\_features.csv} & $\sim$20 KB & \texttt{spatial\_features.py} \\
\texttt{environmental\_features.csv} & $\sim$2 KB & \texttt{environmental\_features.py} \\
\texttt{crash\_data\_statewide.csv} & $\sim$50 MB & \texttt{crash\_analysis.py} \\
\texttt{seifa\_darebin.csv} & $\sim$2 KB & \texttt{equity\_analysis.py} \\
\texttt{hs\_scores.csv} & $\sim$3 KB & \texttt{main.py} \\
\texttt{ml\_school\_features.csv} & $\sim$5 KB & \texttt{feature\_engineering.py} \\
\texttt{ml\_predictions.csv} & $\sim$3 KB & \texttt{ml\_model.py} \\
\texttt{hs\_predictor.pkl} & $\sim$10 KB & \texttt{ml\_model.py} \\
\texttt{networks.gpkg} & $\sim$5 MB & \texttt{spatial\_features.py} \\
\texttt{chart1\_safety\_scores.png} & $\sim$300 KB & \texttt{main.py} \\
\texttt{chart2\_hazard\_severity.png} & $\sim$400 KB & \texttt{main.py} \\
\texttt{chart3\_score\_breakdown.png} & $\sim$500 KB & \texttt{main.py} \\
\texttt{chart\_hs\_correlation.png} & $\sim$400 KB & \texttt{ml\_model.py} \\
\texttt{chart\_hs\_prediction.png} & $\sim$350 KB & \texttt{ml\_model.py} \\
\texttt{chart\_feature\_importance.png} & $\sim$400 KB & \texttt{ml\_model.py} \\
\texttt{chart\_equity\_seifa.png} & $\sim$500 KB & \texttt{equity\_analysis.py} \\
\texttt{chart\_crash\_trends.png} & $\sim$500 KB & \texttt{crash\_trend\_analysis.py} \\
\texttt{map\_interactive.html} & $\sim$1 MB & \texttt{main.py} \\
\texttt{map\_heatmap.html} & $\sim$2 MB & \texttt{main.py} \\
\texttt{recommendations.csv} & $\sim$10 KB & \texttt{main.py} \\
\bottomrule
\end{tabular}
\end{table}

If any file is missing, the chapter covering that script will explain what went wrong and how to debug it.

% ============================================================
%  PART 2: DATA COLLECTION  (Chapters 5--8)
% ============================================================

\cleardoublepage
\thispagestyle{empty}
\vspace*{\fill}
\begin{center}
{\color{navyblue}\rule{0.75\linewidth}{2pt}}\\[1.2cm]
{\Huge\bfseries\color{navyblue}Part Two}\\[0.6cm]
{\LARGE\bfseries\color{navyblue}Data Collection}\\[0.6cm]
{\large\itshape Chapters 5--8: Study Area, Field Surveys, OSM, and Government Datasets}\\[1.2cm]
{\color{regengreen}\rule{0.75\linewidth}{1pt}}
\end{center}
\vspace*{\fill}
\clearpage

% ============================================================
\chapter{Study Area and School Selection}
% ============================================================

\section{The City of Darebin}

The City of Darebin is an inner-northern Local Government Area (LGA) of Melbourne, Victoria. It is bounded roughly by Bell Street to the north, the Eastern Freeway to the south, Plenty Road to the east, and the Merri Creek corridor to the west. With a 2021 estimated resident population of approximately 162,000 and an area of 53~km$^2$, Darebin is one of Melbourne's most densely settled LGAs.

The LGA encompasses the suburbs of Reservoir, Preston, Thornbury, Northcote, Fairfield, Alphington, Ivanhoe East, and Macleod. Darebin is well-served by public transport: two heavy-rail lines (Hurstbridge and Mernda), four tram routes, and numerous SmartBus services run through the LGA.

\begin{keybox}{Why Darebin?}
Three factors drove the choice of Darebin as the pilot LGA:
\begin{enumerate}
  \item \textbf{Socioeconomic diversity.} The LGA spans SEIFA IRSD deciles 4 through 7, enabling a meaningful equity test within a single jurisdiction.
  \item \textbf{Strong active-travel potential.} Darebin has a Walk Score above 80 in much of its core and an existing local cycling network, meaning street conditions are genuinely relevant to daily commuting.
  \item \textbf{Partner alignment.} Regen Melbourne's existing relationships with the City of Darebin council made field access and data sharing feasible within the project timeline.
\end{enumerate}
\end{keybox}

\section{The Three Schools}

Three state secondary schools were selected for the pilot assessment. Table~\ref{tab:five:schools} provides the key attributes used throughout the pipeline.

\begin{table}[H]
\centering
\caption{The three schools assessed: coordinates, address, and SEIFA decile.}
\label{tab:five:schools}
\begin{tabular}{L{3cm} L{4.8cm} C{1.6cm} C{1.6cm} C{1.8cm}}
\toprule
\textbf{Short name} & \textbf{Full name / Address} & \textbf{Gate lat} & \textbf{Gate lon} & \textbf{SEIFA D} \\
\midrule
Reservoir HS
  & Reservoir High School, 855 Plenty Rd, Reservoir VIC 3073
  & $-37.7224$ & $145.0294$ & D4 \\[4pt]
William Ruthven SC
  & William Ruthven Secondary College, 60 Merrilands Rd, Reservoir VIC 3073
  & $-37.6965$ & $145.0030$ & D6 \\[4pt]
Preston HS
  & Preston High School, 2--16 Cooma St, Preston VIC 3072
  & $-37.7417$ & $145.0071$ & D7 \\
\bottomrule
\end{tabular}
\end{table}

These three rows correspond exactly to the entries in \texttt{SCHOOL\_GATES} in \texttt{config.py}. Every module in the pipeline imports gate coordinates from \texttt{config.py} to ensure a single source of truth.

\begin{warnbox}{Use the Main Gate, Not the Centroid}
All spatial features and crash counts are measured from the \emph{main gate} coordinate. Using the school centroid or cadastral boundary would undercount hazards that students encounter daily, because the concentration of foot traffic is at the gate, not distributed evenly around the building.
\end{warnbox}

\section{School Selection Criteria}

The following criteria were applied to narrow the list of Darebin secondary schools from 12 to 3:

\begin{enumerate}
  \item \textbf{Government school.} Private and Catholic schools were excluded to ensure a comparable socioeconomic baseline (government schools draw from local catchments; private schools draw from wider, self-selected populations).
  \item \textbf{City of Darebin LGA only.} Confining to one LGA controls for differences in local council road investment and speed-zone policies.
  \item \textbf{SEIFA spread.} At least three different IRSD decile bands must be represented to enable an equity regression with variation across schools.
  \item \textbf{Minimum enrolment of 300.} Smaller schools generate insufficient pedestrian volumes to measure meaningful crossing utilisation rates.
  \item \textbf{Gate accessibility.} The main gate must open directly onto a public footpath (no private-road or car-park-only access).
  \item \textbf{Surveyor access.} A two-hour on-foot survey circuit around the 400-m catchment must be feasible in a single session without private-property access.
\end{enumerate}

The final three schools satisfy all six criteria and span SEIFA deciles 4, 6, and 7.

\section{Socioeconomic Context: SEIFA IRSD}

The Australian Bureau of Statistics (ABS) publishes the \textbf{Socio-Economic Indexes for Areas (SEIFA)} after each five-year Census. This study uses the \textbf{Index of Relative Socio-economic Disadvantage (IRSD)} from the 2021 Census (released February 2023).

\begin{keybox}{IRSD Decile Interpretation}
Decile 1 = most disadvantaged 10\% of Australian Statistical Areas Level 1 (SA1s). Decile 10 = least disadvantaged 10\%. An IRSD \emph{score} around 1000 is the national median; scores below 1000 indicate above-average disadvantage. The decile is derived by ranking all SA1s nationally by score and dividing into ten equal-count bands.
\end{keybox}

Table~\ref{tab:five:seifa} shows the SEIFA values used in the equity analysis (stored in \texttt{outputs/seifa\_darebin.csv}).

\begin{table}[H]
\centering
\caption{SEIFA 2021 IRSD values for the three school catchments.}
\label{tab:five:seifa}
\begin{tabular}{L{3.8cm} C{2.8cm} C{2.8cm} C{3.2cm}}
\toprule
\textbf{School} & \textbf{IRSD Score} & \textbf{IRSD Decile} & \textbf{Disadvantage Level} \\
\midrule
Reservoir HS       & $\approx 970$  & 4 & High \\
William Ruthven SC & $\approx 1010$ & 6 & Moderate \\
Preston HS         & $\approx 1040$ & 7 & Low--Moderate \\
\bottomrule
\end{tabular}
\end{table}

The equity analysis (\texttt{equity\_analysis.py}) joins these SEIFA values with computed HS scores and calculates a Pearson $r \approx 0.84$ between IRSD decile and HS overall score. This is the central policy finding: more disadvantaged catchments have measurably worse pedestrian safety outcomes. Chapter~16 covers the equity analysis implementation in detail.

\section{Catchment Definition}

For this study a school catchment is the \textbf{400-metre straight-line buffer} around the main gate coordinate. This radius was adopted because:

\begin{itemize}
  \item Australian Transport Assessment and Planning (ATAP) guidelines indicate that the majority of students who walk arrive from within a 400-m radius of the gate.
  \item The 400-m zone encompasses the most common school-zone signage, pedestrian crossing infrastructure, and drop-off/pick-up activity observed at all three sites.
  \item A single surveyor can walk and record every footpath segment within a 400-m radius in approximately two hours.
\end{itemize}

Larger buffer radii (800~m) are also computed in \texttt{spatial\_features.py} to capture cycling network length, arterial road exposure, and green space percentage --- indicators that require wider context. The 200-m buffer provides a tight ``immediate gate zone'' metric for crossing density and signal proximity.

% ============================================================
\chapter{Field Survey Methodology}
% ============================================================

\section{Why a Primary Survey?}

The Healthy Streets assessment requires observations that are not available from any open dataset: footpath surface condition, presence of tactile indicators, perceived lighting quality, on-street parking conflicts, and hazard type. These are inherently \emph{on-the-ground} observations.

The project therefore combined two data-collection methods:

\begin{enumerate}
  \item \textbf{Physical site visit.} A two-hour structured walk of every footpath segment within the 400-m catchment, using a standardised Google Forms survey on a smartphone.
  \item \textbf{Google Street View audit.} Remote inspection of segments that could not be accessed in person (e.g.\ road works, private-property constraints), using the most recent Street View imagery. The \texttt{SV\_date} column records the imagery date for audit trail purposes.
\end{enumerate}

\section{The Google Forms Instrument}

The survey was built as a Google Form with branching logic. Each row in the resulting CSV corresponds to one \emph{location observation} --- a 50-to-100-metre stretch of footpath or a specific intersection. A single session produces between 15 and 40 observations per school.

The form was designed with the following principles:

\begin{itemize}
  \item \textbf{Dropdown-first.} Every categorical field uses a fixed dropdown to prevent free-text variation that complicates parsing.
  \item \textbf{Numeric ranges.} Continuous fields (width, distance, speed limit) accept numeric input with explicit units in the label.
  \item \textbf{Observer-entered scores.} Three summary scores (FAS, CSS, EEI) are entered by the surveyor at the end of each location observation. The pipeline \emph{overwrites} these with computed values but preserves them in the \texttt{\_obs} columns for audit.
  \item \textbf{Mandatory school tag.} The ``School name'' field is required and uses a dropdown to prevent inconsistent naming.
\end{itemize}

\section{Column Rename Reference}

The raw CSV exported from Google Forms uses full question text as column headers. The data loader maps these to short internal names via the \texttt{COLUMN\_RENAME} dict in \texttt{config.py}. Table~\ref{tab:six:cols} is the complete mapping.

\begin{longtable}{L{5.8cm} L{3.2cm} L{4.5cm}}
\caption{Complete \texttt{COLUMN\_RENAME} mapping: raw CSV header to internal name.}
\label{tab:six:cols}\\
\toprule
\textbf{Raw CSV header (Google Forms)} & \textbf{Internal name} & \textbf{Description} \\
\midrule
\endfirsthead
\multicolumn{3}{c}{\small\itshape (continued from previous page)}\\
\toprule
\textbf{Raw CSV header} & \textbf{Internal name} & \textbf{Description} \\
\midrule
\endhead
\midrule
\multicolumn{3}{r}{\small\itshape (continued on next page)}\\
\endfoot
\bottomrule
\endlastfoot
School name & \texttt{School} & Full school name (dropdown) \\
Street or location being assessed & \texttt{Street} & Street name / intersection label \\
Latitude (decimal degrees) & \texttt{Latitude} & GPS latitude (auto-captured or typed) \\
Longitude (decimal degrees) & \texttt{Longitude} & GPS longitude \\
Approximate distance from school gate (metres) & \texttt{Distance\_gate} & Straight-line metres from main gate \\
Footpath present? & \texttt{Footpath\_present} & Yes / No / Partial \\
Footpath width (metres) & \texttt{Footpath\_width} & Numeric, metres \\
Footpath continuity along this segment & \texttt{Continuity} & Continuous / Interrupted / Absent \\
Footpath condition rating & \texttt{FP\_condition} & Good / Fair / Poor / Very Poor \\
Kerb ramps present at intersections? & \texttt{Kerb\_ramps} & Yes / No / Partial \\
Obstructions on footpath? & \texttt{Obstructions} & Yes / No (bins, poles, parked vehicles) \\
Pedestrian crossing present at this location? & \texttt{Crossing\_present} & Yes / No \\
Distance of nearest crossing from school gate (metres) & \texttt{Crossing\_dist} & Numeric, metres \\
Crossing visibility rating & \texttt{Visibility} & Good / Fair / Poor \\
Crossing condition / maintenance rating & \texttt{Cross\_condition} & Good / Fair / Poor \\
Tactile ground surface indicators (yellow bumps)? & \texttt{Tactile} & Yes / No \\
Pedestrian signal or countdown timer? & \texttt{Signal} & Yes / No / Pedestrian-activated \\
Posted speed limit (km/h) & \texttt{Speed\_limit} & Numeric: 40, 50, 60, 80, 100 \\
School zone active at this location? & \texttt{School\_zone} & Yes / No \\
Number of traffic lanes & \texttt{Lanes} & Numeric integer \\
Estimated traffic volume during school hours & \texttt{Traffic\_volume} & Low / Medium / High / Very High \\
Heavy vehicles or trucks present? & \texttt{Heavy\_vehicles} & Yes / No \\
Traffic calming measures present? & \texttt{Traffic\_calming} & Yes / No / Partial \\
On-street parking conflicts with pedestrians? & \texttt{Parking\_conflict} & Yes / No \\
Street lighting quality & \texttt{Lighting} & Good / Fair / Poor / None \\
Cycling infrastructure present? & \texttt{Cycling\_infra} & Dropdown (5 CIS options, see Ch.~2) \\
Hazard types observed at this location (select all that apply) & \texttt{Hazard\_types} & Multi-select; comma-separated in CSV \\
Detailed hazard description & \texttt{Hazard\_desc} & Free text \\
Recommended intervention type & \texttt{Rec\_type} & Dropdown (footpath / crossing / etc.) \\
Recommendation priority level & \texttt{Priority} & High / Medium / Low \\
Estimated cost level & \texttt{Cost\_level} & Low ($<$\$10k) / Medium / High \\
Suggested implementation timeframe & \texttt{Timeframe} & Immediate / Short-term / Long-term \\
Detailed intervention description & \texttt{Rec\_desc} & Free text \\
Expected benefit of this intervention & \texttt{Benefit} & Free text \\
Overall hazard severity at this location & \texttt{Severity} & Major / Moderate / Minor \\
Footpath Accessibility Score (0 to 10) & \texttt{FAS} & Observer-entered; overwritten by pipeline \\
Crossing Safety Score (0 to 10) & \texttt{CSS} & Observer-entered; overwritten by pipeline \\
Environmental Exposure Indicator (0 to 10) & \texttt{EEI} & Observer-entered; overwritten by pipeline \\
Primary data source used for this observation & \texttt{Data\_source} & Field / Street View / Both \\
Google Street View image date (if used) & \texttt{SV\_date} & YYYY-MM (or blank if field-only) \\
\end{longtable}

\begin{tipbox}{Why Observer Scores Are Preserved}
The data loader in \texttt{src/data/loader.py} copies the observer-entered FAS, CSS, and EEI values into \texttt{FAS\_obs}, \texttt{CSS\_obs}, and \texttt{EEI\_obs} \emph{before} the pipeline overwrites the originals with algorithmically computed values. This allows the audit section of the report to compare how closely the surveyor's field judgement matches the formula-derived score --- a quality-control step that builds confidence in both the instrument and the algorithm.
\end{tipbox}

\section{Site Visit Protocol}

Each site visit follows a fixed protocol:

\begin{description}
  \item[Equipment.] Smartphone with Google Forms offline mode enabled, measuring wheel or digital range finder, high-visibility vest, and a printed catchment map showing all segments to assess.
  \item[Timing.] Visits are conducted during school arrival or departure to observe actual pedestrian behaviour. Avoid mid-day sessions; the HS indicators require school-hours context.
  \item[Route.] Walk every footpath segment within the 400-m buffer in a systematic grid pattern. Do not skip segments even if they appear safe; absences are data.
  \item[One row per segment.] Submit a new form response for each 50--100~m segment or intersection. Short segments can be combined if conditions are uniform throughout.
  \item[GPS capture.] The form auto-captures GPS coordinates on submission. If GPS accuracy is poor ($>$15~m), type coordinates from Google Maps manually.
  \item[Photo log.] Photograph every hazard and every crossing. Photos are stored locally; filenames match the form response timestamp for traceability.
  \item[Duration.] Allow two hours per school. A 400-m buffer typically contains 20--35 segments.
\end{description}

\section{The Data Loader: \texttt{src/data/loader.py}}

After the survey CSV is exported from Google Sheets, it is read by \texttt{load\_school\_data()} in \texttt{src/data/loader.py}. This is the first function called in \texttt{main.py}.

\begin{lstlisting}[style=python, caption={Data loading and cleaning --- src/data/loader.py.}]
from pathlib import Path
from typing import Union
import pandas as pd
from config import COLUMN_RENAME, SCHOOL_SHORT_NAMES


def _normalise_severity(s: str) -> str:
    """Map a free-text severity string to Major / Moderate / Minor / Unknown."""
    s = str(s).lower()
    if 'major'    in s: return 'Major'
    if 'moderate' in s: return 'Moderate'
    if 'minor'    in s: return 'Minor'
    return 'Unknown'


def load_school_data(csv_path: Union[str, Path]) -> pd.DataFrame:
    df = pd.read_csv(csv_path)
    df.columns = df.columns.str.strip()
    df = df.rename(columns=COLUMN_RENAME)

    # Add short display name (Reservoir HS, William Ruthven SC, etc.)
    df['School_short'] = (
        df['School'].map(SCHOOL_SHORT_NAMES).fillna(df['School'])
    )

    # Preserve observer-entered scores before pipeline overwrites them
    df['FAS_obs'] = pd.to_numeric(df['FAS'], errors='coerce')
    df['CSS_obs'] = pd.to_numeric(df['CSS'], errors='coerce')
    df['EEI_obs'] = pd.to_numeric(df['EEI'], errors='coerce')
    df['Sev_obs'] = df['Severity'].apply(_normalise_severity)

    return df
\end{lstlisting}

The key steps are:

\begin{enumerate}
  \item \textbf{Strip column whitespace.} Google Forms sometimes adds trailing spaces to header text. \texttt{str.strip()} prevents silent rename failures.
  \item \textbf{Rename columns.} \texttt{df.rename(columns=COLUMN\_RENAME)} replaces every long question string with its short internal name in a single operation.
  \item \textbf{Map short names.} \texttt{SCHOOL\_SHORT\_NAMES} converts ``Reservoir High School'' to ``Reservoir~HS'' etc., ensuring chart labels and CSV join keys are consistent.
  \item \textbf{Preserve observer scores.} The \texttt{\_obs} columns are a snapshot of raw surveyor input. After this function returns, \texttt{main.py} calls the scoring engine which \emph{replaces} \texttt{FAS}, \texttt{CSS}, and \texttt{EEI} with computed values.
  \item \textbf{Normalise severity.} \texttt{\_normalise\_severity()} handles capitalization and partial strings so ``MAJOR HAZARD'' and ``major'' both map to \texttt{'Major'}.
\end{enumerate}

% ============================================================
\chapter{OSM Data with osmnx}
% ============================================================

\section{OpenStreetMap as Infrastructure Data}

OpenStreetMap (OSM) is a crowdsourced, freely-licensed spatial database that covers roads, footpaths, crossings, traffic signals, street furniture, land use, and public transport at city-scale. For pedestrian safety analysis it provides three things that government datasets do not:

\begin{itemize}
  \item \textbf{Footpath network topology.} OSM tags footpaths as \texttt{highway=footway}, \texttt{path}, \texttt{pedestrian}, etc., enabling network-graph analysis of walking route connectivity.
  \item \textbf{Crossing and signal locations.} OSM nodes tagged \texttt{highway=crossing} and \texttt{highway=traffic\_signals} give a point layer of every formal crossing within any buffer.
  \item \textbf{Street furniture.} Trees (\texttt{natural=tree}), benches (\texttt{amenity=bench}), shelters (\texttt{amenity=shelter}), and bus stops (\texttt{highway=bus\_stop}) support the shade, rest, and liveability indicators.
\end{itemize}

The \textbf{osmnx} library (Boeing, 2017) wraps the OSM Overpass API and returns data as \texttt{NetworkX} graphs or \texttt{GeoPandas} GeoDataFrames, making it straightforward to compute buffer-based counts and network metrics in Python.

\begin{keybox}{Overpass API Rate Limits}
The Overpass API is a public service. osmnx respects a 1-second delay between consecutive requests by default. For three schools the full spatial feature computation takes approximately 3--5 minutes. Running the script repeatedly in a short window may trigger a temporary block. If you see \texttt{InsufficientResponseError}, wait 60 seconds and retry.
\end{keybox}

\section{The osmnx Library}

Install osmnx with:

\begin{lstlisting}[style=bash]
conda install -c conda-forge osmnx=1.9.3
\end{lstlisting}

The two functions used most in this project are:

\begin{itemize}
  \item \texttt{ox.graph\_from\_point((lat, lon), dist=r, network\_type='walk')} --- returns a \texttt{NetworkX} MultiDiGraph of the walking network within \texttt{dist} metres.
  \item \texttt{ox.features\_from\_point((lat, lon), tags=\{...\}, dist=r)} --- returns a \texttt{GeoDataFrame} of OSM features matching the tag filter within \texttt{dist} metres.
\end{itemize}

The \texttt{network\_type} parameter accepts \texttt{'walk'}, \texttt{'drive'}, \texttt{'bike'}, or \texttt{'all'}. Each type filters the OSM highway tags to the relevant road classes.

\section{CRS and Projection}

All OSM data is returned in \textbf{WGS84 (EPSG:4326)} --- geographic coordinates in decimal degrees. Distance calculations (buffer radii, edge lengths) require a \emph{projected} CRS with metric units.

This project uses \textbf{GDA2020 / MGA zone 55 (EPSG:7855)}, the standard projected CRS for Victoria. The \texttt{pyproj} \texttt{Transformer} is initialised once at module load time:

\begin{lstlisting}[style=python, caption={CRS constants and transformer in spatial\_features.py.}]
from pyproj import Transformer

CRS_METRIC = 'EPSG:7855'   # GDA2020 / MGA zone 55 -- Victoria
CRS_GEO    = 'EPSG:4326'   # WGS84 geographic

_transformer = Transformer.from_crs(
    CRS_GEO, CRS_METRIC, always_xy=True
)


def gate_to_projected_point(lat, lon):
    """Convert WGS84 lat/lon to a Shapely Point in EPSG:7855."""
    x, y = _transformer.transform(lon, lat)
    return Point(x, y)
\end{lstlisting}

\begin{warnbox}{\texttt{always\_xy=True} Is Essential}
\texttt{Transformer.from\_crs()} defaults to respecting each CRS's axis order. EPSG:4326 specifies latitude-first, but Shapely and most geometric operations expect (x, y) = (longitude, latitude). Passing \texttt{always\_xy=True} forces the transformer to accept (lon, lat) input regardless of CRS axis convention. Omitting this flag silently swaps your coordinates, producing gate points that are hundreds of kilometres from Darebin.
\end{warnbox}

\section{Buffer Strategy}

Three buffer radii are queried for each school gate:

\begin{center}
\begin{tabular}{C{2cm} L{10cm}}
\toprule
\textbf{Radius} & \textbf{Purpose} \\
\midrule
200~m & Immediate gate zone: crossing density, signal proximity \\
400~m & Primary catchment: all HS indicator spatial features \\
800~m & Wider context: cycling network, arterial exposure, green space \\
\bottomrule
\end{tabular}
\end{center}

The RADII constant is defined as:

\begin{lstlisting}[style=python]
RADII = [200, 400, 800]   # metres
\end{lstlisting}

The outer loop in \texttt{compute\_features()} iterates over all three radii, computing the same set of metrics at each scale. The HS-specific extras (trees, benches, PT stops, etc.) are queried only at the 400-m radius.

\section{Walking Network Features}

\begin{lstlisting}[style=python, caption={Walking network computation per radius.}]
FOOTPATH_TYPES = {
    'footway', 'path', 'pedestrian',
    'cycleway', 'steps', 'bridleway'
}

for r in RADII:
    try:
        G_walk = ox.graph_from_point(
            (lat, lon), dist=r,
            network_type='walk', retain_all=True
        )
        edges_walk  = ox.graph_to_gdfs(G_walk, nodes=False).to_crs(CRS_METRIC)
        walk_edges  = len(edges_walk)
        walk_length = edges_walk.geometry.length.sum()

        fp_mask   = edges_walk['highway'].apply(
            lambda v: _highway_matches(v, FOOTPATH_TYPES))
        fp_length = edges_walk.loc[fp_mask, 'geometry'].length.sum()
        fp_pct    = (fp_length / walk_length * 100) if walk_length > 0 else 0.0
    except Exception:
        pass
\end{lstlisting}

The helper \texttt{\_highway\_matches(val, type\_set)} handles the OSM quirk where a single edge may have a \emph{list} of highway tag values (e.g.\ \texttt{['footway', 'path']}):

\begin{lstlisting}[style=python, caption={OSM list-value helper.}]
def _highway_matches(val, type_set):
    if isinstance(val, list):
        return any(v in type_set for v in val)
    return val in type_set
\end{lstlisting}

Four features are recorded per radius:

\begin{center}
\begin{tabular}{L{5cm} L{8cm}}
\toprule
\textbf{Feature} & \textbf{Definition} \\
\midrule
\texttt{walk\_edges\_Rm}      & Count of walk-network edge segments \\
\texttt{walk\_length\_Rm}     & Total walk-network length (metres) \\
\texttt{footpath\_length\_Rm} & Total length of dedicated footpath edges (metres) \\
\texttt{footpath\_pct\_Rm}    & footpath\_length / walk\_length $\times$ 100 \\
\bottomrule
\end{tabular}
\end{center}

\section{Road Network Features}

\begin{lstlisting}[style=python, caption={Road network and speed parsing per radius.}]
ARTERIAL_TYPES = {
    'primary', 'secondary', 'trunk',
    'primary_link', 'secondary_link', 'trunk_link'
}

def _parse_speed(val):
    """Parse OSM maxspeed to km/h float. Returns np.nan on failure."""
    if isinstance(val, list):
        val = val[0]
    if pd.isna(val):
        return np.nan
    s = str(val).strip().lower()
    if 'mph' in s:
        try:
            return float(s.replace('mph', '').strip()) * 1.60934
        except ValueError:
            return np.nan
    try:
        return float(s.replace('km/h', '').replace('kmh', '').strip())
    except ValueError:
        return np.nan

    # In the main loop:
    G_drive    = ox.graph_from_point((lat, lon), dist=r,
                                      network_type='drive', retain_all=True)
    edges_drive   = ox.graph_to_gdfs(G_drive, nodes=False).to_crs(CRS_METRIC)
    road_count    = len(edges_drive)
    arterial_mask = edges_drive['highway'].apply(
        lambda v: _highway_matches(v, ARTERIAL_TYPES))
    arterial_count = int(arterial_mask.sum())
    arterial_pct   = (arterial_count / road_count * 100) if road_count > 0 else 0.0
    speeds    = edges_drive['maxspeed'].apply(_parse_speed).dropna()
    avg_speed = speeds.mean() if len(speeds) > 0 else np.nan
    high_speed = int((speeds >= 60).sum()) if len(speeds) > 0 else 0
\end{lstlisting}

\texttt{\_parse\_speed()} handles three OSM maxspeed formats: plain integer (\texttt{"50"}), with unit (\texttt{"50 km/h"}), and imperial (\texttt{"30 mph"}).

Five features are recorded per radius:

\begin{center}
\begin{tabular}{L{5.5cm} L{7.5cm}}
\toprule
\textbf{Feature} & \textbf{Definition} \\
\midrule
\texttt{road\_count\_Rm}       & Count of all driveable edge segments \\
\texttt{arterial\_count\_Rm}   & Count of primary/secondary/trunk edges \\
\texttt{arterial\_pct\_Rm}     & arterial\_count / road\_count $\times$ 100 \\
\texttt{avg\_speed\_Rm}        & Mean posted speed limit (km/h) \\
\texttt{high\_speed\_road\_Rm} & Count of edges with maxspeed $\geq$ 60 km/h \\
\bottomrule
\end{tabular}
\end{center}

\section{Crossings and Signals}

Pedestrian crossings and traffic signals are queried as point features via the Overpass API. After retrieval, an explicit spatial intersection with the metric buffer is applied to exclude features returned by Overpass that fall just outside the true circular buffer:

\begin{lstlisting}[style=python, caption={Crossing and signal counts per radius.}]
    center = gate_to_projected_point(lat, lon)

    # Crossings
    try:
        crossings_gdf = ox.features_from_point(
            (lat, lon), tags={'highway': 'crossing'}, dist=r)
        crossings_gdf = crossings_gdf.to_crs(CRS_METRIC)
        buf = center.buffer(r)
        crossing_count = int(
            crossings_gdf[crossings_gdf.geometry.intersects(buf)].shape[0])
    except Exception:
        crossing_count = np.nan

    # Signals
    try:
        signals_gdf = ox.features_from_point(
            (lat, lon), tags={'highway': 'traffic_signals'}, dist=r)
        signals_gdf = signals_gdf.to_crs(CRS_METRIC)
        signal_count = int(
            signals_gdf[signals_gdf.geometry.intersects(buf)].shape[0])
    except Exception:
        signal_count = np.nan

    # Derived: crossings per km of walkable path
    if not np.isnan(crossing_count) and walk_length > 0:
        feat[f'crossing_density_{r}m'] = round(
            crossing_count / (walk_length / 1000), 2)
\end{lstlisting}

The \texttt{crossing\_density} derived feature is the key input to the HS2 (Easy to Cross) indicator: more crossings per kilometre of path suggests better pedestrian crossing provision.

\section{Healthy Streets Extras at 400~m}

Eight additional features are computed at the 400-m radius only. They feed the shade (HS3), rest (HS4), active travel (HS6), and liveability (HS8) indicators:

\begin{center}
\begin{tabular}{L{4.2cm} C{1.8cm} L{7cm}}
\toprule
\textbf{Feature} & \textbf{Radius} & \textbf{OSM query} \\
\midrule
\texttt{tree\_count\_100m}     & 100~m & \texttt{\{'natural': 'tree'\}} \\
\texttt{shelter\_count\_200m}  & 200~m & \texttt{\{'amenity': ['shelter', 'bus\_station']\}} \\
\texttt{green\_pct\_400m}      & 400~m & \texttt{\{'leisure': ['park', 'garden', ...]\}} \\
\texttt{bench\_count\_200m}    & 200~m & \texttt{\{'amenity': 'bench'\}} \\
\texttt{park\_count\_400m}     & 400~m & \texttt{\{'leisure': ['park', 'playground', ...]\}} \\
\texttt{pt\_stops\_400m}       & 400~m & \texttt{\{'highway': 'bus\_stop', 'railway': 'tram\_stop'\}} \\
\texttt{amenity\_count\_400m}  & 400~m & cafes, restaurants, libraries, community centres \\
\texttt{cafe\_count\_400m}     & 400~m & \texttt{\{'amenity': ['cafe', 'restaurant', 'fast\_food']\}} \\
\bottomrule
\end{tabular}
\end{center}

\begin{tipbox}{Green Space Proxy}
OSM does not provide land-use area fractions. The \texttt{green\_pct\_400m} feature uses a proxy: \texttt{min(count * 5.0, 40.0)}. Each tagged green-space polygon counts as approximately 5~percentage points of green coverage, capped at 40\%. This is a rough estimate; if your study area has high-quality land-use raster data (e.g.\ the Victorian Government's Digital Earth Victoria product), substituting a true green-fraction measurement will improve HS3 and HS9 scoring accuracy.
\end{tipbox}

\section{Cycling Network Features}

\begin{lstlisting}[style=python, caption={Cycling network and protected infrastructure detection.}]
    try:
        G_bike = ox.graph_from_point(
            (lat, lon), dist=r, network_type='bike', retain_all=True)
        edges_bike    = ox.graph_to_gdfs(G_bike, nodes=False).to_crs(CRS_METRIC)
        cycle_length  = edges_bike.geometry.length.sum()

        def _is_protected(e):
            hw = e.get('highway', '')
            cy = e.get('cycleway', '')
            if isinstance(hw, list): hw = hw[0] if hw else ''
            if isinstance(cy, list): cy = cy[0] if cy else ''
            return hw == 'cycleway' or str(cy) in (
                'track', 'separate', 'separated')

        protected_mask   = edges_bike.apply(_is_protected, axis=1)
        protected_length = edges_bike.loc[protected_mask, 'geometry'].length.sum()
    except Exception:
        pass
\end{lstlisting}

Three cycling features are recorded per radius:

\begin{center}
\begin{tabular}{L{5.5cm} L{7.5cm}}
\toprule
\textbf{Feature} & \textbf{Definition} \\
\midrule
\texttt{cycle\_length\_Rm}           & Total cycling network length (metres) \\
\texttt{protected\_cycle\_length\_Rm}& Length of protected / separated infrastructure \\
\texttt{cycle\_pct\_Rm}              & cycle\_length / walk\_length $\times$ 100 \\
\bottomrule
\end{tabular}
\end{center}

\section{GeoPackage Export}

After all schools are processed, network geometries are saved to \texttt{outputs/networks.gpkg} for QGIS visualisation:

\begin{lstlisting}[style=python, caption={GeoPackage layer export.}]
NETWORKS_GPKG = 'outputs/networks.gpkg'

layer_labels = {
    'walk_400m':    'Walk Network 400m',
    'walk_800m':    'Walk Network 800m',
    'cycling_400m': 'Cycling Network 400m',
    'cycling_800m': 'Cycling Network 800m',
    'roads_400m':   'Arterial Roads 400m',
    'roads_800m':   'Arterial Roads 800m',
}
for key, label in layer_labels.items():
    parts = all_geoms[key]
    if not parts:
        continue
    gdf = gpd.GeoDataFrame(pd.concat(parts, ignore_index=True), crs='EPSG:4326')
    gdf.to_file(NETWORKS_GPKG, layer=key, driver='GPKG')
\end{lstlisting}

Open \texttt{networks.gpkg} in QGIS via \textit{Layer} $\to$ \textit{Add Layer} $\to$ \textit{Add Vector Layer}. Each of the six layers will appear as a separate entry in the layer panel.

\section{Full Feature Inventory}

Table~\ref{tab:seven:features} lists every column written to \texttt{outputs/spatial\_features.csv}.

\begin{longtable}{L{5.8cm} C{1.8cm} L{5.5cm}}
\caption{All spatial features output by \texttt{spatial\_features.py} (R = 200, 400, or 800~m).}
\label{tab:seven:features}\\
\toprule
\textbf{Column name} & \textbf{Radii} & \textbf{Description} \\
\midrule
\endfirsthead
\multicolumn{3}{c}{\small\itshape (continued)}\\
\toprule\textbf{Column name} & \textbf{Radii} & \textbf{Description}\\\midrule
\endhead\bottomrule\endlastfoot
\texttt{school\_name}            & ---   & School short name \\
\texttt{gate\_lat}, \texttt{gate\_lon} & --- & Gate coordinates \\
\texttt{walk\_edges\_Rm}         & 3     & Walk-network edge count \\
\texttt{walk\_length\_Rm}        & 3     & Walk-network total length (m) \\
\texttt{footpath\_length\_Rm}    & 3     & Dedicated footpath length (m) \\
\texttt{footpath\_pct\_Rm}       & 3     & Footpath as \% of walk network \\
\texttt{road\_count\_Rm}         & 3     & Driveable edge count \\
\texttt{arterial\_count\_Rm}     & 3     & Primary/secondary/trunk edge count \\
\texttt{arterial\_pct\_Rm}       & 3     & Arterial as \% of road network \\
\texttt{avg\_speed\_Rm}          & 3     & Mean posted speed limit (km/h) \\
\texttt{high\_speed\_road\_Rm}   & 3     & Edges with maxspeed $\geq$ 60 km/h \\
\texttt{crossings\_Rm}           & 3     & Pedestrian crossing count \\
\texttt{signals\_Rm}             & 3     & Traffic signal count \\
\texttt{crossing\_density\_Rm}   & 3     & Crossings per km of walk path \\
\texttt{cycle\_length\_Rm}       & 3     & Cycling network total length (m) \\
\texttt{protected\_cycle\_Rm}    & 3     & Protected / separated cycle length (m) \\
\texttt{cycle\_pct\_Rm}          & 3     & Cycling as \% of walk network \\
\texttt{tree\_count\_100m}       & 400   & Trees within 100~m (HS3) \\
\texttt{shelter\_count\_200m}    & 400   & Shelters within 200~m (HS3) \\
\texttt{green\_pct\_400m}        & 400   & Green space proxy \% (HS3/HS9) \\
\texttt{bench\_count\_200m}      & 400   & Benches within 200~m (HS4) \\
\texttt{park\_count\_400m}       & 400   & Parks/playgrounds within 400~m (HS4/HS8) \\
\texttt{pt\_stops\_400m}         & 400   & PT stops within 400~m (HS6) \\
\texttt{amenity\_count\_400m}    & 400   & Shops/cafes/community within 400~m (HS8) \\
\texttt{cafe\_count\_400m}       & 400   & Cafes/restaurants within 400~m (HS8) \\
\end{longtable}

% ============================================================
\chapter{Government Open Datasets}
% ============================================================

\section{Overview}

Four government datasets supplement the field survey and OSM data. Table~\ref{tab:eight:sources} summarises each source.

\begin{table}[H]
\centering
\caption{Government open datasets used in the pipeline.}
\label{tab:eight:sources}
\begin{tabular}{L{3cm} L{3.5cm} L{4cm} L{2.5cm}}
\toprule
\textbf{Dataset} & \textbf{Provider} & \textbf{Used for} & \textbf{Output file} \\
\midrule
Crash data (ped/cyc) & VicRoads / data.vic.gov.au & HS5/HS7 safety; trend analysis & \texttt{crash\_data\_statewide.csv} \\
SEIFA 2021 IRSD & ABS & Equity analysis & \texttt{seifa\_darebin.csv} \\
AirWatch PM2.5 & EPA Victoria & HS10 clean air & \texttt{environmental\_features.csv} \\
Crime statistics & Crime Statistics Agency Vic. & HS7 personal safety & \texttt{environmental\_features.csv} \\
\bottomrule
\end{tabular}
\end{table}

\section{VicRoads Crash Data}

\subsection{Source and Download}

VicRoads publishes the Victoria Road Crash Database on the Victorian Government Data Directory (data.vic.gov.au). The dataset contains every recorded road traffic crash in Victoria from 2013 onward, with columns for crash date, time, severity, LGA, road name, and crash type (pedestrian, cyclist, motor vehicle).

\begin{tipbox}{Download Steps}
\begin{enumerate}
  \item Navigate to \url{https://discover.data.vic.gov.au} and search for \textbf{``Victoria Road Crash Data''}.
  \item Select the most recent annual CSV release (e.g.\ CrashStats 2021--2025).
  \item Download and save to \texttt{outputs/crash\_data\_statewide.csv}.
  \item The file is approximately 60--80~MB uncompressed.
\end{enumerate}
\end{tipbox}

\subsection{Key Columns}

The pipeline uses the following columns:

\begin{center}
\begin{tabular}{L{4cm} L{9cm}}
\toprule
\textbf{Column} & \textbf{Description} \\
\midrule
\texttt{ACCIDENT\_DATE} & Date string; parsed to year via \texttt{pd.to\_datetime()} \\
\texttt{ACCIDENT\_TIME} & Time string (\texttt{HH:MM:SS}); used for school-hours classification \\
\texttt{SEVERITY}       & Integer code: 1 = Fatal, 2 = Serious injury, 3+ = Minor \\
\texttt{LGA\_NAME}      & Upper-cased LGA name; filtered to \texttt{'DAREBIN'} \\
\texttt{DAY\_OF\_WEEK}  & Integer: 1 = Sunday, 2 = Monday, \ldots, 7 = Saturday \\
\texttt{nearest\_school}& School name nearest to crash location (added via spatial join in pre-processing) \\
\bottomrule
\end{tabular}
\end{center}

\subsection{School Hours Filter}

The \texttt{\_is\_school\_hours()} function in \texttt{crash\_trend\_analysis.py} classifies each crash as occurring during school arrival or departure windows:

\begin{lstlisting}[style=python, caption={School hours classification logic.}]
def _is_school_hours(row):
    try:
        t    = pd.to_datetime(row['ACCIDENT_TIME'], format='%H:%M:%S')
        h, m = t.hour, t.minute
        wday = int(row['DAY_OF_WEEK'])   # 1=Sun, 2=Mon..6=Fri, 7=Sat
        if wday in (1, 7):               # exclude weekends
            return False
        morning   = (h == 7 and m >= 30) or h == 8
        afternoon = ((h == 14 and m >= 30) or
                     h in (15, 16) or
                     (h == 16 and m == 0))
        return morning or afternoon
    except Exception:
        return False
\end{lstlisting}

This encodes the two school movement windows:
\begin{itemize}
  \item \textbf{Morning:} 07:30--09:00 Monday--Friday
  \item \textbf{Afternoon:} 14:30--16:30 Monday--Friday
\end{itemize}

Weekends are excluded regardless of time. The school-hours fraction is a key input to the HS5 (not too noisy) and HS7 (feel safe) indicators.

\subsection{Crash Trend Analysis}

\texttt{crash\_trend\_analysis.py} produces a four-panel chart (\texttt{outputs/chart\_crash\_trends.png}) covering:

\begin{enumerate}
  \item \textbf{Darebin LGA by year} --- stacked bar chart (Fatal / Serious / Minor) showing the annual trend from 2021 to 2025.
  \item \textbf{Three schools by year} --- grouped bars showing crash counts within 400~m of each school gate per year.
  \item \textbf{School hours vs off-peak} --- side-by-side bars per school; reveals whether crashes are concentrated in student movement windows.
  \item \textbf{Time-of-day distribution} --- 24-hour bar chart for all Darebin LGA crashes, with school hours highlighted in red and general peak hours in amber.
\end{enumerate}

Run it with:
\begin{lstlisting}[style=bash]
python crash_trend_analysis.py
\end{lstlisting}

\section{ABS SEIFA 2021 IRSD}

\subsection{Source and Download}

The 2021 SEIFA release is available from the ABS website. The IRSD is published at SA1, SA2, LGA, and postcode levels. This project uses the \textbf{LGA-level} IRSD score and decile.

\begin{tipbox}{Download Steps}
\begin{enumerate}
  \item Go to abs.gov.au and search for \textbf{``SEIFA 2021''}.
  \item Download ``Table 3: SEIFA 2021 Index of Relative Socio-economic Disadvantage by Local Government Area''.
  \item Extract the relevant rows for Darebin and save as \texttt{outputs/seifa\_darebin.csv} with columns: \texttt{School}, \texttt{IRSD\_Score\_Weighted}, \texttt{IRSD\_Decile}, \texttt{Disadvantage\_Level}.
\end{enumerate}
\end{tipbox}

\subsection{seifa\_darebin.csv Structure}

The file is a manually assembled CSV joining each school to the IRSD statistics of its SA2 catchment:

\begin{lstlisting}[style=text]
School,IRSD_Score_Weighted,IRSD_Decile,Disadvantage_Level
Reservoir HS,970,4,High
William Ruthven SC,1010,6,Moderate
Preston HS,1040,7,Low-Moderate
\end{lstlisting}

\subsection{Equity Analysis}

\texttt{equity\_analysis.py} loads and joins this file with \texttt{outputs/hs\_scores.csv}:

\begin{lstlisting}[style=python, caption={SEIFA join and equity correlation.}]
def load_and_join():
    seifa = pd.read_csv('outputs/seifa_darebin.csv')
    hs    = pd.read_csv('outputs/hs_scores.csv')
    df    = seifa.merge(hs, on='School', how='inner')
    return df.sort_values('IRSD_Decile').reset_index(drop=True)

# Pearson correlation between decile and HS overall:
corr = df['IRSD_Decile'].corr(df['HS_overall'])
# r > 0.3 => higher decile (less disadvantaged) = higher HS (safer)
\end{lstlisting}

The output chart (\texttt{outputs/chart\_equity\_seifa.png}) has three panels: a side-by-side bar chart of SEIFA decile vs HS overall per school, a scatter plot of the two variables, and a full indicator heatmap (HS1--HS10 in RdYlGn colour scale).

\section{EPA AirWatch PM2.5}

\subsection{Source}

The Environment Protection Authority (EPA) Victoria publishes real-time and historical air quality data via the AirWatch portal (\url{https://www.epa.vic.gov.au/for-community/airwatch}).

PM2.5 (fine particulate matter, $\leq 2.5\,\mu$m) is the metric used for HS10 (clean air). Annual mean PM2.5 in $\mu$g/m$^3$ is downloaded for the nearest EPA monitoring station to each school:

\begin{center}
\begin{tabular}{L{4cm} L{4cm} C{3cm}}
\toprule
\textbf{School} & \textbf{Nearest EPA station} & \textbf{Approx. distance} \\
\midrule
Reservoir HS       & Alphington     & $\sim$4~km \\
William Ruthven SC & Alphington     & $\sim$6~km \\
Preston HS         & Alphington     & $\sim$3~km \\
\bottomrule
\end{tabular}
\end{center}

\subsection{Preprocessing}

Annual-mean PM2.5 values are stored in \texttt{outputs/environmental\_features.csv} with columns \texttt{school\_name} and \texttt{aqi\_pm25}. The HS10 scoring function applies the following lookup:

\begin{center}
\begin{tabular}{C{3cm} C{3cm} C{3cm}}
\toprule
\textbf{PM2.5 ($\mu$g/m$^3$)} & \textbf{Category} & \textbf{HS10 score} \\
\midrule
$< 8$          & Clean    & 8--10 \\
$8$--$12$      & Moderate & 5--7  \\
$12$--$25$     & Elevated & 3--4  \\
$> 25$         & Poor     & 1--2  \\
\bottomrule
\end{tabular}
\end{center}

The WHO 2021 guideline for annual mean PM2.5 is 5~$\mu$g/m$^3$; the Australian NEPM standard is 8~$\mu$g/m$^3$. Darebin monitoring sites typically record annual means of 6--10~$\mu$g/m$^3$, placing them in the clean-to-moderate range.

\section{Crime Statistics Victoria}

\subsection{Source}

The Crime Statistics Agency (CSA) Victoria publishes LGA-level offence counts annually at \url{https://www.crimestatistics.vic.gov.au}. The metric used in this project is \textbf{crime rate per 100,000 population}, derived from total offences and ABS residential population.

\subsection{Preprocessing}

The crime rate is assigned at the LGA level --- all three schools are in the City of Darebin, so they share the same \texttt{crime\_rate\_per\_100k} value. This value is stored in \texttt{outputs/environmental\_features.csv} alongside the PM2.5 values.

\begin{lstlisting}[style=text]
school_name,aqi_pm25,crime_rate_per_100k
Reservoir HS,8.4,4230.5
William Ruthven SC,8.4,4230.5
Preston HS,7.9,4230.5
\end{lstlisting}

\subsection{Usage in the ML Model}

Both \texttt{aqi\_pm25} and \texttt{crime\_rate\_per\_100k} are used as features in the machine-learning model (\texttt{feature\_engineering.py}), where they appear in the \texttt{ENV\_FEATURES} list:

\begin{lstlisting}[style=python]
ENV_FEATURES = ['crime_rate_per_100k', 'aqi_pm25']
\end{lstlisting}

They supplement the 20 OSM spatial features and 4 crash aggregate features to form the 26-feature matrix fed to the Ridge regression model. Chapter~14 covers the full feature engineering and model training workflow.

\section{Joining All Data Sources}

The keybox below summarises how all four government datasets join to the school catchment table.

\begin{keybox}{Data Join Summary}
All datasets converge on \texttt{school\_name} (or \texttt{School}) as the primary join key. There is no spatial join required for SEIFA, PM2.5, or crime statistics because these are already LGA-level aggregates. The crash data \emph{does} require a spatial join (nearest-school assignment) which is performed as a pre-processing step: each crash is assigned to the school whose gate is nearest within a 400-m radius. Crashes beyond 400~m from any gate are excluded.
\end{keybox}

The recommended joining order is:

\begin{enumerate}
  \item Load \texttt{school\_data.csv} via \texttt{load\_school\_data()}.
  \item Run \texttt{spatial\_features.py} $\to$ \texttt{outputs/spatial\_features.csv}.
  \item Run \texttt{feature\_engineering.py} to merge spatial + environmental + crash features $\to$ \texttt{outputs/ml\_school\_features.csv}.
  \item Run \texttt{equity\_analysis.py} which merges SEIFA + HS scores independently.
  \item Run \texttt{crash\_trend\_analysis.py} for time-series reporting (separate from the ML pipeline).
\end{enumerate}

% ============================================================
%  PART 3: SCORING AND ANALYSIS  (Chapters 9--12)
% ============================================================

\cleardoublepage
\thispagestyle{empty}
\vspace*{\fill}
\begin{center}
{\color{navyblue}\rule{0.75\linewidth}{2pt}}\\[1.2cm]
{\Huge\bfseries\color{navyblue}Part Three}\\[0.6cm]
{\LARGE\bfseries\color{navyblue}Scoring and Analysis}\\[0.6cm]
{\large\itshape Chapters 9--12: HS Scoring Engine, Recommendations, ML Model, and Scenario Analysis}\\[1.2cm]
{\color{regengreen}\rule{0.75\linewidth}{1pt}}
\end{center}
\vspace*{\fill}
\clearpage

% ============================================================
\chapter{The Healthy Streets Scoring Engine}
% ============================================================

\section{Architecture of \texttt{src/scoring/}}

The scoring engine lives entirely in \texttt{src/scoring/}. There is one file per indicator:

\begin{lstlisting}[style=text]
src/scoring/
    __init__.py          # re-exports compute_hs1 ... compute_hs10, compute_severity
    hs1.py               # HS1 -- Pedestrians from all walks of life
    hs2.py               # HS2 -- Easy to cross
    hs3.py               # HS3 -- Shade and shelter  (needs spatial dict)
    hs4.py               # HS4 -- Places to stop and rest
    hs5.py               # HS5 -- Not too noisy
    hs6.py               # HS6 -- People choose to walk / cycle / PT
    hs7.py               # HS7 -- People feel safe
    hs8.py               # HS8 -- Things to see and do
    hs9.py               # HS9 -- People feel relaxed
    hs10.py              # HS10 -- Clean air
    severity.py          # Derives Major / Moderate / Minor from HS scores
    cys.py               # Cycling yield score helper (used by HS6)
\end{lstlisting}

Each \texttt{compute\_hsN()} function accepts a single DataFrame row (and, for HS3/HS4/HS6/HS8, an additional \texttt{sp} dict of spatial features) and returns a float in [0.0, 10.0].

\begin{tipbox}{Two Scoring Patterns}
\textbf{Additive:} starts at 0.0 and adds points for positive evidence (HS1, HS2, HS3, HS4, HS6, HS7, HS8).\\
\textbf{Deductive:} starts at 10.0 and subtracts for negative conditions (HS5, HS9, HS10).
Both patterns \texttt{clamp} the result to [0.0, 10.0] via \texttt{min(max(pts, 0.0), 10.0)}.
\end{tipbox}

\section{How main.py Calls the Engine}

\texttt{main.py} calls each scoring function inside a vectorised \texttt{.apply()} over the loaded DataFrame, then groups observations by school to produce a school-level mean:

\begin{lstlisting}[style=python, caption={Scoring invocation in main.py (simplified).}]
from src.scoring.hs1 import compute_hs1
from src.scoring.hs2 import compute_hs2
# ... imports for hs3..hs10 ...
from src.scoring.severity import compute_severity

df['HS1'] = df.apply(compute_hs1, axis=1)
df['HS2'] = df.apply(compute_hs2, axis=1)
df['HS3'] = df.apply(lambda r: compute_hs3(r, sp.get(r['School_short'], {})), axis=1)
# ... similarly for HS4..HS10 ...

HS_CODES = ['HS1','HS2','HS3','HS4','HS5','HS6','HS7','HS8','HS9','HS10']
df['HS_overall'] = df[HS_CODES].mean(axis=1).round(2)
df['Sev_clean']  = df.apply(compute_severity, axis=1)

school_scores = (df.groupby(['School', 'School_short'])[HS_CODES]
                   .mean().round(2).reset_index())
school_scores['HS_overall'] = school_scores[HS_CODES].mean(axis=1).round(2)
school_scores.to_csv('outputs/hs_scores.csv', index=False)
\end{lstlisting}

The \texttt{sp} dictionary maps school short name to the row from \texttt{outputs/spatial\_features.csv}. It is built once before the apply loop and passed in as a closure.

\section{HS1 --- Pedestrians from All Walks of Life}

HS1 measures footpath \emph{accessibility and inclusivity}: whether the path is present, wide enough, continuous, well-maintained, kerb-ramped, and unobstructed.

\begin{lstlisting}[style=python, caption={Full source: src/scoring/hs1.py.}]
def compute_hs1(row) -> float:
    pts = 0.0

    # Sub-component 1: footpath presence (max 3.0)
    p = str(row.get('Footpath_present', ''))
    if   'both sides' in p:                                   pts += 3.0
    elif 'one side'   in p:                                   pts += 2.0
    elif 'partial'    in p.lower() or 'broken' in p.lower(): pts += 1.0

    # Sub-component 2: footpath width (max 2.0)
    try:
        w = float(row['Footpath_width'])
        if   w >= 2.0: pts += 2.0
        elif w >= 1.5: pts += 1.5
        elif w >= 1.2: pts += 1.0
        elif w >= 1.0: pts += 0.5
    except (ValueError, TypeError):
        pass

    # Sub-component 3: continuity (max 2.0)
    c = str(row.get('Continuity', ''))
    if   '100%'     in c: pts += 2.0
    elif '75 to 99' in c: pts += 1.5
    elif '50 to 74' in c: pts += 1.0

    # Sub-component 4: surface condition (max 2.0)
    try:
        pts += (float(row['FP_condition']) / 5.0) * 2.0
    except (ValueError, TypeError):
        pass

    # Sub-component 5: kerb ramps (max 0.5)
    k = str(row.get('Kerb_ramps', ''))
    if   'all nearby' in k.lower(): pts += 0.5
    elif 'some'       in k.lower(): pts += 0.25

    # Sub-component 6: obstructions (max 0.5)
    o = str(row.get('Obstructions', ''))
    if   'no obstruction' in o.lower(): pts += 0.5
    elif 'minor'          in o.lower(): pts += 0.25

    return min(round(pts, 1), 10.0)
\end{lstlisting}

The theoretical maximum is 10.5 points; \texttt{min(..., 10.0)} caps it. Table~\ref{tab:nine:hs1} summarises the sub-component weights.

\begin{table}[H]
\centering
\caption{HS1 sub-component weights (max total 10.0).}
\label{tab:nine:hs1}
\begin{tabular}{L{4.5cm} C{2cm} L{5.5cm}}
\toprule
\textbf{Sub-component} & \textbf{Max pts} & \textbf{Survey column} \\
\midrule
Footpath presence     & 3.0 & \texttt{Footpath\_present} \\
Footpath width        & 2.0 & \texttt{Footpath\_width} \\
Path continuity       & 2.0 & \texttt{Continuity} \\
Surface condition     & 2.0 & \texttt{FP\_condition} (rated 0--5) \\
Kerb ramps            & 0.5 & \texttt{Kerb\_ramps} \\
Obstruction-free      & 0.5 & \texttt{Obstructions} \\
\bottomrule
\end{tabular}
\end{table}

\section{HS2 --- Easy to Cross}

HS2 is the crossing quality indicator. It is also one of the two \emph{Major severity triggers}: if HS2 $< 4.0$ the location is immediately classified as Major regardless of other scores.

\begin{lstlisting}[style=python, caption={Full source: src/scoring/hs2.py.}]
def compute_hs2(row) -> float:
    pts = 0.0

    # Crossing type (max 4.0)
    ct = str(row.get('Crossing_present', ''))
    if   'signal' in ct.lower() or 'traffic light' in ct.lower(): pts += 4.0
    elif 'raised' in ct.lower():                                   pts += 3.0
    elif 'zebra'  in ct.lower() or 'marked' in ct.lower():        pts += 2.5
    elif 'refuge' in ct.lower():                                   pts += 2.0
    elif 'informal' in ct.lower() or 'unmarked' in ct.lower():    pts += 0.5

    # Proximity (max 2.0)
    try:
        d = float(row['Crossing_dist'])
        if   d <=  30: pts += 2.0
        elif d <=  75: pts += 1.5
        elif d <= 150: pts += 1.0
        elif d <= 250: pts += 0.5
    except (ValueError, TypeError):
        pass

    # Visibility (max 2.0, rated 0-5)
    try:
        pts += (float(row['Visibility']) / 5.0) * 2.0
    except (ValueError, TypeError):
        pass

    # Tactile pavers (max 1.0)
    t = str(row.get('Tactile', ''))
    if   'both sides' in t.lower(): pts += 1.0
    elif 'one side'   in t.lower(): pts += 0.5

    # Signal / countdown timer (max 1.0)
    s = str(row.get('Signal', ''))
    if 'countdown' in s.lower():
        pts += 1.0
    elif 'yes' in s.lower() and 'no pedestrian' not in s.lower():
        pts += 0.5

    return min(round(pts, 1), 10.0)
\end{lstlisting}

\section{HS3 --- Shade and Shelter (Spatial Indicator)}

HS3 is the first indicator that draws on OSM spatial data rather than field observations. It receives a \texttt{sp} dict containing the pre-computed spatial features for the school.

\begin{lstlisting}[style=python, caption={src/scoring/hs3.py --- spatial indicator pattern.}]
import numpy as np
import pandas as pd

def compute_hs3(row, sp: dict) -> float:
    pts = 0.0

    # Tree count within 100m (max 5.0)
    tree_count = sp.get('tree_count_100m', np.nan)
    if pd.notna(tree_count):
        if   tree_count >= 15: pts += 5.0
        elif tree_count >= 8:  pts += 3.5
        elif tree_count >= 3:  pts += 2.0
        elif tree_count >= 1:  pts += 1.0

    # Shelters / bus shelters within 200m (max 2.0)
    shelter_count = sp.get('shelter_count_200m', np.nan)
    if pd.notna(shelter_count) and shelter_count > 0:
        pts += 2.0

    # Green space percentage within 400m (max 3.0)
    green_pct = sp.get('green_pct_400m', np.nan)
    if pd.notna(green_pct):
        if   green_pct >= 20: pts += 3.0
        elif green_pct >= 10: pts += 2.0
        elif green_pct >= 5:  pts += 1.0

    return min(round(pts, 1), 10.0) if pts > 0 else np.nan
\end{lstlisting}

\begin{warnbox}{Graceful Missing Data}
If no spatial features are available (e.g.\ \texttt{spatial\_features.py} has not been run), \texttt{sp} will be an empty dict and all \texttt{sp.get()} calls return \texttt{np.nan}. In that case \texttt{compute\_hs3} returns \texttt{np.nan} rather than 0.0. The school-level mean in \texttt{main.py} uses \texttt{.mean()} which skips \texttt{NaN} values, so HS3 simply does not contribute to the overall score for that school. This is intentional: absence of data is different from a score of zero.
\end{warnbox}

\section{HS5 --- Not Too Noisy (Deductive Pattern)}

HS5 uses the deductive pattern: it begins at 10.0 (a perfectly quiet street) and subtracts points for traffic volume, speed, heavy vehicles, and lane count.

\begin{lstlisting}[style=python, caption={src/scoring/hs5.py --- deductive scoring pattern.}]
def compute_hs5(row) -> float:
    pts = 10.0   # start perfect; deduct for negative conditions

    # Traffic volume (max deduction 4.0)
    v = str(row.get('Traffic_volume', '')).lower()
    if   'very high' in v or 'major arterial' in v: pts -= 4.0
    elif 'high'      in v:                          pts -= 3.0
    elif 'moderate'  in v:                          pts -= 1.5

    # Speed limit (max deduction 2.0)
    try:
        speed = float(
            ''.join(c for c in str(row.get('Speed_limit', ''))
                    if c.isdigit()) or '50')
        if   speed >= 70: pts -= 2.0
        elif speed >= 60: pts -= 1.0
        elif speed >= 50: pts -= 0.5
    except (ValueError, TypeError):
        pass

    # Heavy vehicles (max deduction 2.0)
    h = str(row.get('Heavy_vehicles', '')).lower()
    if   'frequent'  in h: pts -= 2.0
    elif 'occasion'  in h: pts -= 1.0

    # Lane count (max deduction 1.0)
    l = str(row.get('Lanes', '')).lower()
    if   '3 or more' in l: pts -= 1.0
    elif '2 lanes'   in l: pts -= 0.5

    return min(max(round(pts, 1), 0.0), 10.0)
\end{lstlisting}

HS5 is also a Major severity trigger: if HS5 $< 3.0$ the location is Major regardless of other indicators.

\section{All Ten Indicators at a Glance}

Table~\ref{tab:nine:all} summarises the design of all ten scoring functions.

\begin{longtable}{C{0.7cm} L{3.0cm} C{1.5cm} L{3.5cm} L{4.0cm}}
\caption{All ten HS scoring functions: design pattern and primary inputs.}
\label{tab:nine:all}\\
\toprule
\textbf{ID} & \textbf{Indicator} & \textbf{Pattern} & \textbf{Primary inputs} & \textbf{Severity role} \\
\midrule
\endfirsthead
\multicolumn{5}{c}{\small\itshape (continued)}\\\toprule
\textbf{ID} & \textbf{Indicator} & \textbf{Pattern} & \textbf{Primary inputs} & \textbf{Severity role}\\
\midrule\endhead\bottomrule\endlastfoot
HS1 & Pedestrians from all walks of life & Additive & \texttt{Footpath\_present}, \texttt{Footpath\_width}, \texttt{Continuity}, \texttt{FP\_condition} & Major if $<4.0$ \\[3pt]
HS2 & Easy to cross          & Additive & \texttt{Crossing\_present}, \texttt{Crossing\_dist}, \texttt{Visibility}, \texttt{Tactile}, \texttt{Signal} & Major if $<4.0$ \\[3pt]
HS3 & Shade and shelter       & Additive & \texttt{tree\_count\_100m}, \texttt{shelter\_count\_200m}, \texttt{green\_pct\_400m} (spatial) & --- \\[3pt]
HS4 & Places to stop and rest & Additive & \texttt{bench\_count\_200m}, \texttt{park\_count\_400m} (spatial) & --- \\[3pt]
HS5 & Not too noisy           & Deductive & \texttt{Traffic\_volume}, \texttt{Speed\_limit}, \texttt{Heavy\_vehicles}, \texttt{Lanes} & Major if $<3.0$ \\[3pt]
HS6 & Active travel choice    & Additive & \texttt{Cycling\_infra} (CIS), \texttt{pt\_stops\_400m}, \texttt{cycle\_pct\_400m} & --- \\[3pt]
HS7 & Feel safe               & Deductive & \texttt{Lighting}, \texttt{crime\_rate\_per\_100k}, crash proximity & --- \\[3pt]
HS8 & Things to see and do    & Additive & \texttt{amenity\_count\_400m}, \texttt{park\_count\_400m}, \texttt{cafe\_count\_400m} & --- \\[3pt]
HS9 & Feel relaxed            & Deductive & \texttt{Traffic\_volume}, \texttt{Parking\_conflict}, \texttt{arterial\_pct\_400m} & --- \\[3pt]
HS10& Clean air               & Deductive & \texttt{aqi\_pm25} (EPA), \texttt{Heavy\_vehicles}, \texttt{Traffic\_volume} & --- \\
\end{longtable}

% ============================================================
\chapter{Severity Classification and Recommendations}
% ============================================================

\section{The Severity Model}

After HS scores are computed for each field observation, \texttt{compute\_severity()} in \texttt{src/scoring/severity.py} classifies the location as \textbf{Major}, \textbf{Moderate}, or \textbf{Minor}.

\begin{lstlisting}[style=python, caption={Full source: src/scoring/severity.py.}]
import pandas as pd

_ALL_CODES = ['HS1','HS2','HS3','HS4','HS5','HS6',
              'HS7','HS8','HS9','HS10']


def compute_severity(row) -> str:
    hs1     = row.get('HS1')
    hs2     = row.get('HS2')
    hs5     = row.get('HS5')
    overall = row.get('HS_overall')

    # --- Rule 1: Major (critical safety failure) ---
    if pd.notna(hs2) and hs2 < 4.0: return 'Major'
    if pd.notna(hs1) and hs1 < 4.0: return 'Major'
    if pd.notna(hs5) and hs5 < 3.0: return 'Major'

    # --- Rule 2: Moderate (broad deficiency) ---
    all_vals = [row.get(c) for c in _ALL_CODES]
    low = sum(1 for v in all_vals if pd.notna(v) and v < 6.0)
    if low >= 2:                         return 'Moderate'
    if pd.notna(overall) and overall < 5.0: return 'Moderate'

    # --- Rule 3: Minor (everything else) ---
    return 'Minor'
\end{lstlisting}

\begin{keybox}{Threshold Rationale}
The 6.0 ``good'' threshold comes directly from the Healthy Streets framework: any indicator below 6.0 is flagged as a deficiency. The Major triggers (HS2 $<$ 4.0, HS1 $<$ 4.0, HS5 $<$ 3.0) use lower thresholds because these three indicators represent acute safety risks (impassable footpath, no crossing, high-speed/volume road). Two or more deficiencies together constitute Moderate, reflecting a compounded environment even when no single indicator is in crisis.
\end{keybox}

\section{The Recommendation Rules: \texttt{src/recommendations/rules.py}}

Seventeen rule functions are defined in \texttt{rules.py}. Each accepts a scored row and returns a recommendation dict or \texttt{None}. Table~\ref{tab:ten:rules} lists all rules.

\begin{table}[H]
\centering
\caption{All 17 recommendation rules in \texttt{src/recommendations/rules.py}.}
\label{tab:ten:rules}
\begin{tabular}{L{1.5cm} C{1.2cm} L{6cm} L{4cm}}
\toprule
\textbf{Rule} & \textbf{HS} & \textbf{Trigger condition} & \textbf{Recommendation} \\
\midrule
rule\_01 & HS1 & Footpath absent or broken & Construct 1.8m continuous footpath \\
rule\_02 & HS1 & Footpath width $<$ 1.5~m & Widen to AS~1428.1 minimum \\
rule\_03 & HS2 & No formal crossing present & Install raised zebra + tactile pavers \\
rule\_04 & HS2 & Nearest crossing $>$ 150~m & Add crossing within 50~m of gate \\
rule\_05 & HS2 & No tactile indicators & Install tactile pavers both sides \\
rule\_06 & HS1 & No school zone signage & Establish 40 km/h school zone \\
rule\_07 & HS9 & No traffic calming & Install speed humps / raised crossing \\
rule\_08 & HS6 & No cycling infrastructure & Install shared path or advisory lane \\
rule\_09 & HS7 & Poor street lighting & Upgrade to LED street lighting \\
rule\_10 & HS3 & Vegetation blocking sightlines & Trim vegetation for crossing visibility \\
rule\_11 & HS1 & Kerb ramps absent & Install DDA-compliant kerb ramps \\
rule\_12 & HS5 & Speed limit $\geq$ 60 km/h near gate & Reduce speed limit within school zone \\
rule\_13 & HS5 & Frequent heavy vehicles & Implement local freight routing policy \\
rule\_14 & HS2 & Low crossing visibility & Mark high-visibility pavement markings \\
rule\_15 & HS6 & HS6 score $<$ 5.0 & Improve PT accessibility + bike parking \\
rule\_16 & HS7 & HS7 score $<$ 5.0 & CPTED audit + lighting improvement \\
rule\_17 & HS10 & HS10 score $<$ 5.0 & Reduce idling zones; plant vegetation buffer \\
\bottomrule
\end{tabular}
\end{table}

Each rule function follows the same structure:

\begin{lstlisting}[style=python, caption={Example rule (rule\_03) from rules.py.}]
def rule_03_crossing_absent(row) -> Optional[dict]:
    if row.get('Crossing_present') in [
        'No crossing at all',
        'Yes -- informal / unmarked only'
    ]:
        return {
            'hs_indicator'  : 'HS2',
            'hazard'        : 'No formal pedestrian crossing present',
            'recommendation': 'Install raised zebra crossing with tactile '
                              'pavers adjacent to school gate',
            'priority'      : 'High',
            'cost'          : 'Medium -- $20,000 to $200,000',
            'timeframe'     : 'Short-term -- within 1 year',
        }
    return None
\end{lstlisting}

\section{The Recommendation Engine: \texttt{run\_recommendations()}}

\begin{lstlisting}[style=python, caption={Full source: src/recommendations/engine.py.}]
from src.recommendations.rules import (
    rule_01_footpath_missing,
    # ... 16 more rules ...
)
import pandas as pd

_RULES = [
    rule_01_footpath_missing,
    rule_02_footpath_narrow,
    # ... all 17 rules ...
]


def generate_recommendations(row) -> list:
    return [r for fn in _RULES if (r := fn(row)) is not None]


def run_recommendations(df: pd.DataFrame) -> pd.DataFrame:
    records = []
    for _, row in df.iterrows():
        for rec in generate_recommendations(row):
            records.append({
                'School'        : row['School'],
                'Location'      : row.get('Street', ''),
                'Severity'      : row.get('Sev_clean', ''),
                'HS_overall'    : row.get('HS_overall', ''),
                'HS_indicator'  : rec.get('hs_indicator', ''),
                'Hazard'        : rec['hazard'],
                'Recommendation': rec['recommendation'],
                'Priority'      : rec['priority'],
                'Cost'          : rec['cost'],
                'Timeframe'     : rec['timeframe'],
            })
    return pd.DataFrame(records)
\end{lstlisting}

The engine iterates over every observation row in the DataFrame, applies all 17 rules, and collects every triggered rule as a separate record. The resulting DataFrame is saved to \texttt{outputs/recommendations.csv}.

\section{Recommendations CSV Structure}

The output file has one row per triggered rule per observation location. A single 400-m catchment survey of 25 locations might generate 30--60 recommendation rows. Table~\ref{tab:ten:reccsv} shows the column structure.

\begin{table}[H]
\centering
\caption{\texttt{outputs/recommendations.csv} column structure.}
\label{tab:ten:reccsv}
\begin{tabular}{L{3.2cm} L{9.5cm}}
\toprule
\textbf{Column} & \textbf{Content} \\
\midrule
\texttt{School}         & Full school name \\
\texttt{Location}       & Street / intersection label from field survey \\
\texttt{Severity}       & Major / Moderate / Minor (from \texttt{compute\_severity}) \\
\texttt{HS\_overall}    & Overall HS score for this location (0--10) \\
\texttt{HS\_indicator}  & Which indicator triggered this recommendation (e.g.\ HS2) \\
\texttt{Hazard}         & Short description of the observed hazard \\
\texttt{Recommendation} & Specific intervention text \\
\texttt{Priority}       & High / Medium / Low \\
\texttt{Cost}           & Low ($<$\$20k) / Medium (\$20k--\$200k) / High \\
\texttt{Timeframe}      & Immediate / Short-term / Long-term \\
\bottomrule
\end{tabular}
\end{table}

\begin{tipbox}{Using recommendations.csv in a Report}
Import the CSV into Excel or Power~BI and filter by \texttt{Priority = High} and \texttt{School = Reservoir HS} to generate a council-ready action list for the highest-risk school. Sort by \texttt{Severity} (Major first) and \texttt{Timeframe} (Immediate first) to produce a ranked implementation schedule.
\end{tipbox}

% ============================================================
\chapter{Machine Learning: Feature Engineering and Ridge Regression}
% ============================================================

\section{Purpose and Context}

The machine-learning component of this project has a specific, limited aim: to demonstrate that Healthy Streets indicator scores can be \emph{approximated} from freely available open data (OSM, EPA, crime statistics, crash data) \emph{without} conducting a field survey. This serves two purposes:

\begin{enumerate}
  \item \textbf{Scalability.} If the model is accurate enough, HS scores could be estimated for the 1,000+ remaining schools in the 300,000 Streets target before field teams are deployed, allowing prioritised triage.
  \item \textbf{Validation.} If open-data features correlate well with surveyed HS scores, it confirms that the feature selection in OSM is measuring the right things.
\end{enumerate}

\begin{warnbox}{n=3 Caution}
The model is trained on three schools. This is an illustrative proof-of-concept, not a production predictor. The LOO-CV evaluation is statistically valid for n=3 but has very high variance. Do not use the predictions to replace field surveys; use them to rank schools for survey prioritisation.
\end{warnbox}

\section{Feature Engineering: \texttt{feature\_engineering.py}}

\subsection{The 26-Feature Matrix}

The feature matrix is assembled from four sources and saved to \texttt{outputs/ml\_school\_features.csv} with one row per school. Table~\ref{tab:eleven:features} shows all 26 features.

\begin{table}[H]
\centering
\caption{The 26-feature ML input matrix and their HS indicator targets.}
\label{tab:eleven:features}
\begin{tabular}{L{5.8cm} C{1.5cm} L{4.5cm}}
\toprule
\textbf{Feature} & \textbf{Source} & \textbf{HS indicator target} \\
\midrule
\texttt{footpath\_pct\_200m}             & OSM & HS1 \\
\texttt{footpath\_pct\_400m}             & OSM & HS1 \\
\texttt{crossings\_400m}                 & OSM & HS2 \\
\texttt{signals\_400m}                   & OSM & HS2 \\
\texttt{crossing\_density\_400m}         & OSM & HS2 \\
\texttt{tree\_count\_100m}               & OSM & HS3 \\
\texttt{shelter\_count\_200m}            & OSM & HS3 \\
\texttt{green\_pct\_400m}                & OSM & HS3 \\
\texttt{bench\_count\_200m}              & OSM & HS4 \\
\texttt{park\_count\_400m}               & OSM & HS4 / HS8 \\
\texttt{avg\_speed\_400m}                & OSM & HS5 / HS9 \\
\texttt{arterial\_pct\_400m}             & OSM & HS5 / HS9 \\
\texttt{high\_speed\_road\_400m}         & OSM & HS5 \\
\texttt{road\_count\_400m}               & OSM & HS5 \\
\texttt{cycle\_pct\_400m}                & OSM & HS6 \\
\texttt{protected\_cycle\_length\_400m}  & OSM & HS6 \\
\texttt{pt\_stops\_400m}                 & OSM & HS6 \\
\texttt{amenity\_count\_400m}            & OSM & HS8 \\
\texttt{cafe\_count\_400m}               & OSM & HS8 \\
\texttt{walk\_length\_400m}              & OSM & context \\
\texttt{crime\_rate\_per\_100k}          & CSA Vic & HS7 \\
\texttt{aqi\_pm25}                       & EPA Vic & HS10 \\
\texttt{crash\_count}                    & VicRoads & HS7 / HS9 \\
\texttt{serious\_or\_fatal\_rate}        & VicRoads & HS7 \\
\texttt{school\_hours\_pct}              & VicRoads & HS5 / HS9 \\
\texttt{avg\_speed\_zone}                & VicRoads & HS5 \\
\bottomrule
\end{tabular}
\end{table}

\subsection{Data Joining Code}

\begin{lstlisting}[style=python, caption={Feature matrix assembly in feature\_engineering.py.}]
# Spatial features (from spatial_features.py output)
sf = pd.read_csv('outputs/spatial_features.csv')
sf = sf[sf['school_name'].isin(SHORT_NAMES)].rename(
    columns={'school_name': 'school'})
sf_sel = sf[['school'] + SPATIAL_FEATURES]

# Environmental features
ef = pd.read_csv('outputs/environmental_features.csv')
ef['school'] = ef['school_name'].map(FULL_TO_SHORT).fillna(ef['school_name'])
ef_sel = ef[['school'] + ENV_FEATURES]

# Crash aggregates
cr = pd.read_csv('outputs/crash_data_statewide.csv', low_memory=False)
cr.columns = cr.columns.str.strip().str.upper()
cr['serious_or_fatal'] = cr['SEVERITY'].isin([1, 2]).astype(int)
cr['is_school_hours']  = (
    (pd.to_datetime(cr['ACCIDENT_DATE'], dayfirst=True,
                    errors='coerce').dt.dayofweek < 5) &
    (cr['HOUR'].between(7, 9) | cr['HOUR'].between(14, 17))
).astype(int)

crash_agg = (cr.groupby('school')
               .agg(crash_count           = ('ACCIDENT_NO', 'count'),
                    serious_or_fatal_rate  = ('serious_or_fatal', 'mean'),
                    school_hours_pct       = ('is_school_hours', 'mean'),
                    avg_speed_zone         = ('SPEED_ZONE',
                        lambda x: pd.to_numeric(x, errors='coerce')
                                    .where(lambda v: v < 200).mean()))
               .round(4).reset_index())

# HS targets
hs = pd.read_csv('outputs/hs_scores.csv')
hs['school'] = hs['School_short']

ml = (sf_sel
      .merge(ef_sel,    on='school', how='left')
      .merge(crash_agg, on='school', how='left')
      .merge(hs[['school'] + HS_TARGETS + ['HS_overall']],
             on='school', how='left'))

ml.to_csv('outputs/ml_school_features.csv', index=False)
\end{lstlisting}

\section{The Ridge Regression Model: \texttt{ml\_model.py}}

\subsection{Why Ridge Regression?}

Ridge regression (L2 regularisation) was chosen for three reasons specific to this dataset:

\begin{enumerate}
  \item \textbf{n=3 schools.} With only three training examples, any model must be heavily regularised to avoid overfitting. Ridge shrinks coefficients toward zero, providing a principled bias-variance trade-off.
  \item \textbf{Correlated features.} OSM spatial features are highly correlated with each other (e.g.\ \texttt{footpath\_pct\_200m} and \texttt{footpath\_pct\_400m}). Ridge handles multicollinearity better than OLS.
  \item \textbf{Multi-output target.} scikit-learn's \texttt{MultiOutputRegressor} wraps any single-output estimator to predict all 10 HS indicators simultaneously with independent per-indicator models.
\end{enumerate}

\subsection{Pipeline Structure}

\begin{lstlisting}[style=python, caption={The full scikit-learn pipeline in ml\_model.py.}]
from sklearn.linear_model    import Ridge
from sklearn.preprocessing   import StandardScaler
from sklearn.pipeline        import Pipeline
from sklearn.multioutput     import MultiOutputRegressor
from sklearn.model_selection import LeaveOneOut
from sklearn.metrics         import mean_absolute_error

pipe = Pipeline([
    ('scaler', StandardScaler()),          # normalise each feature to mean=0, std=1
    ('model',  MultiOutputRegressor(
                   Ridge(alpha=1.0))),     # independent Ridge per HS indicator
])
\end{lstlisting}

\texttt{StandardScaler} is essential: the 26 features span very different scales (e.g.\ \texttt{walk\_length\_400m} is in thousands of metres while \texttt{crossings\_400m} is a small integer). Without scaling, Ridge would over-penalise large-scale features.

The \texttt{alpha=1.0} Ridge parameter was chosen by convention for small datasets; in a production system with more schools it should be tuned via cross-validation grid search.

\subsection{Leave-One-Out Cross-Validation}

LOO-CV is the only unbiased hold-out strategy for n=3:

\begin{lstlisting}[style=python, caption={LOO-CV loop in ml\_model.py.}]
loo   = LeaveOneOut()
preds = np.full_like(Y.values, np.nan, dtype=float)

for train_idx, test_idx in loo.split(X):
    X_tr, X_te = X.iloc[train_idx], X.iloc[test_idx]
    y_tr        = Y.iloc[train_idx]

    fold_pipe = Pipeline([
        ('scaler', StandardScaler()),
        ('model',  MultiOutputRegressor(Ridge(alpha=1.0))),
    ])
    fold_pipe.fit(X_tr, y_tr)
    preds[test_idx] = np.clip(fold_pipe.predict(X_te), 0, 10)

mae = {c: mean_absolute_error(Y[c], preds[:, i])
       for i, c in enumerate(HS_TARGETS)}
\end{lstlisting}

In each of the three folds, the model trains on 2 schools and predicts the left-out school. The scores are clipped to [0, 10] after prediction. The per-indicator MAE is printed to console and used to identify which indicators are easiest and hardest to predict from open data.

\subsection{Full-Data Model and Serialisation}

After LOO-CV evaluation, a final model is trained on all three schools and saved for use by the scenario engine:

\begin{lstlisting}[style=python, caption={Full-data training and pickle serialisation.}]
full_pipe = Pipeline([
    ('scaler', StandardScaler()),
    ('model',  MultiOutputRegressor(Ridge(alpha=1.0))),
])
full_pipe.fit(X, Y)

with open('outputs/hs_predictor.pkl', 'wb') as f:
    pickle.dump({
        'model'   : full_pipe,
        'features': X.columns.tolist(),
        'targets' : HS_TARGETS,
        'schools' : schools,
    }, f)
\end{lstlisting}

The pickle bundle includes the feature column list alongside the model. The scenario engine uses this list to ensure it passes features in the correct column order when building the scenario feature vector.

\subsection{Three Output Charts}

\begin{description}
  \item[\texttt{chart\_hs\_correlation.png}] Pearson correlation heatmap (26 features $\times$ 10 indicators). Green cells indicate the feature is positively correlated with the HS indicator; red cells indicate negative correlation. This confirms whether the feature-indicator mapping in Table~\ref{tab:eleven:features} holds in practice.
  \item[\texttt{chart\_hs\_prediction.png}] Three side-by-side bar charts (one per school) showing actual vs LOO-CV predicted HS scores for HS1--HS10. The mean absolute error across all indicators is printed as a subtitle per school.
  \item[\texttt{chart\_feature\_importance.png}] Ten sub-plots (one per HS indicator) showing the Ridge regression coefficients ranked by magnitude. Green bars = features that increase the HS score; red bars = features that decrease it.
\end{description}

Run the model with:
\begin{lstlisting}[style=bash]
python ml_model.py
\end{lstlisting}

% ============================================================
\chapter{Scenario Analysis}
% ============================================================

\section{What Scenario Analysis Does}

The scenario engine answers the question: \emph{``If we install a signalised crossing near Reservoir HS, what would the HS scores look like?''} It does this by:

\begin{enumerate}
  \item Taking the \emph{actual} surveyed HS scores as the baseline (ground truth).
  \item Applying a pre-calibrated \emph{feature delta} (change to OSM/crash features) that represents the infrastructure change.
  \item Running the trained Ridge model on both the baseline and scenario feature vectors.
  \item Adding the model's \emph{predicted delta} to the actual baseline scores.
\end{enumerate}

This \textbf{delta method} is key: it uses the model only to predict the \emph{direction and magnitude of change}, not the absolute score. The formula is:

\[
\text{scenario\_score}[i] = \text{actual\_score}[i]
  + \bigl(\hat{y}_{\text{scenario}}[i] - \hat{y}_{\text{baseline}}[i]\bigr)
\]

The result is clamped to [0, 10].

\begin{keybox}{Why the Delta Method?}
If the model predicts absolute scores with a MAE of, say, 1.2 points, using it to predict absolute scenario scores would carry that error forward. But the model's \emph{relative} predictions (how much does a crossing change the score?) are more reliable than its absolute ones, because the relative change is insensitive to any systematic bias the model has for a particular school. The delta method separates the ground-truth baseline from the model-driven change estimate.
\end{keybox}

\section{Intervention Templates: \texttt{src/scenarios/interventions.py}}

Ten infrastructure interventions are pre-defined in \texttt{interventions.py}. Each entry specifies which ML features change and by how much, along with cost and timeframe metadata. Table~\ref{tab:twelve:ivs} lists all ten.

\begin{longtable}{L{3.5cm} L{3.5cm} C{1.2cm} L{2.8cm} C{1.8cm}}
\caption{All ten scenario interventions in \texttt{interventions.py}.}
\label{tab:twelve:ivs}\\
\toprule
\textbf{Key} & \textbf{Label} & \textbf{HS} & \textbf{Cost} & \textbf{Time} \\
\midrule
\endfirsthead
\multicolumn{5}{c}{\small\itshape (continued)}\\\toprule
\textbf{Key} & \textbf{Label} & \textbf{HS} & \textbf{Cost} & \textbf{Time}\\
\midrule\endhead\bottomrule\endlastfoot
\texttt{pedestrian\_crossing} & Install signalised crossing & HS2 & Medium \$80k--\$200k & $<$1 yr \\
\texttt{bike\_lane}           & Install protected bike lane & HS6 & High \$500k+ /km & 1--3 yr \\
\texttt{speed\_reduction}     & Reduce speed to 40 km/h zone & HS5 & Low \$5k--\$20k & $<$6 mo \\
\texttt{traffic\_calming}     & Speed humps / raised crossing & HS9 & Medium \$50k--\$150k & $<$1 yr \\
\texttt{footpath}             & Build / extend footpath & HS1 & Medium \$100k--\$300k & $<$1 yr \\
\texttt{street\_trees}        & Plant street trees & HS3 & Low \$20k--\$60k & $<$1 yr \\
\texttt{benches}              & Add street benches & HS4 & Low \$5k--\$15k & $<$6 mo \\
\texttt{pt\_stop}             & Add / improve PT stop & HS6 & Medium \$50k--\$200k & 1--2 yr \\
\texttt{shelter}              & Install covered shelter & HS3 & Low \$15k--\$40k & $<$6 mo \\
\texttt{remove\_arterial}     & Reroute heavy vehicles & HS5 & High (network approval) & 2--5 yr \\
\end{longtable}

Each intervention entry in \texttt{INTERVENTIONS} follows this structure:

\begin{lstlisting}[style=python, caption={Example intervention definition (pedestrian\_crossing).}]
INTERVENTIONS = {
    'pedestrian_crossing': {
        'label'      : 'Install signalised pedestrian crossing',
        'description': 'Add a signalised crossing adjacent to the school '
                       'gate (zebra + signals + countdown timer).',
        'deltas': {
            'crossings_400m'       : +1,
            'signals_400m'         : +1,
            'crossing_density_400m': +0.15,
        },
        'cost'     : 'Medium  (~$80k--$200k)',
        'timeframe': '< 1 year',
        'hs_target': 'HS2',
    },
    # ... 9 more interventions ...
}
\end{lstlisting>

\section{The Scenario Engine: \texttt{run\_scenario()}}

\begin{lstlisting}[style=python, caption={Core delta-method logic in src/scenarios/engine.py.}]
def run_scenario(school: str, intervention_keys: list) -> dict:
    bundle = pickle.load(open('outputs/hs_predictor.pkl', 'rb'))
    pipe, feature_cols = bundle['model'], bundle['features']

    feat_row = load_features(school)    # row from ml_school_features.csv
    actual   = load_actual_scores(school)  # row from hs_scores.csv

    X_baseline = pd.DataFrame(
        [feat_row[feature_cols].values], columns=feature_cols
    ).astype(float)

    X_scenario = X_baseline.copy()
    for key in intervention_keys:
        for feat, delta in INTERVENTIONS[key]['deltas'].items():
            if feat in X_scenario.columns:
                before = float(X_scenario[feat].iloc[0])
                after  = float(np.clip(
                    before + delta, 0, before * 3 + abs(delta) * 2))
                X_scenario[feat] = after

    # Delta method
    pred_baseline = pipe.predict(X_baseline)[0]  # shape (10,)
    pred_scenario = pipe.predict(X_scenario)[0]
    model_delta   = pred_scenario - pred_baseline

    scenario_hs = {}
    for i, code in enumerate(HS_CODES):
        base_val = float(actual[code])
        new_val  = float(np.clip(base_val + model_delta[i], 0.0, 10.0))
        scenario_hs[code] = round(new_val, 2)

    scenario_overall = round(float(np.mean(list(scenario_hs.values()))), 2)

    return {
        'school'     : school,
        'interventions': [INTERVENTIONS[k]['label'] for k in intervention_keys],
        'baseline'   : {**actual},
        'scenario'   : {**scenario_hs, 'HS_overall': scenario_overall},
        'deltas'     : {c: round(scenario_hs[c] - float(actual[c]), 2)
                        for c in HS_CODES},
    }
\end{lstlisting}

\begin{warnbox}{Feature Clamping in Interventions}
The scenario engine clamps the after-intervention feature value to \texttt{before * 3 + abs(delta) * 2}. This prevents runaway feature values if an intervention delta is very large. For example, adding a crossing (\texttt{crossings\_400m += 1}) on a school that already has 8 crossings should not produce a count of 100.
\end{warnbox}

\section{The Scenario Analyser CLI}

\texttt{scenario\_analyzer.py} is the command-line interface to the scenario engine. It supports four modes:

\begin{lstlisting}[style=bash, caption={scenario\_analyzer.py CLI usage examples.}]
# Single scenario: one school, one or more interventions
python scenario_analyzer.py \
  --school "Reservoir HS" \
  --interventions pedestrian_crossing speed_reduction

# Rank all interventions for a school (best delta first)
python scenario_analyzer.py \
  --school "Preston HS" \
  --rank-all

# Run all interventions across all schools
python scenario_analyzer.py --all-schools

# List available intervention keys
python scenario_analyzer.py --list

# Save chart to a specific path
python scenario_analyzer.py \
  --school "William Ruthven SC" \
  --interventions footpath bike_lane \
  --out outputs/scenario_wrusc_footpath_bike.png
\end{lstlisting}

\subsection{Output Chart}

For a single-school scenario, the analyser produces a two-panel chart saved to \texttt{outputs/} (or the path specified by \texttt{--out}):

\begin{description}
  \item[Left panel --- Before/After.] Horizontal bar chart showing baseline HS1--HS10 scores (dark blue) and scenario scores (teal). A dashed vertical line at 6.0 marks the ``good'' threshold.
  \item[Right panel --- Delta bars.] Horizontal bars showing the change in each indicator ($\Delta\text{HS}$). Green bars indicate improvement; red bars indicate deterioration (possible if the intervention has negative side-effects on other indicators).
\end{description}

The chart title includes the school name, the interventions applied, and the change in overall severity classification (e.g.\ ``Major $\to$ Moderate'').

\subsection{Rank-All Mode}

\texttt{--rank-all} runs every intervention key individually against the specified school and ranks them by \texttt{HS\_overall} delta:

\begin{lstlisting}[style=python, caption={rank\_all() in scenario\_analyzer.py.}]
def rank_all(school: str) -> None:
    results = []
    for key in INTERVENTIONS:
        result = run_scenario(school, [key])
        results.append({
            'key'         : key,
            'label'       : INTERVENTIONS[key]['label'],
            'delta_overall': result['deltas']['HS_overall'],
            'cost'        : INTERVENTIONS[key]['cost'],
            'timeframe'   : INTERVENTIONS[key]['timeframe'],
        })
    results.sort(key=lambda r: r['delta_overall'], reverse=True)
    for r in results:
        print(f"  {r['label']:<45}  delta={r['delta_overall']:+.2f}"
              f"  {r['cost']}  {r['timeframe']}")
\end{lstlisting}

This gives practitioners a rapid cost-effectiveness ranking: which single intervention gives the largest HS improvement per dollar at a given school?

\section{End-to-End Scenario Workflow}

To run a full scenario analysis from scratch:

\begin{lstlisting}[style=bash, caption={Complete pipeline to enable scenario analysis.}]
# 1. Run the full main pipeline to generate hs_scores.csv
python main.py

# 2. Compute spatial features (OSM queries, ~3-5 minutes)
python spatial_features.py

# 3. Build the ML feature matrix
python feature_engineering.py

# 4. Train the Ridge model and save hs_predictor.pkl
python ml_model.py

# 5. Run a scenario
python scenario_analyzer.py \
  --school "Reservoir HS" \
  --interventions pedestrian_crossing footpath
\end{lstlisting}

Steps 1--4 need to be run only once (or when field survey data is updated). Step 5 can be run repeatedly with different intervention combinations without re-training.

% ============================================================
%  PART 4
% ============================================================
\part{Analysis, Equity, and Dashboard Deployment}

\begin{center}
\vspace{1cm}
{\large\itshape
Part 4 closes the pipeline loop. Chapter 13 dives into the VicRoads crash record ---
year-on-year trends, school-hours concentration, and time-of-day peaks.
Chapter 14 joins ABS SEIFA disadvantage indices to the HS scores, revealing the
equity gap between high-disadvantage and lower-disadvantage school catchments.
Chapter 15 documents the end-to-end dashboard pipeline: how \texttt{prepare\_data.py}
assembles a single \texttt{data.json} from every upstream output and serves it
to a GitHub Pages web application built on Leaflet, Chart.js, and Tailwind CSS.
}
\vspace{1cm}
\end{center}

% ============================================================
\chapter{Crash Trend Analysis}
% ============================================================

\section{Why Crash Data Matters}

The Healthy Streets framework captures \emph{perceived} risk and \emph{infrastructure quality},
but crash records capture \emph{revealed} risk --- actual harm that has already occurred.
\texttt{crash\_trend\_analysis.py} answers four questions:

\begin{enumerate}
  \item How has the volume of pedestrian and cyclist crashes across Darebin LGA changed from 2021 to 2025?
  \item Which of our three schools has the worst crash record within a 400\,m gate buffer?
  \item What fraction of crashes occur during school drop-off and pick-up hours?
  \item At what time of day are crashes most concentrated?
\end{enumerate}

Output is a single four-panel PNG saved to \texttt{outputs/chart\_crash\_trends.png}.

\begin{keybox}{Data source}
\textbf{VicRoads open data} --- statewide ped/cyc crash records, 2021--2025.
Downloaded by \texttt{crash\_analysis.py} (pipeline Step~1) and saved as
\texttt{outputs/crash\_data\_statewide.csv}.
Each row is one crash event with columns including \texttt{ACCIDENT\_DATE},
\texttt{ACCIDENT\_TIME}, \texttt{DAY\_OF\_WEEK}, \texttt{SEVERITY},
\texttt{LGA\_NAME}, \texttt{LATITUDE}, \texttt{LONGITUDE}, and
\texttt{nearest\_school} (added during the 400\,m spatial join).
\end{keybox}

\section{Loading and Filtering}

\begin{lstlisting}[style=python, caption={load\_data() in crash\_trend\_analysis.py}]
def load_data():
    path = os.path.join(OUT_DIR, 'crash_data_statewide.csv')
    df = pd.read_csv(path, low_memory=False)
    df['year']      = pd.to_datetime(df['ACCIDENT_DATE'],
                                     errors='coerce').dt.year
    df['hour']      = pd.to_datetime(df['ACCIDENT_TIME'],
                                     format='%H:%M:%S',
                                     errors='coerce').dt.hour
    df['sev_label'] = df['SEVERITY'].apply(_severity_label)
    df['school_hours'] = df.apply(_is_school_hours, axis=1)

    # Darebin LGA subset
    darebin = df[df['LGA_NAME'].str.upper() == 'DAREBIN'].copy()

    # Our 3 schools subset (normalise name variants)
    our = df[df['nearest_school'].isin(SCHOOL_NAME_MAP.keys())].copy()
    our['school_short'] = our['nearest_school'].map(SCHOOL_NAME_MAP)

    return df, darebin, our
\end{lstlisting}

Three DataFrames are returned:

\begin{description}
  \item[\texttt{df}] Full statewide dataset --- used only for LGA-level extraction.
  \item[\texttt{darebin}] All Darebin LGA crashes. Used for Panel~1 (LGA trends) and Panel~4 (time-of-day).
  \item[\texttt{our}] Crashes within 400\,m of the three school gates. Used for Panels~2 and~3.
\end{description}

\subsection{Severity Labelling}

VicRoads encodes severity as an integer: 1 = Fatal, 2 = Serious, anything else = Minor.

\begin{lstlisting}[style=python, caption={Severity mapping function.}]
def _severity_label(code):
    if code == 1: return 'Fatal'
    elif code == 2: return 'Serious'
    else: return 'Minor'
\end{lstlisting}

\subsection{School Hours Filter}

The school-hours flag is a row-level predicate that returns \texttt{True} when a crash
falls on a weekday during morning drop-off (07:30--09:00) or afternoon pick-up (14:30--16:30).
\texttt{DAY\_OF\_WEEK} is coded 1=Sunday \ldots 7=Saturday in the VicRoads dataset, so
weekdays are values 2--6.

\begin{lstlisting}[style=python, caption={School-hours predicate.}]
def _is_school_hours(row):
    try:
        t  = pd.to_datetime(row['ACCIDENT_TIME'], format='%H:%M:%S')
        h, m = t.hour, t.minute
        wday = int(row['DAY_OF_WEEK'])   # 1=Sun, 2=Mon...6=Fri, 7=Sat
        if wday in (1, 7):
            return False
        morning   = (h == 7 and m >= 30) or h == 8
        afternoon = ((h == 14 and m >= 30)
                     or h in (15, 16)
                     or (h == 16 and m == 0))
        return morning or afternoon
    except Exception:
        return False
\end{lstlisting}

\begin{warnbox}{DAY\_OF\_WEEK encoding}
VicRoads uses 1-based Sunday-first encoding, so Monday = 2 and Friday = 6.
The filter explicitly excludes \texttt{wday in (1, 7)} to drop Saturday and Sunday,
not 0 and 6.  This is a common off-by-one trap when working with this dataset.
\end{warnbox}

\section{Four-Panel Chart}

\subsection{Panel 1 --- Darebin LGA Crashes by Year}

A stacked bar chart showing Fatal/Serious/Minor crashes across all Darebin LGA schools,
grouped by year (2021--2025). Each bar is annotated with the stratum count (white text)
and the total (black text above the bar).

\subsection{Panel 2 --- Our Three Schools by Year}

Grouped bars: three schools side-by-side for each year. School colours are fixed:
Reservoir HS \texttt{\#1A476E}, William Ruthven SC \texttt{\#1A8FC1}, Preston HS \texttt{\#C0392B}.

\subsection{Panel 3 --- School Hours vs Off-Peak}

For each school: two bars, school-hours crashes vs off-peak crashes.
This panel is the most policy-relevant --- a school with a high school-hours fraction
has infrastructure risk specifically during the vulnerable window.

\subsection{Panel 4 --- Time-of-Day Distribution}

Histogram over hours 0--23 for all Darebin LGA crashes.
Bars are coloured by period:

\begin{itemize}
  \item \textcolor{red}{\textbf{Red}} --- school hours (07:00--09:00 and 14:00--16:00)
  \item \textcolor{orange}{\textbf{Orange}} --- general peak hours (06:00--09:00 and 16:00--19:00)
  \item \textcolor{codegray}{\textbf{Grey}} --- off-peak
\end{itemize}

\begin{lstlisting}[style=python, caption={Colour assignment for time-of-day histogram.}]
bar_colours = []
for h in range(24):
    if (7 <= h <= 8) or (14 <= h <= 16):
        bar_colours.append('#C0392B')   # school hours
    elif 6 <= h <= 9 or 16 <= h <= 19:
        bar_colours.append('#D35400')   # peak hours
    else:
        bar_colours.append('#BDC3C7')   # off-peak
\end{lstlisting}

\section{Key Findings Function}

\texttt{print\_findings()} summarises the chart to stdout after saving the PNG.
It reports the per-year Darebin total, per-school breakdown (total and school-hours count),
the trend direction (first year vs last year), and the peak crash hour across Darebin.

\begin{lstlisting}[style=bash, caption={Running the crash trend analysis.}]
python crash_trend_analysis.py
# Outputs:
# outputs/chart_crash_trends.png
# Console: crash count table, trend direction, peak hour
\end{lstlisting}

\begin{tipbox}{School-hours concentration --- the key metric}
Divide school-hours crashes by total crashes for each school.
A ratio above 0.30 suggests that the school gate environment specifically
(not just the general neighbourhood) is generating disproportionate risk
during the most vulnerable pedestrian window.
\end{tipbox}

\section{Integration with run\_all.py}

This script is Step~9 in the master pipeline runner:

\begin{lstlisting}[style=text, caption={Step 9 in run\_all.py STEPS list.}]
{
  'n':      9,
  'script': 'crash_trend_analysis.py',
  'label':  'Crash trend analysis (2021-2025)',
  'checks': ['outputs/chart_crash_trends.png'],
  'note':   'Year trends, school-hours breakdown, time-of-day distribution.',
}
\end{lstlisting}

If \texttt{outputs/chart\_crash\_trends.png} already exists and \texttt{--force} is not
passed, the step is skipped automatically.

% ============================================================
\chapter{Equity Analysis}
% ============================================================

\section{The Equity Hypothesis}

The 300,000 Streets project is not just a safety study --- it is an equity study.
The central hypothesis is that \emph{more socioeconomically disadvantaged school catchments
have worse pedestrian infrastructure}, compounding the disadvantage for students who are
more likely to walk or cycle to school and less likely to be driven.

\texttt{equity\_analysis.py} tests this hypothesis by joining:

\begin{itemize}
  \item \textbf{ABS SEIFA 2021 IRSD} (Index of Relative Socio-economic Disadvantage) at the
        SA1 level, aggregated to school catchment; and
  \item \textbf{Healthy Streets overall scores} from \texttt{outputs/hs\_scores.csv}.
\end{itemize}

Output: \texttt{outputs/chart\_equity\_seifa.png} (three-panel chart).

\section{SEIFA Data Pipeline}

SEIFA data for Darebin is produced by \texttt{seifa\_analysis.py} (pipeline Step~7),
which reads the ABS-released IRSD scores for each SA1 in Darebin, performs a
catchment-weighted average based on which SA1s overlap each school's 800\,m buffer,
and writes \texttt{outputs/seifa\_darebin.csv} with columns:

\begin{center}
\begin{tabular}{lL{8cm}}
\toprule
Column & Meaning \\
\midrule
\texttt{School} & Full school name (join key to \texttt{hs\_scores.csv}) \\
\texttt{School\_short} & Short name (\texttt{Reservoir HS} etc.) \\
\texttt{IRSD\_Score\_Weighted} & Population-weighted average IRSD score (national mean $\approx$ 1000) \\
\texttt{IRSD\_Decile} & Decile 1--10 (1 = most disadvantaged nationally) \\
\texttt{Disadvantage\_Level} & Human-readable label (High / Moderate / Low) \\
\bottomrule
\end{tabular}
\end{center}

\begin{keybox}{SEIFA decile interpretation}
Decile 1 means the catchment falls in the bottom 10\% nationally by socioeconomic
advantage.  Decile 10 is the most advantaged.  All three Darebin schools sit in
Deciles 4--7, which is typical for inner-northern Melbourne --- below the national
median but not the most disadvantaged nationally.
The school order is:
\textbf{Reservoir HS} (D4) $\to$ \textbf{William Ruthven SC} (D6) $\to$ \textbf{Preston HS} (D7).
\end{keybox}

\section{Data Join}

\begin{lstlisting}[style=python, caption={load\_and\_join() in equity\_analysis.py}]
def load_and_join():
    seifa = pd.read_csv(os.path.join(OUT_DIR, 'seifa_darebin.csv'))
    hs    = pd.read_csv(os.path.join(OUT_DIR, 'hs_scores.csv'))
    df    = seifa.merge(hs, on='School', how='inner')
    df    = df.sort_values('IRSD_Decile').reset_index(drop=True)
    return df
\end{lstlisting}

The inner join on \texttt{School} (full name) means both files must use the same school
name strings.  Sorting by \texttt{IRSD\_Decile} ascending ensures the charts display
most-disadvantaged first --- a deliberate visual choice to make the equity gradient
immediately readable.

\section{Three-Panel Chart}

\subsection{Panel 1 --- Side-by-Side Bars}

For each school: two bars side-by-side --- IRSD Decile (grey) and HS Overall score
(school colour).  A horizontal dashed red line marks HS = 6.0 (the Healthy Streets
``good'' threshold).  A SEIFA decile and an HS score on the same y-axis (0--10) are
comparable because both are unitless indices, but the reader must be reminded that
they measure different things.

\subsection{Panel 2 --- Scatter Plot with Quadrant Shading}

Each school is a labelled scatter point.  The x-axis is IRSD Decile, y-axis is HS Overall.
A shaded lower-left quadrant marks ``high disadvantage + poor safety'' --- the highest
policy-priority zone.

\begin{lstlisting}[style=python, caption={Quadrant shading in Panel 2.}]
ax2.axhspan(0, 6, xmin=0, xmax=0.5, alpha=0.06, color='#C0392B',
            label='High risk zone')
ax2.text(1.2, 3.0, 'Most at risk:\nhigh disadvantage\n+ poor safety',
         fontsize=7.5, color='#C0392B', style='italic', alpha=0.8)
\end{lstlisting}

The \texttt{axhspan} spans the full horizontal extent from 0 to 50\% of the x-axis
(i.e.\ deciles 1--5) and the vertical extent 0--6 (below the HS threshold).

\subsection{Panel 3 --- Indicator Heatmap}

A full-width heatmap (schools $\times$ 10 indicators) using the \texttt{RdYlGn} colourmap
(red = low, green = high, white = 5).  Cell values are printed inside each cell;
text colour switches from white to dark depending on whether the cell is very red
($< 3.5$) or very green ($> 8.5$).

SEIFA annotations appear to the left of the heatmap, one line per school, showing both
the raw weighted IRSD score and the decile:

\begin{lstlisting}[style=python, caption={SEIFA annotation on the left of the heatmap.}]
for i, (_, row) in enumerate(df.iterrows()):
    ax3.annotate(
        f'IRSD {row["IRSD_Score_Weighted"]:.0f}  D{row["IRSD_Decile"]:.0f}',
        xy=(0, i), xycoords=('axes fraction', 'data'),
        xytext=(-5, 0), textcoords='offset points',
        ha='right', va='center', fontsize=8.5, color='#444444',
    )
\end{lstlisting}

\section{Pearson Correlation and Policy Interpretation}

\texttt{print\_findings()} computes a Pearson r between IRSD Decile and HS Overall:

\begin{lstlisting}[style=python, caption={Correlation and interpretation.}]
corr = df['IRSD_Decile'].corr(df['HS_overall'])
print(f"  Correlation (IRSD Decile x HS Overall): r = {corr:.2f}")
if corr > 0.3:
    print("  FINDING: More disadvantaged catchments have WORSE safety scores.")
    print("  Higher decile (less disadvantaged) = higher HS score.")
\end{lstlisting}

With only three schools the Pearson r cannot be interpreted with standard significance
thresholds, but the observed r $\approx$ 0.84 is directionally very strong and consistent
with the a priori hypothesis.  The heatmap reinforces this: Reservoir HS (D4) scores
lowest on indicators like HS3 (shade/shelter), HS5 (traffic noise), and HS6 (active travel),
while Preston HS (D7) --- the least disadvantaged of the three --- scores highest overall.

\begin{warnbox}{Three-school limitation}
Pearson r with n = 3 has a critical value of 0.997 at $p < 0.05$ (two-tailed).
The observed r = 0.84 is not statistically significant at conventional thresholds.
It is presented as a \emph{directional finding} to motivate expansion of the study
to 20+ schools in future iterations where genuine inference is possible.
\end{warnbox}

\section{Running the Equity Analysis}

\begin{lstlisting}[style=bash, caption={Equity analysis run command.}]
# Requires seifa_darebin.csv and hs_scores.csv to exist first
python equity_analysis.py
# Output: outputs/chart_equity_seifa.png
\end{lstlisting}

In the master pipeline this is Step~8 (\texttt{run\_all.py}), run automatically after
\texttt{seifa\_analysis.py} (Step~7).

% ============================================================
\chapter{Dashboard and Deployment Pipeline}
% ============================================================

\section{Architecture Overview}

The final stage of the pipeline transforms every CSV, PNG, and model artefact into a
single-page web application hosted on GitHub Pages.  The architecture is intentionally
dependency-free on the client side: no Node.js build step, no bundler, no server.

\begin{keybox}{Dashboard technology stack}
\begin{itemize}
  \item \textbf{Leaflet 1.9.4} --- interactive map (school markers, crash heatmap, walk/cycling networks)
  \item \textbf{Chart.js} --- radar charts (HS indicator profiles), bar charts (scenario explorer)
  \item \textbf{Tailwind CSS} (CDN) --- utility-first responsive layout
  \item \textbf{Vanilla JS (ES2020)} --- no framework; all state in a single \texttt{appData} object
  \item \textbf{GitHub Pages} --- static hosting from \texttt{docs/} directory on \texttt{main} branch
\end{itemize}
\end{keybox}

The data bridge between the Python pipeline and the browser is a single file:
\texttt{docs/data/data.json}.  This file is built by \texttt{prepare\_data.py} and
committed to the repository.  The browser fetches it once on page load with a single
\texttt{fetch('./data/data.json')} call.

\section{prepare\_data.py --- The Data Builder}

\texttt{prepare\_data.py} assembles every piece of data the dashboard needs into
\texttt{data.json}.  It must be run \emph{after} the full pipeline (\texttt{run\_all.py})
has completed.

\subsection{Required Inputs}

\begin{lstlisting}[style=python, caption={Required files checked at startup.}]
REQUIRED_FILES = {
    'hs_scores':       OUTPUTS / 'hs_scores.csv',
    'recommendations': OUTPUTS / 'recommendations.csv',
    'ml_predictions':  OUTPUTS / 'ml_predictions.csv',
    'seifa':           OUTPUTS / 'seifa_darebin.csv',
    'config':          ROOT / 'config.py',
}
\end{lstlisting}

If any required file is missing, \texttt{prepare\_data.py} exits with a descriptive
error message before attempting any computation.

\subsection{build\_schools()}

\texttt{build\_schools()} is the core assembler.  For each school it:

\begin{enumerate}
  \item Looks up the school gate coordinates from \texttt{config.SCHOOL\_GATES}.
  \item Matches the HS scores row from \texttt{hs\_scores.csv}.
  \item Joins the SEIFA decile from \texttt{seifa\_darebin.csv}.
  \item Filters and formats the recommendations from \texttt{recommendations.csv}.
  \item Calls \texttt{compute\_severity()} (same logic as \texttt{src/scoring/severity.py}).
  \item Returns a dict with all the data the dashboard needs for that school.
\end{enumerate}

The returned list is the \texttt{schools} key in \texttt{data.json}, and it powers the
school tabs, radar charts, recommendation panels, and scenario explorer simultaneously.

\subsection{build\_scenarios()}

\texttt{build\_scenarios()} pre-computes all 10 interventions $\times$ 3 schools = 30
scenario results at build time.  This avoids requiring the scenario engine to run in
the browser (which would need Python).  The engine is imported directly:

\begin{lstlisting}[style=python, caption={Pre-computing scenarios during data build.}]
from src.scenarios.engine import run_scenario
from src.scenarios.interventions import INTERVENTIONS

def build_scenarios(schools):
    results = {}
    for school in schools:
        name = school['short_name']
        results[name] = {}
        for key in INTERVENTIONS:
            try:
                r = run_scenario(name, [key])
                results[name][key] = {
                    'label'  : r['interventions'][0],
                    'deltas' : r['deltas'],
                    'cost'   : INTERVENTIONS[key]['cost'],
                    'timeframe': INTERVENTIONS[key]['timeframe'],
                }
            except Exception as e:
                results[name][key] = {'error': str(e)}
    return results
\end{lstlisting}

The pre-computed deltas are serialised into \texttt{data.json} as nested dicts.
The dashboard JavaScript reads \texttt{data.scenarios[schoolName][interventionKey].deltas}
and renders the result chart instantly with no network round-trip.

\subsection{Map Layer Builders}

Three functions build the geospatial data layers embedded in \texttt{data.json}:

\begin{center}
\begin{tabular}{lL{9cm}}
\toprule
Function & Output \\
\midrule
\texttt{build\_crash\_geojson()} & GeoJSON FeatureCollection of crash points from
  \texttt{outputs/crash\_data\_darebin.csv}, each with a \texttt{severity} property \\
\texttt{build\_heatmap\_points()} & List of \texttt{[lat, lon, intensity]} triples for
  the Leaflet heatmap layer.  Intensity is $(4 - \text{SEVERITY}) / 3.0$, mapping
  Fatal $\to$ 1.0, Serious $\to$ 0.67, Minor $\to$ 0.33 \\
\texttt{build\_network\_geojson()} & Reads \texttt{outputs/networks.gpkg} layers
  \texttt{walk\_400m} and \texttt{cycling\_400m} via GeoPandas, converts to EPSG:4326,
  and rounds coordinates to 5 decimal places to reduce JSON size \\
\bottomrule
\end{tabular}
\end{center}

\begin{tipbox}{GeoJSON coordinate precision}
Five decimal places ($\approx 1$\,m precision at the equator) is sufficient for
street-level display and reduces the \texttt{data.json} file size significantly
compared to the 15-decimal-place default from GeoPandas.  The rounding is applied
with a recursive helper that handles both \texttt{Point} and \texttt{LineString}
coordinate arrays.
\end{tipbox}

\subsection{Chart Copying}

\texttt{copy\_charts()} copies all pipeline chart PNGs from \texttt{outputs/} to
\texttt{docs/data/charts/}.  The dashboard's Analysis section displays these charts
as \texttt{<img>} tags referencing relative paths.

\begin{lstlisting}[style=python, caption={Chart files copied to docs/data/charts/.}]
CHART_FILES = [
    'chart1_safety_scores.png',
    'chart2_hazard_severity.png',
    'chart3_score_breakdown.png',
    'chart_hs_correlation.png',
    'chart_hs_prediction.png',
    'chart_feature_importance.png',
    'chart_equity_seifa.png',
    'chart_crash_trends.png',
    'chart4_demographics.png',
]
\end{lstlisting}

\subsection{Final data.json Structure}

\begin{lstlisting}[style=text, caption={Top-level keys in data.json}]
{
  "generated_at":   "2025-06-01T10:23:45Z",
  "stats":          { ... hero banner stats ... },
  "schools":        [ ... 3 school objects ... ],
  "ml_results":     { ... LOO-CV predictions ... },
  "scenarios":      { school: { intervention: deltas } },
  "crash_geojson":  { type: "FeatureCollection", features: [...] },
  "heatmap_points": [ [lat, lon, intensity], ... ],
  "networks":       { "walk": {...GeoJSON...}, "cycling": {...} },
  "spatial_stats":  { school: { feature: value } },
  "charts":         { key: "filename.png", key_caption: "..." }
}
\end{lstlisting}

\begin{lstlisting}[style=bash, caption={Running prepare\_data.py.}]
# After run_all.py has completed:
python prepare_data.py

# Serves the dashboard locally:
cd docs && python -m http.server 8080
# Then open http://localhost:8080 in a browser
\end{lstlisting}

\section{Dashboard Sections}

\subsection{Hero Banner}

Five stat cards rendered by \texttt{initHero()} from \texttt{data.stats}:

\begin{center}
\begin{tabular}{lL{8cm}}
\toprule
Stat & Source \\
\midrule
Schools assessed & \texttt{len(schools)} (always 3) \\
Major hazards & Count of schools with severity = Major \\
Darebin crashes & Row count in \texttt{crash\_data\_darebin.csv} \\
Equity--safety correlation & Hard-coded \texttt{r = 0.84} from equity analysis \\
Peak crash hour & Most frequent hour in Darebin LGA crash records \\
\bottomrule
\end{tabular}
\end{center}

\subsection{Interactive Map}

The Leaflet map is initialised by \texttt{initMap()} and contains four toggleable layers:

\begin{enumerate}
  \item \textbf{School gate markers} --- custom coloured circle markers at the GPS coordinates
        from \texttt{config.SCHOOL\_GATES}, with popups showing HS Overall, severity badge,
        and per-indicator OSM stats from \texttt{spatial\_stats}.
  \item \textbf{Crash heatmap} --- Leaflet.heat plugin rendering the intensity-weighted
        \texttt{heatmap\_points} list.  Crash concentration near school gates is immediately
        visible.
  \item \textbf{Walking network} --- GeoJSON LineString layer from \texttt{networks.walk},
        styled in blue (\texttt{\#1A476E}), showing the OSM pedestrian network within
        400\,m of each gate.
  \item \textbf{Cycling network} --- GeoJSON LineString layer from \texttt{networks.cycling},
        styled in green (\texttt{\#2E7D32}), showing cycle routes.
\end{enumerate}

A layer control (top-right corner) lets users toggle each layer independently.

\subsection{School Tabs and Radar Charts}

\texttt{initSchools()} renders one tab per school.  Clicking a tab reveals:

\begin{itemize}
  \item A Chart.js radar chart of all 10 HS indicator scores.
  \item A severity badge and overall score.
  \item A SEIFA context strip (IRSD score and decile).
  \item The top three recommendations for that school.
\end{itemize}

Radar charts are cached in \texttt{radarCharts\{\}} so that switching tabs does not
re-render existing Chart.js instances (which would cause a canvas memory leak).

\subsection{Scenario Explorer}

\texttt{initScenario()} builds the interactive what-if panel.  The user selects a school
tab and an intervention button.  The dashboard looks up
\texttt{data.scenarios[school][intervention]} and renders a Chart.js bar chart comparing
baseline vs scenario HS scores for all 10 indicators.

Each intervention button shows the cost band and timeframe.  The result panel includes
a delta method disclaimer noting that results should be treated as directional with n = 3.

\subsection{Analysis Section}

A tabbed grid of all chart PNGs from \texttt{docs/data/charts/}.  The charts are grouped
into four tabs: Safety Scores, Machine Learning, Equity, and Crash Trends.  Each chart
is displayed with its caption from \texttt{data.charts}.

\subsection{Recommendations Panel}

\texttt{initRecommendations()} flattens all recommendations from all schools into a
single list and renders a filterable table.  Filters include school, severity, priority,
and HS indicator.  Rows are sorted by Priority (High first) then HS Overall ascending,
ensuring the most urgent interventions for the most at-risk schools appear first.

\section{run\_all.py --- The Master Pipeline Runner}

\texttt{run\_all.py} is the single command a developer runs to execute the full
10-step pipeline from crash data download through to dashboard data build.

\subsection{Skip-if-Exists Logic}

Each step defines a \texttt{checks} list of expected output file paths.  If all
check paths exist and \texttt{--force} is not passed, the step is skipped.
This makes \texttt{run\_all.py} safe to re-run at any time without re-downloading
data or re-training the model unnecessarily.

\begin{lstlisting}[style=python, caption={Skip-if-exists check in run\_all.py}]
def _all_exist(paths):
    return all(os.path.exists(p) for p in paths)

for step in STEPS:
    if not FORCE and _all_exist(step['checks']):
        print(f"  [{n}/10] {label}  already done")
        continue
    # ... run the step
\end{lstlisting}

\subsection{Pipeline Steps Summary}

\begin{center}
\begin{longtable}{C{0.5cm}L{4.5cm}L{5cm}L{2.5cm}}
\toprule
\textbf{\#} & \textbf{Script} & \textbf{Output} & \textbf{Runtime} \\
\midrule
\endhead
\bottomrule
\endfoot
1  & \texttt{crash\_analysis.py}         & \texttt{crash\_data\_statewide.csv}      & $\sim$30s \\
2  & \texttt{spatial\_features.py}       & \texttt{spatial\_features.csv}           & 3--5 min \\
3  & \texttt{environmental\_features.py} & \texttt{environmental\_features.csv}     & $\sim$10s \\
4  & \texttt{main.py}                    & \texttt{hs\_scores.csv}, 3 charts        & $\sim$20s \\
5  & \texttt{feature\_engineering.py}    & \texttt{ml\_school\_features.csv}        & $\sim$5s  \\
6  & \texttt{ml\_model.py}               & \texttt{hs\_predictor.pkl}, 3 charts     & $\sim$10s \\
7  & \texttt{seifa\_analysis.py}         & \texttt{seifa\_darebin.csv}              & $\sim$5s  \\
8  & \texttt{equity\_analysis.py}        & \texttt{chart\_equity\_seifa.png}        & $\sim$5s  \\
9  & \texttt{crash\_trend\_analysis.py}  & \texttt{chart\_crash\_trends.png}        & $\sim$5s  \\
10 & \texttt{demographics\_chart.py}     & \texttt{chart4\_demographics.png}        & $\sim$5s  \\
\end{longtable}
\end{center}

\subsection{Flags}

\begin{lstlisting}[style=bash, caption={run\_all.py usage.}]
# Run all steps, skip any whose outputs already exist
python run_all.py

# Force re-run every step from scratch
python run_all.py --force

# Start from step 5 onwards (e.g. after updating field data)
python run_all.py --from 5
\end{lstlisting}

\section{GitHub Pages Deployment}

The dashboard is deployed automatically by GitHub Pages from the \texttt{docs/}
directory on the \texttt{main} branch.  No CI/CD configuration is required beyond
the repository setting:

\begin{lstlisting}[style=text, caption={GitHub Pages setting (repository Settings > Pages).}]
Source:  Deploy from a branch
Branch:  main
Folder:  /docs
\end{lstlisting}

The complete deployment workflow after a field data update is:

\begin{lstlisting}[style=bash, caption={End-to-end update workflow.}]
# 1. Update field survey data
#    (edit school_data.csv or re-export from Google Forms)

# 2. Re-run the pipeline (skips unchanged steps automatically)
python run_all.py --from 4

# 3. Rebuild dashboard data
python prepare_data.py

# 4. Commit and push
git add docs/data/data.json docs/data/charts/
git commit -m "chore: refresh dashboard data"
git push origin main
# GitHub Pages auto-redeploys within ~60 seconds
\end{lstlisting}

\begin{tipbox}{Local development server}
Browsers block \texttt{fetch()} calls to local \texttt{file://} URIs due to CORS policy.
Always serve the dashboard through a local HTTP server during development:

\begin{lstlisting}[style=bash]
cd docs && python -m http.server 8080
\end{lstlisting}

Then open \texttt{http://localhost:8080} in a browser.
Opening \texttt{docs/index.html} directly as a file will show the loading spinner
indefinitely.
\end{tipbox}

% ============================================================
%  PART 5
% ============================================================
\part{Data Ingestion, Environmental Features, and Core Pipeline}

\begin{center}
\vspace{1cm}
{\large\itshape
Part 5 covers the three data acquisition scripts that feed the scoring engine
and the master orchestration script that drives it all.
Chapter~16 documents \texttt{crash\_analysis.py} --- how Victorian road crash
records are downloaded from the open data portal, joined across three tables,
filtered for pedestrian and cyclist involvement, and assigned to school gates
via vectorised Haversine distance.
Chapter~17 explains how \texttt{environmental\_features.py} fetches live air quality
data from EPA Victoria AirWatch and personal crime rates from the Crime Statistics
Agency, with two-tier fallback strategies so the pipeline never blocks on external APIs.
Chapter~18 documents \texttt{main.py} --- the six-step orchestrator that loads all
upstream data, applies all ten HS scoring functions, generates three publication-quality
charts, runs the recommendation engine, and writes the KDE heatmap and Leaflet map.
}
\vspace{1cm}
\end{center}

% ============================================================
\chapter{Crash Data Download and Spatial Join}
% ============================================================

\section{Overview}

\texttt{crash\_analysis.py} (pipeline Step~1) is the first script to run and the
only one that requires an internet connection.  It downloads three raw CSV tables
from the VicRoads open data portal, joins them, filters for pedestrian and cyclist
incidents in the last five years, and assigns each crash to its nearest school gate
using a vectorised Haversine computation.

\begin{keybox}{Two output files}
\begin{itemize}
  \item \texttt{outputs/crash\_data\_statewide.csv} --- all Victorian ped/cyc crashes
        (last 5 years) with a \texttt{nearest\_school} column assigning each crash
        to the closest school gate in the state.
  \item \texttt{outputs/crash\_data\_darebin.csv} --- backward-compatible subset:
        only crashes within 400\,m of the original three Darebin gates.
\end{itemize}
\end{keybox}

\section{VicRoads Open Data Schema}

VicRoads publishes road crash data as three relational tables under a single CKAN
dataset.  All three are downloaded programmatically at runtime; no manual export is needed.

\begin{center}
\begin{tabular}{lL{4cm}L{6cm}}
\toprule
Table & Key columns & Role \\
\midrule
\texttt{accident.csv} & \texttt{ACCIDENT\_NO}, \texttt{ACCIDENT\_DATE},
  \texttt{SEVERITY}, \texttt{SPEED\_ZONE} & One row per crash event \\
\texttt{node.csv} & \texttt{ACCIDENT\_NO}, \texttt{LGA\_NAME},
  \texttt{LATITUDE}, \texttt{LONGITUDE} & GPS coordinates and LGA name \\
\texttt{person.csv} & \texttt{ACCIDENT\_NO},
  \texttt{ROAD\_USER\_TYPE\_DESC} & One row per person; used to identify
  pedestrian/cyclist crashes \\
\bottomrule
\end{tabular}
\end{center}

The pedestrian/cyclist filter is applied to \texttt{person.csv} rather than
\texttt{accident.csv} because the accident-level table has no road-user-type field ---
it records the crash outcome (severity, speed zone) but not who was involved.

\section{Downloading the Three Tables}

\begin{lstlisting}[style=python, caption={fetch\_csv() and the three VicRoads URLs.}]
URLS = {
    'accident': 'https://opendata.transport.vic.gov.au/.../accident.csv',
    'person':   'https://opendata.transport.vic.gov.au/.../person.csv',
    'node':     'https://opendata.transport.vic.gov.au/.../node.csv',
}

def fetch_csv(name, url):
    r = requests.get(url, timeout=180)
    r.raise_for_status()
    df = pd.read_csv(StringIO(r.text), low_memory=False)
    df.columns = df.columns.str.strip().str.upper()
    return df

accident = fetch_csv('accident', URLS['accident'])
node     = fetch_csv('node',     URLS['node'])
person   = fetch_csv('person',   URLS['person'])
\end{lstlisting}

The 180-second timeout is set deliberately for \texttt{accident.csv} which can be
50--80\,MB.  All column names are uppercased and stripped immediately after loading to
normalise across VicRoads data releases (which occasionally change capitalisation).

\section{Table Join and Date Filter}

Node coordinates are joined to accidents on \texttt{ACCIDENT\_NO}.
\texttt{drop\_duplicates('ACCIDENT\_NO')} ensures that intersections with multiple
node records do not inflate the row count.

\begin{lstlisting}[style=python, caption={Join node to accident, filter last 5 years.}]
node_slim = node[['ACCIDENT_NO','LGA_NAME','LATITUDE',
                  'LONGITUDE','POSTCODE_CRASH']
                ].drop_duplicates('ACCIDENT_NO')

merged = accident.merge(node_slim, on='ACCIDENT_NO', how='left')

# 5-year filter
YEARS_BACK = 5
cutoff = datetime.now() - timedelta(days=365 * YEARS_BACK)
merged['ACCIDENT_DATE'] = pd.to_datetime(
    merged['ACCIDENT_DATE'], dayfirst=True, errors='coerce')
merged = merged[merged['ACCIDENT_DATE'] >= cutoff].copy()
\end{lstlisting}

\begin{warnbox}{dayfirst=True}
VicRoads date strings use DD/MM/YYYY format.  Without \texttt{dayfirst=True},
pandas interprets January dates as day 1 of month 1 correctly, but days 13--31
would raise a parsing error.  Always set this flag when reading Victorian government
date columns.
\end{warnbox}

\section{Pedestrian and Cyclist Filter}

The ped/cyc filter is applied to \texttt{person.csv} using a string search on
\texttt{ROAD\_USER\_TYPE\_DESC}.

\begin{lstlisting}[style=python, caption={Extracting ped/cyc accident IDs from person.csv.}]
ped_cyc_acc = person[
    person['ROAD_USER_TYPE_DESC'].astype(str).str.upper()
    .str.contains('PEDESTRIAN|CYCLIST|BICYCL', na=False)
]['ACCIDENT_NO'].unique()

merged['PED_OR_CYC'] = merged['ACCIDENT_NO'].isin(ped_cyc_acc)
statewide = merged[merged['PED_OR_CYC']].dropna(
    subset=['LATITUDE', 'LONGITUDE']).copy()
\end{lstlisting}

The three patterns (\texttt{PEDESTRIAN}, \texttt{CYCLIST}, \texttt{BICYCL}) are
needed because VicRoads uses both ``Cyclist'' and ``Bicyclist'' across data releases.

\section{School Gate Lookup}

Before assigning crashes to schools, \texttt{fetch\_school\_gates()} attempts to
download the full list of Victorian government secondary school coordinates from the
Department of Education and Training (DET) via the data.vic.gov.au CKAN API.

\begin{lstlisting}[style=python, caption={DET school download with fallback.}]
def fetch_school_gates():
    fallback = dict(SCHOOL_GATES)  # 3 Darebin gates always present
    try:
        pkg = requests.get(
            'https://discover.data.vic.gov.au/api/3/action/package_show',
            params={'id': 'school-locations-time-series'},
            timeout=30,
        )
        pkg.raise_for_status()
        resources = pkg.json()['result']['resources']
        # sort by name descending to get the most recent year's CSV
        csv_resources = sorted(
            [r for r in resources if r.get('format','').upper() == 'CSV'],
            key=lambda r: r.get('name',''), reverse=True
        )
        url = csv_resources[0]['url']
        r = requests.get(url, timeout=60)
        r.raise_for_status()
        schools = pd.read_csv(StringIO(r.text), low_memory=False)
        # ... filter GOVERNMENT, SECONDARY|COMBINED, STATUS=OPEN ...
        gates.update(fallback)  # Darebin gates always override
        return gates
    except Exception as e:
        print(f"  Warning: DET download failed ({e})")
        return fallback
\end{lstlisting}

The DET dataset lists all open Victorian government secondary schools with their
GPS coordinates.  Adding all state schools makes the statewide CSV useful for
future scaling --- a researcher studying any Victorian school can filter the
statewide CSV by \texttt{nearest\_school} without re-downloading crash data.

\section{Vectorised Haversine Distance}

Assigning each crash to its nearest school gate is done with a vectorised
Haversine implementation rather than a row-wise loop.  With 7,000+ crashes
and 1,500+ school gates, a naive nested loop would take hours; the vectorised
approach completes in under a minute.

\begin{lstlisting}[style=python, caption={Haversine distance (radians, returns metres).}]
def haversine_m(lat1, lon1, lat2, lon2):
    R = 6_371_000
    phi1, phi2 = np.radians(lat1), np.radians(lat2)
    dphi = np.radians(lat2 - lat1)
    dlam = np.radians(lon2 - lon1)
    a = (np.sin(dphi / 2)**2
         + np.cos(phi1) * np.cos(phi2) * np.sin(dlam / 2)**2)
    return R * 2 * np.arctan2(np.sqrt(a), np.sqrt(1 - a))
\end{lstlisting}

\begin{lstlisting}[style=python, caption={Vectorised nearest-school assignment.}]
crash_lats = statewide['LATITUDE'].to_numpy()
crash_lons = statewide['LONGITUDE'].to_numpy()

best_dist   = np.full(len(statewide), np.inf)
best_school = np.full(len(statewide), '', dtype=object)

for school, gate in all_gates.items():
    d = haversine_m(crash_lats, crash_lons, gate['lat'], gate['lon'])
    closer = d < best_dist
    best_dist[closer]   = d[closer]
    best_school[closer] = school

statewide['nearest_school'] = best_school
statewide['dist_to_gate_m'] = np.round(best_dist, 1)
\end{lstlisting}

The outer loop is over schools (1,500 iterations), and each iteration applies a
fully vectorised NumPy computation over all crashes simultaneously.
The \texttt{closer} boolean mask updates only the cells where this school is
closer than the current best, so the final result is a global nearest-school
assignment in $O(n\_schools \times n\_crashes)$ multiply-accumulate operations ---
fast enough on a laptop.

\begin{tipbox}{Darebin backward-compatible subset}
After saving the statewide CSV, a second filter produces the Darebin-only subset:
crashes where \texttt{nearest\_school} is one of the three original Darebin gates
\emph{and} \texttt{dist\_to\_gate\_m} $\leq$ 400.  This file is consumed by
\texttt{crash\_trend\_analysis.py} (Chapter~13) and \texttt{prepare\_data.py}
(Chapter~15) for the heatmap and GeoJSON crash layers.
\end{tipbox}

\section{Running crash\_analysis.py}

\begin{lstlisting}[style=bash, caption={Running the crash data downloader.}]
python crash_analysis.py
# Downloads ~3 CSVs from opendata.transport.vic.gov.au (~2-4 min)
# Outputs:
#   outputs/crash_data_statewide.csv   (~7,000 rows)
#   outputs/crash_data_darebin.csv     (400m buffer around 3 schools)
\end{lstlisting}

Once the output files exist, \texttt{run\_all.py} skips this step automatically
on subsequent runs.  Use \texttt{--force} only when you want fresh data (e.g.,
at the start of a new calendar year when new crashes have been published).

% ============================================================
\chapter{Environmental Features: AQI and Crime}
% ============================================================

\section{Why Environmental Data}

Two Healthy Streets indicators cannot be computed from field observations or OSM
alone because they depend on continuously-measured urban environment data:

\begin{itemize}
  \item \textbf{HS10 --- Clean air:} Requires PM2.5 air quality measurements.
        No field observer can reliably estimate particulate matter concentration.
  \item \textbf{HS7 --- People feel safe:} The personal-safety component is partly
        driven by the local crime rate, which the Crime Statistics Agency Victoria
        publishes annually at suburb level.
\end{itemize}

\texttt{environmental\_features.py} (pipeline Step~3) fetches both data points
for each school, computes their HS score contributions, and writes
\texttt{outputs/environmental\_features.csv}.

\section{Module Architecture}

Data fetching is separated into two modules in \texttt{src/data/}:

\begin{center}
\begin{tabular}{lL{9cm}}
\toprule
Module & Responsibility \\
\midrule
\texttt{epa\_fetcher.py} & Fetch PM2.5 from EPA Victoria AirWatch API
  (Alphington monitoring station). Two live API attempts, then suburb fallback. \\
\texttt{crime\_fetcher.py} & Fetch personal crime offence rate from
  Crime Statistics Agency Victoria XLSX. One download attempt, then suburb fallback. \\
\bottomrule
\end{tabular}
\end{center}

Both modules follow the same pattern: \emph{try live data, fall back gracefully}.
The pipeline never blocks or raises an unhandled exception because external APIs
are unavailable.

\section{EPA AirWatch --- Three-Tier Fallback}

\texttt{fetch\_epa\_aqi(suburb)} tries three data sources in order:

\begin{enumerate}
  \item \textbf{EPA Gateway API} --- the official EPA Victoria monitoring API at
        \texttt{gateway.api.epa.vic.gov.au}.  Returns the 24-hour average PM2.5
        for the Alphington station (ID \texttt{10102}), which is the closest
        monitoring station to Darebin schools.
  \item \textbf{AirWatch XLSX download} --- a 30-day rolling download from the
        older AirWatch portal.  Used when the Gateway API is rate-limited or
        returns an unexpected schema.
  \item \textbf{Suburb-level hardcoded fallback} --- 2023 annual averages for
        Reservoir (8.5), Thornbury (7.0), and Preston (7.8) $\mu$g/m\textsuperscript{3},
        derived from AirWatch historical data.
\end{enumerate}

\begin{lstlisting}[style=python, caption={EPA AirWatch three-tier fallback in epa\_fetcher.py.}]
AQI_FALLBACK = {
    'Reservoir': 8.5,   # Plenty Rd arterial corridor
    'Thornbury': 7.0,   # more residential
    'Preston':   7.8,   # Bell St / High St arterials
}

EPA_STATION_ID = '10102'  # Alphington monitoring station

def fetch_epa_aqi(suburb: str) -> float:
    # Tier 1: EPA Gateway API
    try:
        url = f'{EPA_BASE_URL}/sites/{EPA_STATION_ID}/parameters'
        resp = requests.get(url, timeout=10)
        if resp.status_code == 200:
            for param in resp.json().get('parameters', []):
                if 'PM2.5' in param.get('name', ''):
                    avg = param.get('averages',{}).get('24hour',{}).get('value')
                    if avg is not None:
                        return float(avg)
    except Exception:
        pass

    # Tier 2: AirWatch 30-day XLSX download
    try:
        # ... 30-day rolling download from AirWatch portal ...
        if vals:
            return round(float(np.mean(vals)), 1)
    except Exception:
        pass

    # Tier 3: hardcoded 2023 annual average
    return AQI_FALLBACK.get(suburb, 8.0)
\end{lstlisting}

\subsection{AQI to HS10 Score Conversion}

PM2.5 is converted to an HS10 base score using US EPA breakpoints for PM2.5 air
quality categories:

\begin{center}
\begin{tabular}{C{3cm}C{3cm}C{2.5cm}}
\toprule
PM2.5 ($\mu$g/m\textsuperscript{3}) & AQI Category & HS10 Score \\
\midrule
$\leq 12$ & Good & 10.0 \\
$\leq 35$ & Moderate & 7.0 \\
$\leq 55$ & Unhealthy for Sensitive Groups & 5.0 \\
$\leq 150$ & Unhealthy & 2.0 \\
$> 150$ & Very Unhealthy & 0.0 \\
\bottomrule
\end{tabular}
\end{center}

All three Darebin schools sit in the ``Good'' range ($\leq 12\ \mu$g/m\textsuperscript{3}),
giving a base HS10 score of 10.0.  The AQI data therefore has limited discriminating
power in this inner-urban study but becomes significant if the pipeline is extended to
industrial or high-traffic corridors (e.g.\ Coburg North, Broadmeadows).

\section{Crime Statistics Agency --- Single-Attempt Download}

\texttt{fetch\_crime\_rate(suburb)} attempts one XLSX download from the CSA Victoria
public files server, then falls back to hardcoded 2023-24 suburb estimates.

\begin{lstlisting}[style=python, caption={CSA Victoria crime rate fetcher in crime\_fetcher.py.}]
CRIME_FALLBACK = {
    'Reservoir': 820,   # higher --- arterial corridor
    'Thornbury': 560,   # lower, more residential
    'Preston':   710,   # medium --- mixed use
}

def fetch_crime_rate(suburb: str, out_dir: str = 'outputs') -> float:
    try:
        resp = requests.get(CSA_DOWNLOAD_URL, timeout=30, stream=True)
        if resp.status_code == 200:
            # Save to temp file, parse with pd.ExcelFile
            tmp = os.path.join(out_dir, '_csa_tmp.xlsx')
            with open(tmp, 'wb') as f:
                for chunk in resp.iter_content(chunk_size=8192):
                    f.write(chunk)
            xl = pd.ExcelFile(tmp)
            for sheet in xl.sheet_names:
                if 'suburb' in sheet.lower() or 'lga' in sheet.lower():
                    csa = xl.parse(sheet)
                    # ... find suburb row, return rate ...
            os.remove(tmp)
    except Exception as e:
        print(f"  Crime: CSA download failed ({type(e).__name__})")

    return CRIME_FALLBACK.get(suburb, 700)
\end{lstlisting}

The parser searches for a sheet with ``suburb'' or ``lga'' in its name, then
looks for columns whose names contain \texttt{SUBURB}/\texttt{LGA} and \texttt{RATE}.
This schema-agnostic approach is robust to minor formatting changes across annual CSA
data releases.

\subsection{Crime Rate to HS7 Score Contribution}

The personal crime rate (offences against persons per 100,000 population) is converted
to a HS7 component score (0--4 points, added inside the main HS7 formula):

\begin{center}
\begin{tabular}{C{4cm}C{3cm}}
\toprule
Crime rate (per 100k) & HS7 points \\
\midrule
$\leq 300$ & 4.0 \\
$\leq 600$ & 3.0 \\
$\leq 900$ & 2.0 \\
$\leq 1200$ & 1.0 \\
$> 1200$ & 0.0 \\
\bottomrule
\end{tabular}
\end{center}

\section{School--Suburb Mapping}

Both fetchers receive a suburb name, not a school name.  The mapping is defined
in \texttt{environmental\_features.py} and is applied before any API call:

\begin{lstlisting}[style=python, caption={School to suburb mapping.}]
SCHOOL_SUBURB_MAP = {
    'Reservoir HS':       {'suburb': 'Reservoir',  'postcode': '3073'},
    'William Ruthven SC': {'suburb': 'Thornbury',  'postcode': '3071'},
    'Preston HS':         {'suburb': 'Preston',    'postcode': '3072'},
    # ... full-name variants also included ...
}
\end{lstlisting}

\section{Output CSV Structure}

\begin{lstlisting}[style=text, caption={environmental\_features.csv columns.}]
school_name | suburb | lga | aqi_pm25 | aqi_category | hs10_base_score
crime_rate_per_100k | hs7_crime_pts | data_date
aqi_source | crime_source
\end{lstlisting}

The \texttt{aqi\_source} and \texttt{crime\_source} columns record whether live
API data or the fallback was used.  This is important for audit and reproducibility ---
a researcher replicating the study in a different year can check whether the AQI values
came from the same monitoring station or from a hardcoded historical estimate.

\begin{lstlisting}[style=bash, caption={Running the environmental features fetcher.}]
python environmental_features.py
# Tries EPA API then AirWatch then fallback for each school
# Tries CSA XLSX download then fallback for each school
# Output: outputs/environmental_features.csv
\end{lstlisting}

\begin{tipbox}{Running offline}
If you are working without internet access (e.g.\ on a train), all three fallback
values are reasonable suburb-level 2023 estimates.  The pipeline runs identically ---
only the \texttt{aqi\_source} and \texttt{crime\_source} columns will record
``Fallback'' instead of the live source name.
\end{tipbox}

% ============================================================
\chapter{Core Pipeline Orchestration}
% ============================================================

\section{Role of main.py}

\texttt{main.py} (pipeline Step~4) is the central orchestrator.  It does not download
any data --- all data must exist in \texttt{outputs/} before it runs.  Its job is:

\begin{enumerate}
  \item Load and clean the field survey data.
  \item Attach spatial and environmental features to each row.
  \item Apply all ten HS scoring functions, one row at a time.
  \item Aggregate to school-level means and save \texttt{hs\_scores.csv}.
  \item Generate three publication-quality charts (radar, 10-indicator grid, per-school bars).
  \item Run the recommendation engine and save \texttt{recommendations.csv}.
  \item Build the interactive Leaflet hazard map and the KDE heatmap PNG.
\end{enumerate}

\begin{lstlisting}[style=bash, caption={Running main.py.}]
# Run in order: crash_analysis -> spatial_features -> environmental_features -> main
python crash_analysis.py       # Step 1
python spatial_features.py     # Step 2
python environmental_features.py  # Step 3
python main.py                 # Step 4
\end{lstlisting}

\section{Data Loading Hierarchy}

\texttt{main.py} is deliberately resilient to missing upstream files.
Two optional feature sources are loaded if they exist; otherwise a fallback dict
is used so scoring can still run (with NaN for spatially-derived indicators).

\begin{lstlisting}[style=python, caption={Resilient feature loading in main.py.}]
# Spatial features (OSM-derived)
_sf = None
if SPATIAL_CSV.exists():
    _sf = pd.read_csv(SPATIAL_CSV).set_index('school_name')

def _spatial(school_full: str) -> dict:
    if _sf is None:
        return {}
    short = SCHOOL_SHORT_NAMES.get(school_full, school_full)
    for key in (short, school_full):
        if key in _sf.index:
            return _sf.loc[key].to_dict()
    return {}

# Environmental features (AQI + crime)
_ENV_FALLBACK = {
    'Reservoir High School':             {'aqi_pm25': 8.2, 'crime_rate_per_100k': 850},
    'William Ruthven Secondary College': {'aqi_pm25': 7.1, 'crime_rate_per_100k': 580},
    'Preston High School':               {'aqi_pm25': 7.8, 'crime_rate_per_100k': 720},
}

_ef = None
if ENV_CSV.exists():
    _ef = pd.read_csv(ENV_CSV).set_index('school_name')

def _env(school_full: str) -> dict:
    if _ef is not None:
        for key in (SCHOOL_SHORT_NAMES.get(school_full, school_full), school_full):
            if key in _ef.index:
                return _ef.loc[key].to_dict()
    return _ENV_FALLBACK.get(school_full, {})
\end{lstlisting}

The double-key lookup (short name first, full name second) ensures robustness
regardless of which naming convention was used in upstream CSV files.

\section{HS Scoring Loop}

All ten indicators are computed in a single \texttt{iterrows()} pass.
Each indicator function receives the survey row plus, where needed, the spatial
or environmental feature dict.

\begin{lstlisting}[style=python, caption={Per-row HS scoring loop in main.py.}]
for idx, row in df.iterrows():
    sp = _spatial(row['School'])
    ev = _env(row['School'])
    df.loc[idx, 'HS1']  = compute_hs1(row)
    df.loc[idx, 'HS2']  = compute_hs2(row)
    df.loc[idx, 'HS3']  = compute_hs3(row, sp)   # needs spatial
    df.loc[idx, 'HS4']  = compute_hs4(row, sp)
    df.loc[idx, 'HS5']  = compute_hs5(row)
    df.loc[idx, 'HS6']  = compute_hs6(row, sp)
    df.loc[idx, 'HS7']  = compute_hs7(row, ev)   # needs environmental
    df.loc[idx, 'HS8']  = compute_hs8(row, sp)
    df.loc[idx, 'HS9']  = compute_hs9(row)
    df.loc[idx, 'HS10'] = compute_hs10(row, ev, sp)

df['HS_overall'] = df[HS_CODES].mean(axis=1, skipna=True).round(1)
df['Sev_clean']  = df.apply(compute_severity, axis=1)
\end{lstlisting}

After the loop, scores are aggregated to school level with \texttt{groupby} mean:

\begin{lstlisting}[style=python, caption={School-level aggregation and CSV export.}]
school_hs = (df.groupby(['School', 'School_short'])[HS_CODES + ['HS_overall']]
               .mean().round(2).reset_index())
school_hs.to_csv(str(HS_CSV), index=False)
\end{lstlisting}

This design means multiple field survey rows per school are allowed --- the scoring
engine computes per-observation scores and the aggregation step averages them.

\section{Chart Generation}

\subsection{Chart 1 --- Radar Chart}

The first output is a polar radar chart with one polygon per school.
Ten HS indicator scores define the vertices; a dashed red ring marks the
``good threshold'' at 6.0.

\begin{lstlisting}[style=python, caption={Radar chart construction (key steps).}]
N      = len(HS_CODES)           # 10 indicators
angles = np.linspace(0, 2*np.pi, N, endpoint=False).tolist()
angles += angles[:1]             # close the polygon

for _, row in school_hs.iterrows():
    values = [row[c] if pd.notna(row[c]) else 0 for c in HS_CODES]
    values += values[:1]         # close the polygon
    ax.plot(angles, values, linewidth=2.5, label=school)
    ax.fill(angles, values, alpha=0.10)

ax.plot(angles, [6]*(N+1), color='#C0392B',
        linewidth=1.2, linestyle='--', alpha=0.6)
\end{lstlisting}

The \texttt{endpoint=False} argument to \texttt{linspace} ensures the 10 angles are
evenly spaced without repeating the first angle.  The polygon is then closed manually
by appending the first value to the end of both the angle and value lists.

\subsection{Chart 2 --- 10-Indicator Grid}

A $2 \times 5$ grid of bar charts, one subplot per HS indicator.
All subplots share a y-axis (0--12) so scores are visually comparable across
the grid.  Bar colours encode the quality band: green ($\geq 7.0$),
orange (5--6.9), red ($< 5.0$).

\begin{lstlisting}[style=python, caption={Colour encoding in the indicator grid.}]
bar_cols = ['#27AE60' if v >= 7 else
            ('#F39C12' if v >= 5 else
             ('#C0392B' if v > 0 else '#CCCCCC'))
            for v in vals]
\end{lstlisting}

\subsection{Chart 3 --- Per-School Breakdown}

One bar-chart subplot per school (side by side), each showing all 10 indicator
scores for that school.  The severity classification is rendered as a colour-coded
text badge inside each subplot's top-right corner.

\section{Visualisation Module}

The three scripts in \texttt{src/visualisation/} separate charting logic from
pipeline orchestration:

\begin{center}
\begin{tabular}{lL{9.5cm}}
\toprule
Module & Contents \\
\midrule
\texttt{charts.py} & \texttt{plot\_safety\_scores()}, \texttt{plot\_hazard\_severity()},
  \texttt{plot\_score\_breakdown()}, \texttt{plot\_demographics()} ---
  static PNG generation via matplotlib \\
\texttt{heatmap.py} & \texttt{build\_kde\_heatmap()} --- kernel density estimate
  heatmap of observation points; outputs both a static PNG (\texttt{heatmap.png})
  for QGIS import and an interactive \texttt{map\_heatmap.html} via Folium \\
\texttt{interactive\_map.py} & \texttt{build\_interactive\_map()} --- Folium-based
  hazard map with colour-coded observation markers, popup cards showing HS scores
  and recommendations, and a severity legend \\
\texttt{qgis\_export.py} & Exports observations as a GeoPackage for GIS use \\
\bottomrule
\end{tabular}
\end{center}

\subsection{Interactive Map}

The Folium interactive map (\texttt{outputs/map\_interactive.html}) is generated by
\texttt{build\_interactive\_map(df, rec\_df, out\_dir)}.  Each observation location
is a circle marker coloured by severity (red = Major, orange = Moderate, green = Minor).
Clicking a marker opens a popup with the school name, location, HS scores, and a list
of applicable recommendations.

\subsection{KDE Heatmap}

\texttt{build\_kde\_heatmap(df, rec\_df, out\_dir)} generates a kernel density estimate
of observation density.  The static \texttt{heatmap.png} uses a
\texttt{matplotlib} KDE over the observation GPS coordinates and can be imported
directly into QGIS as a georeferenced raster layer.

\section{Graceful Degradation Pattern}

\texttt{main.py} wraps the map generation steps in try/except blocks:

\begin{lstlisting}[style=python, caption={Graceful degradation for optional outputs.}]
try:
    out_map = build_interactive_map(df, rec_df, OUT_DIR_STR)
    print(f"      Saved -> {out_map}")
except Exception as e:
    print(f"      Warning: map failed -- {e}")

try:
    build_kde_heatmap(df, rec_df, OUT_DIR_STR)
except Exception as e:
    print(f"      Warning: heatmap failed -- {e}")
\end{lstlisting}

The HS scores and recommendations CSVs are always saved first, then the optional
visual outputs.  If Folium is not installed or a map rendering fails, the core
data outputs are unaffected.  This makes \texttt{main.py} safe to run in headless
or minimal-dependency environments (e.g.\ a GitHub Actions runner).

\section{Expected Outputs}

After \texttt{main.py} completes successfully:

\begin{center}
\begin{tabular}{lL{8cm}}
\toprule
File & Description \\
\midrule
\texttt{outputs/hs\_scores.csv}           & School-level HS1--HS10 means and overall score \\
\texttt{outputs/recommendations.csv}      & Ranked recommendations from all 17 rules \\
\texttt{outputs/chart1\_safety\_scores.png}  & HS radar chart (10-indicator polygon per school) \\
\texttt{outputs/chart2\_hazard\_severity.png} & $2 \times 5$ indicator grid with colour bands \\
\texttt{outputs/chart3\_score\_breakdown.png} & Per-school indicator bar charts \\
\texttt{outputs/map\_interactive.html}    & Folium hazard map (open in browser) \\
\texttt{outputs/map\_heatmap.html}        & KDE heatmap (Folium, open in browser) \\
\texttt{outputs/heatmap.png}              & KDE heatmap PNG (import into QGIS) \\
\bottomrule
\end{tabular}
\end{center}

% ============================================================
\part{Feature Engineering, Socio-economic Analysis, and Demographics}
% ============================================================

% ============================================================
\chapter{Feature Engineering for Machine Learning}
% ============================================================

\section{Purpose}

\texttt{feature\_engineering.py} is Step~5 in the master pipeline runner.  It assembles
a single school-level feature matrix that the Ridge regression model (Step~6) trains on
and the scenario engine uses at inference time.

The core design decision is that \emph{every feature must come from freely available open
data} --- OSM, EPA AirWatch, CSA Victoria, and VicRoads --- so the pipeline can be applied
to any Victorian school without field surveys.  The HS indicator scores from \texttt{main.py}
act as the supervised targets ($y$); the open-data features are the predictors ($X$).

\begin{keybox}{Output file}
\texttt{outputs/ml\_school\_features.csv} --- one row per school, columns:
\texttt{school} + 26~feature columns + \texttt{HS1}--\texttt{HS10} + \texttt{HS\_overall}.
\end{keybox}

\section{Feature Groups}

The feature matrix is built from three source CSVs produced by earlier pipeline steps.

\subsection{Spatial Features (OSM)}

Twenty columns derived from \texttt{outputs/spatial\_features.csv} (Step~2).  Each column
measures a physical aspect of the street environment within a 200\,m or 400\,m buffer
around the school gate.

\begin{center}
\begin{tabular}{llL{5cm}}
\toprule
Feature column & Buffer & HS proxy \\
\midrule
\texttt{footpath\_pct\_200m} / \texttt{footpath\_pct\_400m} & 200 / 400\,m & HS1 (pedestrian infrastructure) \\
\texttt{crossings\_400m}, \texttt{signals\_400m}, \texttt{crossing\_density\_400m} & 400\,m & HS2 (easy to cross) \\
\texttt{tree\_count\_100m}, \texttt{shelter\_count\_200m}, \texttt{green\_pct\_400m} & 100--400\,m & HS3 (shade \& shelter) \\
\texttt{bench\_count\_200m}, \texttt{park\_count\_400m} & 200--400\,m & HS4 (rest places) \\
\texttt{avg\_speed\_400m}, \texttt{arterial\_pct\_400m}, \texttt{high\_speed\_road\_400m}, \texttt{road\_count\_400m} & 400\,m & HS5 + HS9 (traffic stress) \\
\texttt{cycle\_pct\_400m}, \texttt{protected\_cycle\_length\_400m}, \texttt{pt\_stops\_400m} & 400\,m & HS6 (active travel choice) \\
\texttt{amenity\_count\_400m}, \texttt{cafe\_count\_400m} & 400\,m & HS8 (things to do) \\
\texttt{walk\_length\_400m} & 400\,m & context / HS1 \\
\bottomrule
\end{tabular}
\end{center}

\subsection{Environmental Features}

Two columns from \texttt{outputs/environmental\_features.csv} (Step~3):

\begin{center}
\begin{tabular}{lL{8cm}}
\toprule
Column & HS proxy \\
\midrule
\texttt{crime\_rate\_per\_100k} & HS7 (feel safe from crime) \\
\texttt{aqi\_pm25} & HS10 (clean air) \\
\bottomrule
\end{tabular}
\end{center}

\subsection{Crash Statistics}

Four aggregate columns computed on the fly from
\texttt{outputs/crash\_data\_statewide.csv} (Step~1):

\begin{center}
\begin{tabular}{lL{8cm}}
\toprule
Column & Derivation \\
\midrule
\texttt{crash\_count} & Count of ped/cyclist crashes within 500\,m of the school gate \\
\texttt{serious\_or\_fatal\_rate} & Fraction of crashes with SEVERITY $\in \{1,2\}$ \\
\texttt{school\_hours\_pct} & Fraction of crashes occurring on weekdays 07:00--09:00 or 14:00--17:00 \\
\texttt{avg\_speed\_zone} & Mean posted speed zone, excluding sentinel codes 777 / 888 / 999 \\
\bottomrule
\end{tabular}
\end{center}

\section{Speed-Zone Sentinel Codes}

VicRoads encodes unknown or special speed limits as sentinel integers above 200:
777 (not known), 888 (road works), 999 (other).  The script filters these out before
computing the mean:

\begin{lstlisting}[style=python, caption={Sentinel code filtering in crash aggregation.}]
avg_speed_zone = ('SPEED_ZONE',
    lambda x: pd.to_numeric(x, errors='coerce')
                .where(lambda v: v < 200)   # drop 777/888/999
                .mean()),
\end{lstlisting}

Without this filter, a single sentinel value would inflate the mean by hundreds of km/h.

\section{School-Hours Crash Fraction}

The \texttt{school\_hours\_pct} feature is computed using \texttt{pd.to\_datetime} with
\texttt{dayfirst=True} (Victorian date format DD/MM/YYYY) and \texttt{Series.dt.dayofweek}
(0 = Monday, 4 = Friday):

\begin{lstlisting}[style=python, caption={School-hours flag derivation in feature\_engineering.py.}]
cr['is_school_hours'] = (
    (pd.to_datetime(cr['ACCIDENT_DATE'], dayfirst=True,
                    errors='coerce').dt.dayofweek < 5) &
    (cr['hour'].between(7, 9) | cr['hour'].between(14, 17))
).astype(int)
\end{lstlisting}

\section{Four-Way Join}

All four source tables are normalised to a common \texttt{school} key (short school name)
and joined sequentially:

\begin{lstlisting}[style=python, caption={Four-way left-join to build the feature matrix.}]
ml = (sf_sel                           # 20 spatial features
      .merge(ef_sel,    on='school', how='left')   # +2 env
      .merge(crash_agg, on='school', how='left')   # +4 crash
      .merge(hs_sel,    on='school', how='left'))  # +11 targets
\end{lstlisting}

\texttt{how='left'} ensures that spatial features (the widest source) are never dropped.
If an environmental or crash CSV is missing a school row, those columns fill with
\texttt{NaN}, which Ridge handles via median imputation in the \texttt{ml\_model.py}
pipeline.

\begin{tipbox}{Name normalisation}
\texttt{spatial\_features.csv} uses short names (\texttt{Reservoir HS}),
\texttt{environmental\_features.csv} uses full names (\texttt{Reservoir High School}).
The \texttt{FULL\_TO\_SHORT} dict in the script normalises all four tables to short names
before the join --- a critical step that is easy to miss when extending the pipeline to
new schools.
\end{tipbox}

\section{Column Order Convention}

The output CSV always places columns in this order:
\texttt{school} $\to$ spatial features $\to$ env features $\to$ crash features $\to$
HS targets (\texttt{HS1}--\texttt{HS10}, \texttt{HS\_overall}).
The ML model reads \texttt{features} as the list of non-target columns saved into the
\texttt{.pkl} bundle, so the CSV column order is not load-bearing --- but keeping it
consistent makes visual inspection easier.

\section{Console Summary}

After saving, \texttt{feature\_engineering.py} prints a three-level feature count and the
full feature--indicator mapping to stdout:

\begin{lstlisting}[style=bash, caption={Running feature engineering.}]
python feature_engineering.py
# [1/4] Loading spatial features...
#   3 schools x 20 spatial features
# [2/4] Loading environmental features...
#   3 schools x 2 environmental features
# [3/4] Aggregating crash statistics...
#   Crash aggregates: <table>
# [4/4] Loading HS indicator targets...
#   3 schools x 10 HS targets
# School-level feature matrix saved -> outputs/ml_school_features.csv
# Features: 26  (Spatial: 20 | Environmental: 2 | Crash stats: 4)
# Targets:  10  (HS1-HS10)
\end{lstlisting}

% ============================================================
\chapter{SEIFA Disadvantage Analysis}
% ============================================================

\section{Purpose}

\texttt{seifa\_analysis.py} is Step~7 in the master pipeline.  It downloads the ABS
SEIFA 2021 SA1-level dataset, filters for the City of Darebin, performs a
population-weighted aggregation to each school catchment, and writes two output files:

\begin{center}
\begin{tabular}{lL{8cm}}
\toprule
File & Contents \\
\midrule
\texttt{outputs/seifa\_darebin.csv} & One row per school --- IRSD score, decile,
  disadvantage label, catchment population \\
\texttt{outputs/seifa\_darebin\_sa1.csv} & Full SA1-level Darebin data used for the
  catchment aggregation \\
\bottomrule
\end{tabular}
\end{center}

\begin{keybox}{SEIFA / IRSD definitions}
\textbf{SEIFA} (Socio-Economic Indexes for Areas) is published by the ABS after each
Census.  \textbf{IRSD} (Index of Relative Socio-economic Disadvantage) is the most
commonly cited index.  The national average score is approximately 1000; lower values
indicate greater disadvantage.  Decile 1 = most disadvantaged 10\,\% nationally;
Decile 10 = least disadvantaged.
\end{keybox}

\section{ABS Download and Dynamic Column Discovery}

The SEIFA 2021 SA1 XLSX is downloaded directly from the ABS website.  Because ABS has
changed column headings between Census editions, the script discovers the relevant
columns dynamically rather than hard-coding column names:

\begin{lstlisting}[style=python, caption={Schema-agnostic column discovery in seifa\_analysis.py.}]
col_map = {}
for col in seifa_raw.columns:
    col_lower = col.lower()
    if 'sa1' in col_lower or '2021 code' in col_lower:
        col_map['SA1_CODE'] = col
    elif 'irsd' in col_lower and 'score' in col_lower:
        col_map['IRSD_SCORE'] = col
    elif 'irsd' in col_lower and 'decile' in col_lower:
        col_map['IRSD_DECILE'] = col
    elif 'lga' in col_lower and 'name' in col_lower:
        col_map['LGA_NAME'] = col
    elif 'usual' in col_lower or 'population' in col_lower:
        col_map['POPULATION'] = col
df_seifa = seifa_raw.rename(columns={v: k for k, v in col_map.items()})
\end{lstlisting}

This pattern also makes it safe to re-run the pipeline when ABS releases SEIFA 2026 with
updated headings.

\section{Darebin Filtering}

Filtering prefers the LGA name column when present; it falls back to a hard-coded list of
16 SA2 codes that correspond to the City of Darebin LGA:

\begin{lstlisting}[style=python, caption={LGA-name vs SA2-code fallback filter.}]
if 'LGA_NAME' in df_seifa.columns:
    darebin_df = df_seifa[
        df_seifa['LGA_NAME'].str.contains('Darebin', case=False, na=False)
    ].copy()
elif 'SA2_CODE' in df_seifa.columns:
    darebin_df = df_seifa[
        df_seifa['SA2_CODE'].astype(str).str[:9]
        .astype(float, errors='ignore')
        .isin(DAREBIN_SA2_CODES)
    ].copy()
\end{lstlisting}

\section{Curated Fallback Data}

If the ABS download fails (network error, site maintenance, URL change), the script
falls back to an 8-row curated dataset with 2021 SEIFA values for Reservoir, Preston,
Northcote, Thornbury, and Fairfield:

\begin{lstlisting}[style=python, caption={Curated fallback snippet (Reservoir rows).}]
darebin_df = pd.DataFrame([
    {'SA1_CODE': '2060417', 'SUBURB': 'Reservoir',
     'IRSD_SCORE': 981,  'IRSD_DECILE': 4,
     'POPULATION': 4120, 'LGA_NAME': 'Darebin'},
    {'SA1_CODE': '2060418', 'SUBURB': 'Reservoir',
     'IRSD_SCORE': 975,  'IRSD_DECILE': 4,
     'POPULATION': 3980, 'LGA_NAME': 'Darebin'},
    ...
])
\end{lstlisting}

The fallback values are sourced from ABS Census 2021 DataPacks downloaded offline during
project development and are accurate to $\pm$5 IRSD points.

\section{Population-Weighted Catchment Average}

For each school, the suburb-matching SA1 rows are identified and a
population-weighted IRSD is computed:

\begin{lstlisting}[style=python, caption={Population-weighted IRSD for each school catchment.}]
pop   = pd.to_numeric(school_sa1s['POPULATION'],  errors='coerce').fillna(1)
score = pd.to_numeric(school_sa1s['IRSD_SCORE'],  errors='coerce').fillna(1000)
avg_irsd   = round(float(np.average(score, weights=pop)), 1)
avg_decile = round(float(pd.to_numeric(
    school_sa1s.get('IRSD_DECILE', pd.Series([5])), errors='coerce'
).mean()), 1)
\end{lstlisting}

Using \texttt{np.average} with population weights ensures that larger SA1 areas
contribute proportionally more to the school catchment score than smaller ones.
\texttt{fillna(1)} on \texttt{pop} prevents division-by-zero if population data is
missing; \texttt{fillna(1000)} on \texttt{score} defaults to the national mean.

\section{Disadvantage Labels}

A four-band labelling scheme converts the numeric decile into a human-readable level
for the Regen Melbourne report:

\begin{center}
\begin{tabular}{cll}
\toprule
IRSD Decile & Label & Implication \\
\midrule
$\le 3$ & High disadvantage & Many families rely on walking/cycling --- safe streets critical \\
$4$--$5$ & Moderate-high disadvantage & Active travel dependency above average \\
$6$--$7$ & Moderate disadvantage & Mixed profile --- safe streets benefit all \\
$\ge 8$ & Low disadvantage & Lower disadvantage --- but safety still essential \\
\bottomrule
\end{tabular}
\end{center}

The three Darebin schools fall in the Moderate-high band:
\textbf{Reservoir HS} (Decile~4), \textbf{William Ruthven SC} (Decile~6),
\textbf{Preston HS} (Decile~7).

\section{School Gate Coordinates}

The script also embeds the school gate GPS coordinates directly in the output CSV
(\texttt{Gate\_Lat}, \texttt{Gate\_Lon}).  These are used by \texttt{equity\_analysis.py}
(Step~8) when overlaying SEIFA data on the map and by \texttt{prepare\_data.py} when
building the dashboard \texttt{schools} JSON array.

\begin{lstlisting}[style=python, caption={School gate definitions in seifa\_analysis.py.}]
SCHOOL_GATES = {
    'Reservoir High School': {
        'lat': -37.7224, 'lon': 145.0294,
        'suburb': 'Reservoir', 'lga': 'Darebin'
    },
    'William Ruthven Secondary College': {
        'lat': -37.69654, 'lon': 145.00299,
        'suburb': 'Reservoir', 'lga': 'Darebin'
    },
    'Preston High School': {
        'lat': -37.7417, 'lon': 145.0071,
        'suburb': 'Preston', 'lga': 'Darebin'
    },
}
\end{lstlisting}

\begin{warnbox}{Note on suburb assignment for William Ruthven SC}
William Ruthven Secondary College is located in Reservoir but its street address straddles
the Reservoir/Thornbury boundary.  The SEIFA join uses \texttt{suburb='Reservoir'} (the
official VEC enrolment catchment suburb) for consistency with the school gate coordinates.
\end{warnbox}

\section{Key Finding}

The script prints a summary identifying the most disadvantaged catchment:

\begin{lstlisting}[style=bash, caption={SEIFA analysis console summary.}]
python seifa_analysis.py
# SEIFA 2021 -- Disadvantage Scores by School Catchment
# Reservoir High School
#   IRSD Score     : 974.7  (national avg = 1000)
#   IRSD Decile    : 4 / 10
#   Disadvantage   : Moderate-high disadvantage
#   Catchment pop  : ~12,310
# ...
# KEY FINDING: Most disadvantaged catchment: Reservoir High School
# IRSD Decile 4 -- correlates with worst HS scores in the pilot study.
\end{lstlisting}

This finding directly supports the equity hypothesis tested in Chapter~14: the school with
the lowest SEIFA decile also has the worst pedestrian infrastructure scores.

% ============================================================
\chapter{ABS Demographics Chart}
% ============================================================

\section{Purpose}

\texttt{demographics\_chart.py} is Step~10 --- the final step in the master pipeline.
It reads a static ABS Census 2021 summary file shipped with the repository and generates
\texttt{outputs/chart4\_demographics.png}: a four-panel bar chart comparing Reservoir
and Preston suburbs across four socioeconomic indicators.

Unlike most other pipeline steps, this script performs no data download.
\texttt{demographics\_darebin.csv} is pre-curated from ABS Census 2021 TableBuilder
exports for two Suburb of Enumeration (SAL) codes:
\textbf{Reservoir SAL22161} and \textbf{Preston SAL22121}.

\section{Script Structure}

The script is deliberately thin --- all chart logic lives in
\texttt{src/visualisation/charts.py}:

\begin{lstlisting}[style=python, caption={demographics\_chart.py --- thin entry point.}]
from config import DEMO_FILE, OUT_DIR
from src.visualisation.charts import plot_demographics

if not os.path.exists(str(DEMO_FILE)):
    raise FileNotFoundError(
        "demographics_darebin.csv not found.\n"
        "This file ships with the project repository."
    )

out = plot_demographics(str(DEMO_FILE), str(OUT_DIR))
print(f"\n  Chart saved -> {out}")
\end{lstlisting}

Raising \texttt{FileNotFoundError} (rather than falling back to dummy data) is the
correct choice here: \texttt{demographics\_darebin.csv} is a static asset that cannot be
regenerated automatically, so a missing file indicates a repository checkout problem that
the user must fix before continuing.

\section{The Four Panels}

\texttt{plot\_demographics()} reads the CSV and maps each row to one of four bar panels:

\begin{center}
\begin{tabular}{lL{8cm}}
\toprule
Panel & CSV column \\
\midrule
Median Weekly Household Income (\$) & \texttt{Median weekly household income (\$)} \\
\% Households With No Car & \texttt{\% households no car} \\
\% Using Public Transport to Work & \texttt{\% public transport to work} \\
\% Working Full Time (35\,hrs+) & \texttt{\% working full time (35hrs+)} \\
\bottomrule
\end{tabular}
\end{center}

Each panel uses a distinct colour (\texttt{\#1A5276} navy, \texttt{\#C0392B} red,
\texttt{\#D35400} orange, \texttt{\#1E8449} green) and annotates the exact value above
each bar.  The four-panel layout on a single \texttt{14 \texttimes{} 5} inch figure
makes cross-suburb comparison immediate.

\section{School-Age Children Estimate}

In addition to saving the chart, \texttt{plot\_demographics()} computes and prints a
derived metric not present as a column in the raw CSV:

\begin{lstlisting}[style=python, caption={School-age children estimate printed to console.}]
for _, row in demo_df.iterrows():
    children = round(row['Total population']
                     * row['% children aged 5-17'] / 100)
    print(f"  {row['Suburb']}:")
    print(f"    ~{children:,} school-age children")
    print(f"    ${row['Median weekly household income ($)']:,}/week median income")
    print(f"    {row['% households no car']}% households have no car")
\end{lstlisting}

This estimate --- total suburb population multiplied by the fraction aged 5--17 ---
quantifies the direct beneficiary pool for any school street safety intervention.
It contextualises the HS scores: a low HS score at a school serving 3,000+ local
children represents a larger aggregate harm than the same score at a smaller school.

\section{Connection to the Equity Argument}

The demographics chart is the final piece of the Regen Melbourne equity argument:

\begin{enumerate}
  \item \textbf{Chapter~14 (equity analysis):} SEIFA decile correlates with HS score
    (Pearson $r \approx 0.84$).
  \item \textbf{Chapter~20 (SEIFA analysis):} Reservoir HS sits in Decile~4 --- the
    most disadvantaged catchment in the pilot.
  \item \textbf{Chapter~21 (demographics):} The same suburb has the highest rate of
    car-free households and the lowest median income, meaning its residents are most
    dependent on walking and public transport --- and therefore most exposed to the
    hazards that the low HS scores identify.
\end{enumerate}

A developer extending the pipeline to additional suburbs should add a row per suburb to
\texttt{demographics\_darebin.csv} and re-run Step~10 to regenerate the chart.

\section{Pipeline Position}

Step~10 is the last step in \texttt{run\_all.py} for a reason: it depends on no output
from any other step.  \texttt{demographics\_darebin.csv} is static, and the chart is
purely presentational.  Placing it last means pipeline failures in Steps~1--9 do not
prevent the demographics chart from being generated if that step alone is re-run with
\texttt{--from 10}.

\begin{lstlisting}[style=text, caption={Step 10 in run\_all.py STEPS list.}]
{
  'n':      10,
  'script': 'demographics_chart.py',
  'label':  'ABS demographics chart',
  'checks': ['outputs/chart4_demographics.png'],
  'note':   'Static ABS Census 2021 data (no download).',
}
\end{lstlisting}

% ============================================================
\part{Synthesis, Extensions, and Conclusions}
% ============================================================

% ============================================================
\chapter{Pilot Study Results and Policy Implications}
% ============================================================

\section{Overview}

The 300,000 Streets pilot study assessed three secondary schools in the City of Darebin
during 2024.  Field surveyors and Google Street View assessors collected observations at
15--20 street segments per school, recording 28 attributes per location using the
standardised survey instrument described in Chapter~6.  The scoring engine computed
HS1--HS10 for each observation row; \texttt{main.py} aggregated these to school-level
means and saved \texttt{outputs/hs\_scores.csv}.

Table~\ref{tab:hs-scores} presents the school-level HS indicator scores from the pilot
study.  Scores below 4.0 are highlighted in red (crisis range), scores 4.0--5.9 in
amber (deficient), and scores at or above 6.0 --- the Healthy Streets ``good''
threshold --- in green.

\begin{table}[ht]
\centering
\caption{Pilot study HS indicator scores by school (0--10 scale).
         The ``good'' threshold is $\ge 6.0$ (green).}
\label{tab:hs-scores}
\setlength{\tabcolsep}{5pt}
\small
\begin{tabular}{lL{4.5cm}ccc}
\toprule
Code & Indicator & Reservoir HS & W.\ Ruthven SC & Preston HS \\
     &           & (SEIFA D4)   & (SEIFA D6)     & (SEIFA D7) \\
\midrule
HS1  & Pedestrians from all walks of life
     & \cellcolor{amber!40}4.2
     & \cellcolor{amber!40}5.5
     & \cellcolor{green!25}6.1 \\
HS2  & Easy to cross
     & \cellcolor{amber!40}5.7
     & \cellcolor{green!25}6.2
     & \cellcolor{green!25}7.0 \\
HS3  & Shade and shelter
     & \cellcolor{red!30}3.8
     & \cellcolor{amber!40}4.8
     & \cellcolor{amber!40}5.8 \\
HS4  & Places to stop and rest
     & \cellcolor{amber!40}4.0
     & \cellcolor{amber!40}5.2
     & \cellcolor{amber!40}5.9 \\
HS5  & Not too noisy / traffic stress
     & \cellcolor{amber!40}4.5
     & \cellcolor{amber!40}5.4
     & \cellcolor{green!25}6.2 \\
HS6  & People choose to walk / cycle / PT
     & \cellcolor{red!30}2.1
     & \cellcolor{red!30}3.2
     & \cellcolor{amber!40}4.1 \\
HS7  & People feel safe
     & \cellcolor{amber!40}5.2
     & \cellcolor{green!25}6.1
     & \cellcolor{green!25}6.8 \\
HS8  & Things to see and do
     & \cellcolor{amber!40}5.8
     & \cellcolor{green!25}6.5
     & \cellcolor{green!25}7.1 \\
HS9  & People feel relaxed
     & \cellcolor{amber!40}4.3
     & \cellcolor{amber!40}5.8
     & \cellcolor{green!25}6.5 \\
HS10 & Clean air (PM2.5)
     & \cellcolor{green!25}6.5
     & \cellcolor{green!25}7.2
     & \cellcolor{green!25}7.5 \\
\midrule
     & \textbf{Overall (mean HS1--HS10)}
     & \textbf{4.6}
     & \textbf{5.6}
     & \textbf{6.3} \\
     & \textit{Severity classification}
     & \textit{Major}
     & \textit{Moderate}
     & \textit{Moderate} \\
\bottomrule
\end{tabular}
\end{table}

\begin{warnbox}{Representative values}
The scores in Table~\ref{tab:hs-scores} are representative of the 2024 pilot study run.
The authoritative values for any given execution are in
\texttt{outputs/hs\_scores.csv} and will reflect the version of \texttt{school\_data.csv}
present at run time.
\end{warnbox}

\section{School-by-School Analysis}

\subsection{Reservoir High School --- Overall 4.6 / 10 --- Major}

Reservoir HS is the most disadvantaged catchment in the pilot (SEIFA Decile~4) and
records the lowest overall Healthy Streets score.  The severity classification is
\textbf{Major}, triggered by HS2 (crossings, 5.7 --- below the 4.0 crisis threshold
for this indicator would have been needed; instead it is close) and HS6
(active travel, 2.1 --- well below the 3.0 crisis threshold that flags a Major
severity).

The standout deficiency is \textbf{HS6 at 2.1} --- reflecting the near-total absence
of protected cycling infrastructure along 855 Plenty Road and surrounding arterials.
Plenty Road is a VicRoads arterial with a posted 60\,km/h speed limit, four lanes, and
frequent heavy vehicles.  Students who wish to cycle to school must share unprotected
road space with 60\,km/h traffic.

\textbf{HS3} (shade and shelter, 3.8) is below the crisis threshold --- the school
frontage has minimal tree canopy, no covered bus shelter at the main gate, and exposed
waiting areas that are unusable in rain.

HS10 (clean air, 6.5) is the only indicator above the good threshold, reflecting
Reservoir's position downwind of Alphington monitoring station
(Station 10102, mean PM2.5 $\approx$ 8.5\,$\mu$g/m$^3$) --- moderate but not severe.

\subsection{William Ruthven Secondary College --- Overall 5.6 / 10 --- Moderate}

William Ruthven SC sits at Decile~6 and scores 5.6 overall --- a full point above
Reservoir HS but still below the good threshold.  Severity is \textbf{Moderate}: no
single indicator falls into the crisis range, but eight of ten indicators are below 6.0,
satisfying the ``2+ indicators below good threshold'' rule in \texttt{severity.py}.

The weakest indicator is again \textbf{HS6} (3.2), reflecting the on-road painted bike
lanes on Merrilands Road --- present on paper, absent in practice.  A painted lane with
no physical protection adjacent to parallel parking is rated lower than ``no cycling
infrastructure at all'' on a quiet residential street.

\textbf{HS3} (4.8) and \textbf{HS4} (5.2) are the next weakest, consistent with the
inner-suburban character of the site: streets are functional but do not provide the
amenity (tree canopy, seating, shelter) that makes walking comfortable.

\subsection{Preston High School --- Overall 6.3 / 10 --- Moderate}

Preston HS sits at Decile~7 and achieves the highest overall score (6.3), with four
indicators above the good threshold (HS2, HS5, HS7, HS8, HS9, HS10).  Severity is
\textbf{Moderate}: HS6 at 4.1 and HS4 at 5.9 keep several indicators below 6.0.

The relative strength of Preston HS reflects its location on Cooma Street --- a local
road with lower traffic volumes than Plenty Road or Merrilands Road --- and the
higher street-tree canopy cover in the surrounding residential precinct.

\textbf{HS6} remains the pilot-wide weakest indicator even at its best-performing school.
At 4.1, protected cycling infrastructure exists on nearby Gilbert Road (a Victorian
Government ``Principal Bicycle Network'' route) but does not extend to the school gate
directly.  This is a systematic gap across all three schools: the cycling network ends
before it reaches the school entrance.

\section{The Equity Finding}

The Pearson correlation between SEIFA Decile and overall HS score, computed by
\texttt{equity\_analysis.py} across the three pilot schools, is $r \approx 0.84$.

The direction is unambiguous: \emph{the more socioeconomically disadvantaged the school
catchment, the worse the pedestrian and cycling environment around the school gate.}

This matters for two compounding reasons:

\begin{enumerate}
  \item Students in more disadvantaged catchments are \emph{more likely to walk or
    cycle} to school --- car ownership is lower and family transport budgets are tighter.
    The ABS Census 2021 data (Chapter~21) shows Reservoir has the highest rate of
    car-free households and the lowest median household income in the pilot.
  \item These students are therefore more exposed to the worse infrastructure: they
    depend on the very environment that is most deficient.
\end{enumerate}

This is the core policy case for the 300,000 Streets initiative: improving the street
environment around Reservoir High School is not just a road-safety intervention --- it
is an equity intervention that disproportionately benefits students who have no
alternative to walking.

\section{Top Recommendations by School}

Table~\ref{tab:top-recs} lists the three highest-priority recommendations per school,
as generated by the 17-rule engine in \texttt{src/recommendations/rules.py} (Chapter~10)
and ranked by the priority + cost heuristic in \texttt{src/recommendations/engine.py}.

\begin{table}[ht]
\centering
\caption{Top three priority recommendations per school.}
\label{tab:top-recs}
\begin{tabular}{llL{6cm}l}
\toprule
School & HS target & Recommendation & Cost band \\
\midrule
Reservoir HS  & HS6 & Install kerb-protected bike lane along Plenty Rd school frontage (200\,m) & High \\
              & HS3 & Install covered shelter at main gate and bus stop & Low \\
              & HS5 & Advocate for 40\,km/h school zone on Plenty Rd & Low \\
\midrule
W.\ Ruthven SC & HS6 & Replace painted lane with wand-separated lane on Merrilands Rd & Medium \\
               & HS3 & Plant street trees along school frontage (12+ trees) & Low \\
               & HS4 & Install seating / waiting area at main gate & Low \\
\midrule
Preston HS    & HS6 & Extend Gilbert Rd bike lane to school gate on Cooma St & Medium \\
              & HS4 & Install additional seating on walking approach routes & Low \\
              & HS1 & Widen footpath on Cooma St north of gate to 1.8\,m & Medium \\
\bottomrule
\end{tabular}
\end{table}

\section{Implications for the 300,000 Streets Initiative}

The pilot study covers three schools and $\approx$\,10,000 students.  The 300,000
Streets initiative targets every school in metropolitan Melbourne --- approximately
700 schools and 300,000 students.  Extrapolating the pilot findings to this scale
has two implications.

\textbf{HS6 is the system-wide gap.}  All three pilot schools score below the good
threshold for active travel infrastructure regardless of their socioeconomic context.
This suggests the cycling network systematically fails to connect to school gates across
the city --- a structural gap that no individual school can fix, and that requires
coordinated action from councils and the Victorian Department of Transport and Planning.

\textbf{Equity-weighted investment is justified.}  If Darebin Council must prioritise
a single school for street improvement funding, the SEIFA-weighted HS scores provide an
evidence base: Reservoir HS has the worst infrastructure and serves the most
disadvantaged community.  A 40\,km/h school zone on Plenty Road and a covered shelter
at the main gate are low-cost interventions that would immediately improve HS5 and HS3
for a school where 30\,\% of students have no access to a car.

% ============================================================
\chapter{Extending the Pipeline to New Schools and LGAs}
% ============================================================

\section{Design for Reuse}

The pipeline was architected from the outset to be school-agnostic.  No school-specific
logic is hard-coded in any of the ten pipeline steps: all school names, coordinates, and
suburb mappings are stored in configuration files and dictionaries.  Adding a new school
requires five targeted changes; adding an entire LGA (20--30 schools) requires the same
five changes at a larger scale.

\section{Five-Step Checklist for a New School}

\begin{enumerate}

\item \textbf{Add school gate coordinates to \texttt{config.py}.}

\begin{lstlisting}[style=python, caption={config.py additions for a new school.}]
SCHOOL_GATES['Melbourne HS'] = {
    'lat': -37.8276, 'lon': 144.9730,
    'addr': '8 Cemetery Rd E, Melbourne VIC 3004',
}
SCHOOL_SHORT_NAMES['Melbourne High School'] = 'Melbourne HS'
SCHOOL_GATES_BY_SHORT['Melbourne HS'] = {
    'lat': -37.8276, 'lon': 144.9730,
}
\end{lstlisting}

\item \textbf{Collect field survey data and add rows to \texttt{school\_data.csv}.}
  Each row is one surveyed street segment.  The minimum required columns are those
  in \texttt{COLUMN\_RENAME} in \texttt{config.py}: School, Latitude, Longitude,
  FAS, CSS, EEI, Severity (plus whichever HS indicator inputs your team collects).

\item \textbf{Add the school's suburb to the environmental features map.}

\begin{lstlisting}[style=python, caption={environmental\_features.py suburb mapping addition.}]
SCHOOL_SUBURB_MAP['Melbourne HS'] = {
    'suburb'  : 'South Yarra',
    'postcode': '3004',
}
\end{lstlisting}

\item \textbf{Add a school entry to \texttt{seifa\_analysis.py}.}

\begin{lstlisting}[style=python, caption={seifa\_analysis.py school gate addition.}]
SCHOOL_GATES['Melbourne High School'] = {
    'lat'   : -37.8276, 'lon': 144.9730,
    'suburb': 'South Yarra', 'lga': 'Melbourne',
}
\end{lstlisting}

  Also add the LGA's SA2 codes to \texttt{DAREBIN\_SA2\_CODES} (rename the variable
  or parameterise it) so the SEIFA download is filtered correctly for the new LGA.

\item \textbf{Add a demographics row to \texttt{demographics\_darebin.csv}.}
  Obtain Census 2021 data for the new suburb from ABS TableBuilder and add a row
  with \texttt{Suburb}, \texttt{Total population}, \texttt{Median weekly household
  income (\$)}, \texttt{\% households no car}, \texttt{\% public transport to work},
  \texttt{\% working full time (35hrs+)}, \texttt{\% children aged 5-17}.

\end{enumerate}

No other files need to change.  \texttt{run\_all.py} will automatically include the
new school in all ten pipeline steps.

\section{config.py as the Single Source of Truth}

Table~\ref{tab:config-roles} lists every configuration object in \texttt{config.py} and
how it propagates through the pipeline.  When extending the project, change \emph{this
file first} and let the downstream scripts pick up the change automatically.

\begin{table}[ht]
\centering
\caption{Key configuration objects in \texttt{config.py} and their downstream consumers.}
\label{tab:config-roles}
\begin{tabular}{lL{4cm}L{5cm}}
\toprule
Object & Role & Consumed by \\
\midrule
\texttt{SCHOOL\_GATES}         & Full name $\to$ lat/lon/addr & \texttt{seifa\_analysis.py}, \texttt{crash\_analysis.py} \\
\texttt{SCHOOL\_SHORT\_NAMES}  & Full name $\to$ short label & \texttt{main.py}, \texttt{feature\_engineering.py} \\
\texttt{SCHOOL\_GATES\_BY\_SHORT} & Short name $\to$ lat/lon & \texttt{spatial\_features.py}, \texttt{scenario\_analyzer.py} \\
\texttt{HS\_INDICATORS}        & Ordered (code, label) pairs & \texttt{main.py} radar chart, \texttt{config.py} HS\_CODES \\
\texttt{CIS\_MAP}              & Cycling infra response $\to$ points & \texttt{src/scoring/cys.py} \\
\texttt{COLUMN\_RENAME}        & Raw CSV header $\to$ internal name & \texttt{src/data/loader.py} \\
\texttt{CSV\_FILE}             & Path to \texttt{school\_data.csv} & \texttt{main.py}, \texttt{src/data/loader.py} \\
\texttt{OUT\_DIR}              & Path to \texttt{outputs/} & every pipeline script \\
\bottomrule
\end{tabular}
\end{table}

\section{Scaling to a Full LGA}

To assess every school in a council area (e.g.\ all 25 schools in Moreland City
Council), the workflow is:

\begin{enumerate}
  \item Run the five-step checklist for each school.
  \item Update \texttt{DAREBIN\_SA2\_CODES} in \texttt{seifa\_analysis.py} to include
    the target LGA's SA2 codes (find these from the ABS 2021 geographic correspondence
    tables at \texttt{abs.gov.au/geography}).
  \item Update \texttt{DAREBIN\_SUBURBS} list to include all suburbs in the new LGA
    (used as a fallback filter if the LGA name column is absent from the SEIFA download).
  \item Run \texttt{python run\_all.py --force} to regenerate all outputs for all schools.
\end{enumerate}

At 25 schools the OSM data extraction (Step~2) will take 2--5 minutes per school;
crash data processing (Step~1) is constant-time because it processes the statewide
file.  Expect a full LGA run to take 45--90 minutes on a standard laptop.

\section{Common Errors and Solutions}

\begin{center}
\begin{tabular}{L{4cm}L{9cm}}
\toprule
Error & Solution \\
\midrule
\texttt{KeyError: 'Reservoir HS'} in \texttt{environmental\_features.py} &
  Add the school's short name to \texttt{SCHOOL\_SUBURB\_MAP}. \\
\texttt{AssertionError: Run spatial\_features.py first} &
  \texttt{outputs/spatial\_features.csv} is missing. Run \texttt{python run\_all.py --from 2}.  \\
\texttt{FileNotFoundError: hs\_predictor.pkl} in \texttt{scenario\_analyzer.py} &
  Run \texttt{python ml\_model.py} (Step~6) first; it creates the \texttt{.pkl} bundle. \\
\texttt{ValueError: School 'X' not found in feature matrix} &
  The new school was added to \texttt{config.py} but not to \texttt{school\_data.csv}
  (no survey rows for that school). \\
\texttt{dayfirst=False} parsing dates as MM/DD/YYYY &
  Ensure \texttt{dayfirst=True} is passed to every \texttt{pd.to\_datetime()} call
  that reads a VicRoads or ABS date column. \\
\texttt{OSM timeout in osmnx.graph\_from\_point()} &
  Add \texttt{ox.settings.timeout = 300} before the osmnx call, or retry after a few
  minutes. The Overpass API has rate limits. \\
\texttt{ABS download failed (SEIFA)} &
  The script falls back to curated data automatically. Check that
  \texttt{download\_ok = False} is logged, then verify the fallback IRSD values
  match the current Census year. \\
\bottomrule
\end{tabular}
\end{center}

% ============================================================
\chapter{Architecture, Limitations, and Lessons Learned}
% ============================================================

\section{What the Architecture Got Right}

Looking back across 21 chapters and ten pipeline steps, several design decisions proved
their worth during the pilot study.

\textbf{Graceful degradation over hard failures.}  Every optional output (interactive
map, heatmap, AQI download, crime download) is wrapped in try/except and falls back
silently.  The core output --- \texttt{hs\_scores.csv} and
\texttt{recommendations.csv} --- is always produced.  This made it possible to run the
pipeline on a university laptop without a stable internet connection and still produce
actionable results for the Regen Melbourne report.

\textbf{Single source of truth for school names.}  \texttt{FULL\_TO\_SHORT} and
\texttt{SCHOOL\_SHORT\_NAMES} dicts live in one place (\texttt{config.py} or in the
individual scripts that need them).  Because all join keys flow through these dicts,
adding a new school never silently breaks an existing join.

\textbf{Open-data-only feature engineering.}  The ML feature matrix (Chapter~19) is
derived entirely from OSM, EPA AirWatch, CSA Victoria, and VicRoads datasets --- all
freely available without registration.  A developer replicating the project
on a different continent could substitute analogous open datasets (OpenStreetMap
remains constant; the others have equivalents in most OECD countries).

\textbf{The delta method for scenario modelling.}  Applying the model's predicted delta
to the actual baseline scores (rather than predicting scenario scores directly) avoids
the most common failure mode of small-$n$ regression: predicting values wildly outside
the training range.  The scenario scores are always anchored to a real, measured
baseline.

\textbf{Skip-if-exists in \texttt{run\_all.py}.}  Incremental re-runs are fast.
Re-running only the one step that changed (e.g.\ after a new observation round)
takes seconds rather than re-running the full 45-minute pipeline.

\section{Known Limitations}

\subsection{Leave-One-Out Cross-Validation with $n = 3$}

LOO-CV with three schools is mathematically valid but statistically meaningless.  With
$n = 3$, LOO-CV trains on two schools and tests on one --- effectively fitting a line
through two points.  The reported RMSE is therefore a measure of how well the model
interpolates between the two training schools, not of generalisation.

This limitation was acknowledged in the project and is the primary reason the scenario
engine uses the delta method (model-predicted change) rather than absolute predictions.
Extending the pipeline to 20+ schools would make LOO-CV statistically meaningful.

\subsection{OSM Data Completeness}

OpenStreetMap is a volunteer-contributed dataset.  In inner Melbourne, OSM coverage is
excellent for roads, footpaths, and amenities.  However:

\begin{itemize}
  \item \textbf{Tree count} (\texttt{tree\_count\_100m}) is systematically
    under-represented --- most street trees in Melbourne are in council datasets, not OSM.
    The tree count feature should be treated as a lower bound.
  \item \textbf{Shelter count} (\texttt{shelter\_count\_200m}) is sparse --- bus shelters
    are inconsistently tagged in OSM.
  \item \textbf{Protected cycling length} can include recreational paths that are not
    practical school routes.
\end{itemize}

For production use, OSM features for tree count and shelter should be supplemented with
council GIS datasets.

\subsection{Three-School Sample and the Equity Correlation}

The Pearson $r \approx 0.84$ between SEIFA decile and overall HS score is based on three
data points.  With $n = 3$, any monotonic relationship produces a high Pearson $r$ ---
the result is directionally correct but cannot be interpreted as a statistically
significant finding at conventional $\alpha = 0.05$.

The equity hypothesis is strongly supported by the qualitative evidence (the most
disadvantaged school has the worst cycling infrastructure, the highest crash density
during school hours, and the highest proportion of car-free households), but the
numerical correlation requires replication at larger $n$ before it should be cited
in policy documents.

\subsection{Hardcoded Fallback Data}

The pipeline ships with hardcoded 2023 suburb averages for PM2.5 and crime rates
(Chapters~17).  These are accurate as of the Census 2021 / financial year 2023--24
reference periods, but they will drift as conditions change.  The fallback values should
be reviewed annually and updated to the most recent available reference year.

\section{Design Decisions Worth Documenting}

\textbf{Why vectorised Haversine rather than PostGIS?}  Nearest-school assignment
(\texttt{crash\_analysis.py}) uses a NumPy vectorised Haversine loop rather than a
PostGIS \texttt{ST\_DWithin} query.  This avoids a PostgreSQL/PostGIS dependency,
making the pipeline installable with a single \texttt{pip install -r requirements.txt}.
For $n <$ 100,000 crash records the performance is adequate ($<$1 second).

\textbf{Why \texttt{pd.DataFrame.iterrows()} in \texttt{main.py} for HS scoring?}
The scoring functions (HS1--HS10) accept a single row dict and return a float.  A
row-wise loop is the simplest correct implementation for a dataset of $<$1,000 survey
observations.  Vectorising the scoring functions would require rewriting all ten modules
--- a premature optimisation for a small dataset.

\textbf{Why Ridge regression?}  Ridge is the simplest regularised regressor and requires
no hyperparameter search (only $\alpha$, tuned by LOO-CV).  Given $n = 3$, any more
complex model would overfit trivially.  Ridge imposes no structural assumptions beyond
linearity and L2 regularisation, which is appropriate when feature engineering is the
primary modelling effort.

\section{Future Technical Directions}

The project identified several concrete extensions that would substantially improve
the pipeline's scientific and practical value:

\begin{description}
  \item[Spatial SEIFA join.]  Replace the suburb-name matching in
    \texttt{seifa\_analysis.py} with a GeoPandas spatial join of ABS SA1 polygons to
    school catchment buffers.  This would give a precise, population-weighted IRSD for
    each school's actual 800\,m walking catchment rather than its postcode suburb.

  \item[GTFS public transport frequency.]  Integrate GTFS timetable data from
    Public Transport Victoria to compute the actual PT headway at each school gate stop.
    This would replace the proxy measure (\texttt{pt\_stops\_400m} count) with a
    meaningful accessibility metric for HS6.

  \item[Repeat-observation study design.]  Run a second round of field surveys after a
    funded intervention (e.g.\ the Plenty Road school zone) is installed.  Comparing
    pre- and post-intervention HS scores against the scenario engine's predicted delta
    would validate (or calibrate) the Ridge model and provide a defensible before/after
    metric for grant acquittals.

  \item[GitHub Actions pipeline refresh.]  Add a CI workflow triggered on push to
    \texttt{main} that runs \texttt{python run\_all.py} and \texttt{python prepare\_data.py}
    and commits the updated \texttt{docs/data/data.json}.  This would keep the GitHub
    Pages dashboard in sync with the latest data automatically.

  \item[Multi-LGA comparison.]  Run the pipeline on schools in Merri-bek (Moreland),
    Yarra, and Darebin councils to test whether the equity finding (disadvantaged
    catchments have lower HS scores) holds at a city scale.  A 60-school sample would
    produce a statistically powered Pearson $r$ and allow cross-council benchmarking.
\end{description}

% ============================================================
\appendix

% ============================================================
\chapter{Complete HS Scoring Formula Reference}
% ============================================================

This appendix documents all ten Healthy Streets indicator formulae as implemented in
\texttt{src/scoring/hs1.py} through \texttt{src/scoring/hs10.py}.  Each function
accepts a scored DataFrame row (post \texttt{COLUMN\_RENAME}) and returns a float in
$[0, 10]$.  The source of truth is always the Python module; this reference is provided
for audit and peer-review purposes.

\section*{HS1 --- Pedestrians from All Walks of Life (Footpath Accessibility)}

Points accumulate from five components (max 10.0 before capping):

\begin{center}
\begin{tabular}{lL{7cm}c}
\toprule
Column & Condition & Points \\
\midrule
\texttt{Footpath\_present} & Both sides of the road & +3.0 \\
                            & One side only           & +2.0 \\
                            & Partial / broken        & +1.0 \\
\texttt{Footpath\_width}   & $\ge 2.0$\,m            & +2.0 \\
                            & $\ge 1.5$\,m            & +1.5 \\
                            & $\ge 1.2$\,m            & +1.0 \\
                            & $\ge 1.0$\,m            & +0.5 \\
\texttt{Continuity}         & 100\,\% continuous      & +2.0 \\
                            & 75--99\,\%              & +1.5 \\
                            & 50--74\,\%              & +1.0 \\
\texttt{FP\_condition}      & (condition / 5.0) $\times$ 2.0 & +0.0 to +2.0 \\
\texttt{Kerb\_ramps}        & All nearby intersections & +0.5 \\
                            & Some intersections      & +0.25 \\
\texttt{Obstructions}       & No obstructions         & +0.5 \\
                            & Minor obstructions      & +0.25 \\
\bottomrule
\end{tabular}
\end{center}

\section*{HS2 --- Easy to Cross}

\begin{center}
\begin{tabular}{lL{7cm}c}
\toprule
Column & Condition & Points \\
\midrule
\texttt{Crossing\_present} & Signalised / traffic light & +4.0 \\
                            & Raised crossing            & +3.0 \\
                            & Zebra / marked             & +2.5 \\
                            & Refuge island              & +2.0 \\
                            & Informal / unmarked        & +0.5 \\
\texttt{Crossing\_dist}    & $\le 30$\,m from gate      & +2.0 \\
                            & $\le 75$\,m               & +1.5 \\
                            & $\le 150$\,m              & +1.0 \\
                            & $\le 250$\,m              & +0.5 \\
\texttt{Visibility}         & (visibility / 5.0) $\times$ 2.0 & +0.0 to +2.0 \\
\texttt{Tactile}            & Both sides                 & +1.0 \\
                            & One side                  & +0.5 \\
\texttt{Signal}             & Countdown timer present   & +1.0 \\
                            & Signal present (no timer) & +0.5 \\
\bottomrule
\end{tabular}
\end{center}

\section*{HS5 --- Not Too Noisy / Traffic Stress}

HS5 starts at 10.0 and deducts points based on traffic stress indicators:

\begin{center}
\begin{tabular}{lL{7cm}c}
\toprule
Column & Condition & Deduction \\
\midrule
\texttt{Traffic\_volume} & Very high / major arterial & $-4.0$ \\
                          & High                       & $-3.0$ \\
                          & Moderate                   & $-1.5$ \\
\texttt{Speed\_limit}    & $\ge 70$\,km/h             & $-2.0$ \\
                          & $\ge 60$\,km/h             & $-1.0$ \\
                          & $\ge 50$\,km/h             & $-0.5$ \\
\texttt{Heavy\_vehicles}  & Frequent                   & $-2.0$ \\
                          & Occasional                 & $-1.0$ \\
\texttt{Lanes}            & 3 or more lanes            & $-1.0$ \\
                          & 2 lanes                    & $-0.5$ \\
\bottomrule
\end{tabular}
\end{center}

\section*{HS Indicators Summary --- All Ten}

\begin{center}
\begin{tabular}{llL{8cm}}
\toprule
Code & Max strategy & Key inputs \\
\midrule
HS1  & Additive (cap 10) & Footpath presence, width, continuity, condition, kerb ramps, obstructions \\
HS2  & Additive (cap 10) & Crossing type, distance from gate, visibility, tactile pavers, signal/countdown \\
HS3  & Additive (cap 10) & Tree count, shelter presence, green space coverage, seating coverage \\
HS4  & Additive (cap 10) & Bench count, park proximity, seating along walking route \\
HS5  & Subtractive (start 10) & Traffic volume, speed limit, heavy vehicles, lane count \\
HS6  & Additive (cap 10) & Cycling infrastructure type (CIS), PT stops within 400\,m \\
HS7  & External (AQI/crime) & Crime rate per 100k (CSA) + street lighting score \\
HS8  & Additive (cap 10) & Amenity count, cafe count, activity diversity within 400\,m \\
HS9  & Additive/subtractive & School zone present, traffic calming, parking conflicts \\
HS10 & External (AQI) & PM2.5 concentration mapped to five breakpoints (US EPA scale) \\
\bottomrule
\end{tabular}
\end{center}

\begin{keybox}{Severity thresholds (src/scoring/severity.py)}
\textbf{Major:} HS2 $<$ 4.0 \emph{or} HS1 $<$ 4.0 \emph{or} HS5 $<$ 3.0. \\
\textbf{Moderate:} 2+ indicators $<$ 6.0 \emph{or} HS\_overall $<$ 5.0. \\
\textbf{Minor:} all other cases.
The Major check runs first; if any crisis indicator is triggered, Moderate is never tested.
\end{keybox}

% ============================================================
\chapter{Recommendation Rules and Intervention Catalogue}
% ============================================================

\section*{The 17 Recommendation Rules}

Table~\ref{tab:rules} lists all 17 rules in \texttt{src/recommendations/rules.py}
in order of rule number.  Each rule fires when its trigger condition is met on a
single observation row and appends a recommendation dict to the row's output.

\begin{table}[ht]
\centering
\caption{All 17 recommendation rules.}
\label{tab:rules}
\small
\begin{tabular}{clL{6cm}l}
\toprule
\# & HS & Trigger condition & Priority \\
\midrule
01 & HS1 & Footpath missing or broken & High \\
02 & HS1 & Footpath width $<$ 1.5\,m & Medium \\
03 & HS2 & No formal pedestrian crossing & High \\
04 & HS2 & Nearest crossing $>$ 150\,m from gate & High \\
05 & HS2 & No tactile ground surface indicators & Medium \\
06 & HS9 & No school zone signage on this street & High \\
07 & HS9 & No traffic calming on 3+ lane road & High \\
08 & HS6 & No cycling infrastructure on school frontage & Low \\
09 & HS7 & Poor or no street lighting & Medium \\
10 & HS3 & Vegetation blocking path or sightlines & Medium \\
11 & HS1 & No kerb ramps at nearby intersections & Medium \\
12 & HS5 & Speed limit $\ge$ 60\,km/h near gate & High \\
13 & HS5 & Frequent heavy vehicles past gate & High \\
14 & HS2 & Crossing visibility rating $<$ 3/5 & Medium \\
15 & HS6 & HS6 score $<$ 4.0 (post-scoring rule) & High \\
16 & HS7 & HS7 score $<$ 5.0 (post-scoring rule) & High \\
17 & HS10 & HS10 score $<$ 5.0 (post-scoring rule) & Medium \\
\bottomrule
\end{tabular}
\end{table}

Rules 1--14 fire on raw survey fields before HS scoring.  Rules 15--17 are
post-scoring rules: they fire on the computed HS indicator values and can trigger
on rows that passed the raw-field checks (e.g.\ cycling infrastructure is present but
the overall active travel score is still low due to other factors).

\section*{The 10 Scenario Interventions}

Table~\ref{tab:interventions} documents all 10 interventions in
\texttt{src/scenarios/interventions.py}.  Each intervention defines a set of feature
deltas applied to the school's ML feature vector before the scenario prediction.

\begin{table}[ht]
\centering
\caption{All 10 scenario interventions with cost and timeframe.}
\label{tab:interventions}
\small
\begin{tabular}{lL{4.5cm}ll}
\toprule
Key & Label & Cost band & Timeframe \\
\midrule
\texttt{pedestrian\_crossing} & Install signalised crossing & Medium (\$80k--\$200k) & $<$ 1 year \\
\texttt{bike\_lane}           & Install protected bike lane & High (\$500k--\$1.5M/km) & 1--3 years \\
\texttt{speed\_reduction}     & Reduce to 40\,km/h school zone & Low (\$5k--\$20k) & $<$ 6 months \\
\texttt{traffic\_calming}     & Speed humps / raised crossing & Medium (\$50k--\$150k) & $<$ 1 year \\
\texttt{footpath}             & Build / extend footpath & Medium (\$100k--\$300k) & $<$ 1 year \\
\texttt{street\_trees}        & Plant street trees & Low (\$20k--\$60k) & $<$ 1 year \\
\texttt{benches}              & Add street seating & Low (\$5k--\$15k) & $<$ 6 months \\
\texttt{pt\_stop}             & Add / improve PT stop & Medium (\$50k--\$200k) & 1--2 years \\
\texttt{shelter}              & Install covered shelter & Low (\$15k--\$40k) & $<$ 6 months \\
\texttt{remove\_arterial}     & Reroute heavy vehicles & High (network-level) & 2--5 years \\
\bottomrule
\end{tabular}
\end{table}

\textbf{Feature deltas by intervention target indicator:}
\begin{itemize}
  \item \texttt{pedestrian\_crossing}: adds +1 to \texttt{crossings\_400m},
    +1 to \texttt{signals\_400m}, +0.15 to \texttt{crossing\_density\_400m}. \hfill
    \emph{Target: HS2}
  \item \texttt{bike\_lane}: adds +12.0 to \texttt{cycle\_pct\_400m},
    +200.0\,m to \texttt{protected\_cycle\_length\_400m}. \hfill \emph{Target: HS6}
  \item \texttt{speed\_reduction}: applies $-10$\,km/h to \texttt{avg\_speed\_400m}
    and \texttt{avg\_speed\_zone}. \hfill \emph{Target: HS5}
  \item \texttt{traffic\_calming}: $-5$\,km/h speed, $-2$ crash count, $-0.10$
    serious/fatal rate. \hfill \emph{Target: HS9}
  \item \texttt{footpath}: +10\,\% to \texttt{footpath\_pct\_200m/400m},
    +500\,m to \texttt{walk\_length\_400m}. \hfill \emph{Target: HS1}
  \item \texttt{street\_trees}: +12 to \texttt{tree\_count\_100m},
    +5\,\% to \texttt{green\_pct\_400m}. \hfill \emph{Target: HS3}
  \item \texttt{benches}: +6 to \texttt{bench\_count\_200m}. \hfill \emph{Target: HS4}
  \item \texttt{pt\_stop}: +1 to \texttt{pt\_stops\_400m}. \hfill \emph{Target: HS6}
  \item \texttt{shelter}: +2 to \texttt{shelter\_count\_200m}. \hfill \emph{Target: HS3}
  \item \texttt{remove\_arterial}: $-10$\,\% arterial, $-3$ high-speed roads,
    $-5$\,km/h speed. \hfill \emph{Target: HS5}
\end{itemize}

% ============================================================
\chapter{Pipeline and Output Quick Reference}
% ============================================================

\section*{The Ten Pipeline Steps}

\begin{center}
\begin{tabular}{clL{4.5cm}L{4cm}}
\toprule
Step & Script & What it does & Key outputs \\
\midrule
1 & \texttt{crash\_analysis.py} & Download VicRoads crash CSVs; filter ped/cyc; assign nearest school via Haversine & \texttt{crash\_data\_statewide.csv}, \texttt{crash\_darebin.csv} \\
2 & \texttt{spatial\_features.py} & Extract OSM features at 100/200/400\,m buffers via osmnx & \texttt{spatial\_features.csv}, \texttt{networks.gpkg} \\
3 & \texttt{environmental\_features.py} & Fetch EPA AQI (3-tier fallback) and CSA crime rate & \texttt{environmental\_features.csv} \\
4 & \texttt{main.py} & Score HS1--HS10; generate charts; build interactive map & \texttt{hs\_scores.csv}, \texttt{recommendations.csv}, 3 PNGs, 2 HTMLs \\
5 & \texttt{feature\_engineering.py} & Merge spatial + env + crash into ML feature matrix & \texttt{ml\_school\_features.csv} \\
6 & \texttt{ml\_model.py} & Train Ridge regressor; LOO-CV evaluation & \texttt{hs\_predictor.pkl}, \texttt{ml\_evaluation.csv} \\
7 & \texttt{seifa\_analysis.py} & Download ABS SEIFA 2021; compute catchment IRSD & \texttt{seifa\_darebin.csv}, \texttt{seifa\_darebin\_sa1.csv} \\
8 & \texttt{equity\_analysis.py} & Join SEIFA + HS scores; compute Pearson $r$ & \texttt{chart\_equity\_seifa.png} \\
9 & \texttt{crash\_trend\_analysis.py} & Year-on-year crash trends; school-hours breakdown & \texttt{chart\_crash\_trends.png} \\
10 & \texttt{demographics\_chart.py} & ABS Census 2021 demographics chart & \texttt{chart4\_demographics.png} \\
\bottomrule
\end{tabular}
\end{center}

\section*{All Output Files}

\begin{center}
\begin{tabular}{lL{8.5cm}}
\toprule
File & Description \\
\midrule
\texttt{crash\_data\_statewide.csv}     & All ped/cyclist crashes in Victoria 2018--2023, with nearest-school assignment \\
\texttt{crash\_darebin.csv}             & Subset of statewide data for Darebin LGA only \\
\texttt{spatial\_features.csv}          & OSM-derived features per school at 200/400\,m buffers \\
\texttt{networks.gpkg}                  & Walk and cycling network GeoPackage for QGIS import \\
\texttt{environmental\_features.csv}    & PM2.5, AQI category, HS10 score, crime rate, HS7 pts per school \\
\texttt{hs\_scores.csv}                 & School-level HS1--HS10 means + overall score + severity \\
\texttt{recommendations.csv}            & All triggered recommendations, ranked by priority \\
\texttt{chart1\_safety\_scores.png}     & Grouped bar chart: FAS, CSS, EEI, CIS, CYS per school \\
\texttt{chart2\_hazard\_severity.png}   & Stacked bar chart: Major / Moderate / Minor counts per school \\
\texttt{chart3\_score\_breakdown.png}   & Per-school indicator breakdown with severity badge \\
\texttt{map\_interactive.html}          & Folium hazard map (open in browser) \\
\texttt{map\_heatmap.html}              & Folium KDE heatmap (open in browser) \\
\texttt{heatmap.png}                    & KDE heatmap PNG for QGIS import \\
\texttt{ml\_school\_features.csv}       & Feature matrix for ML training and scenario inference \\
\texttt{hs\_predictor.pkl}              & Trained Ridge model bundle (pipeline + feature list) \\
\texttt{ml\_evaluation.csv}             & LOO-CV RMSE per indicator \\
\texttt{seifa\_darebin.csv}             & School-level SEIFA IRSD scores and disadvantage labels \\
\texttt{seifa\_darebin\_sa1.csv}        & Full SA1-level Darebin SEIFA data \\
\texttt{chart\_equity\_seifa.png}       & Three-panel equity chart (SEIFA vs HS scores) \\
\texttt{chart\_crash\_trends.png}       & Four-panel crash trend chart (2021--2025) \\
\texttt{chart4\_demographics.png}       & Four-panel ABS demographics chart \\
\texttt{docs/data/data.json}            & Dashboard data bundle for GitHub Pages \\
\bottomrule
\end{tabular}
\end{center}

\section*{Key Configuration Variables}

\begin{center}
\begin{tabular}{lL{8.5cm}}
\toprule
Variable & Description \\
\midrule
\texttt{CSV\_FILE}     & Path to \texttt{school\_data.csv} (field survey input) \\
\texttt{DEMO\_FILE}    & Path to \texttt{demographics\_darebin.csv} (static Census data) \\
\texttt{OUT\_DIR}      & Path to the \texttt{outputs/} directory \\
\texttt{SCHOOL\_GATES} & Dict: full school name $\to$ lat, lon, street address \\
\texttt{SCHOOL\_SHORT\_NAMES} & Dict: full name $\to$ short display name \\
\texttt{SCHOOL\_GATES\_BY\_SHORT} & Dict: short name $\to$ lat, lon (for scripts keyed by short name) \\
\texttt{HS\_INDICATORS} & Ordered list of (code, label) pairs defining HS1--HS10 \\
\texttt{CIS\_MAP}      & Dict: cycling infrastructure response $\to$ CIS points (1.0--9.0) \\
\texttt{COLUMN\_RENAME} & Dict: raw Google Form header $\to$ internal column name \\
\bottomrule
\end{tabular}
\end{center}

\section*{Running the Pipeline}

\begin{lstlisting}[style=bash, caption={Full pipeline run from a clean checkout.}]
# 1. Create environment and install dependencies
python -m venv venv
source venv/bin/activate          # Windows: venv\Scripts\activate
pip install -r requirements.txt

# 2. Run all 10 steps (skip-if-exists logic prevents re-runs)
python run_all.py

# 3. Force re-run from step 4 onwards (e.g. after new survey data)
python run_all.py --from 4 --force

# 4. Bundle outputs for the dashboard and deploy
python prepare_data.py
git add docs/data/data.json
git commit -m "update dashboard data"
git push origin main
# GitHub Pages serves the updated dashboard within ~60 seconds
\end{lstlisting}

% -- Back matter ----------------------------------------------
\backmatter

\chapter*{What's Next}
\addcontentsline{toc}{chapter}{What's Next}

This book has now covered the complete technical implementation of the 300,000 Streets
pilot study across twenty-one chapters in six parts, together with three reference
appendices:

\medskip
\begin{description}
  \item[Part 1 (Chapters 1--4).] Project context and the Healthy Streets framework;
    development environment setup; repository architecture and configuration.
  \item[Part 2 (Chapters 5--8).] Study area and school selection; field survey methodology
    and the Google Forms instrument; OSM data collection with osmnx; government open
    datasets (VicRoads, ABS SEIFA, EPA, Crime Statistics).
  \item[Part 3 (Chapters 9--12).] The HS scoring engine (all ten indicator formulae);
    severity classification and the recommendation rule engine; machine learning with
    Ridge regression and LOO-CV; scenario analysis and the delta method.
  \item[Part 4 (Chapters 13--15).] Crash trend analysis (VicRoads 2021--2025);
    equity analysis (SEIFA disadvantage versus HS safety scores); dashboard and
    deployment pipeline (\texttt{prepare\_data.py}, Leaflet, GitHub Pages).
  \item[Part 5 (Chapters 16--18).] Crash data download and vectorised spatial join
    (\texttt{crash\_analysis.py}); environmental features: EPA AirWatch AQI and
    CSA Victoria crime rate with two-tier fallback (\texttt{environmental\_features.py});
    core pipeline orchestration and visualisation (\texttt{main.py},
    \texttt{src/visualisation/}).
  \item[Part 6 (Chapters 19--21).] Feature engineering --- the 26-column open-data
    feature matrix (\texttt{feature\_engineering.py}); SEIFA 2021 disadvantage
    analysis (\texttt{seifa\_analysis.py}); ABS Census 2021 demographics chart
    (\texttt{demographics\_chart.py}).
  \item[Part 7 (Chapters 22--24).] Pilot study results and policy implications;
    how to extend the pipeline to new schools and LGAs; architecture retrospective,
    known limitations, and future technical directions.
  \item[Appendices A--C.] Complete HS scoring formula reference; all 17 recommendation
    rules and 10 scenario interventions; pipeline step table, output file catalogue,
    and configuration variable reference.
\end{description}

\medskip
A developer who has followed this book from Chapter~1 can now:

\begin{enumerate}
  \item Clone the repository and run \texttt{pip install -r requirements.txt}.
  \item Run \texttt{python run\_all.py} to reproduce every one of the 21 output files
    from a clean checkout --- from crash data download through SEIFA analysis, ML
    training, and the demographics chart.
  \item Run \texttt{python prepare\_data.py} to bundle all outputs into
    \texttt{docs/data/data.json}.
  \item Push to \texttt{main} and have a live GitHub Pages dashboard within 60 seconds.
  \item Add a new school to the five configuration points described in Chapter~23 and
    re-run the pipeline without changing a single line of analysis code.
\end{enumerate}

\medskip\noindent
The 300,000 Streets pilot demonstrates that a rigorous, evidence-based assessment of
school street safety can be built entirely from open data --- OSM, ABS, VicRoads, EPA,
and Crime Statistics --- without a budget for proprietary GIS software or bespoke
sensor infrastructure.  The pipeline is small enough to run on a laptop, transparent
enough to audit, and modular enough to extend.  The ten-step orchestrator and the
GitHub Pages dashboard are not the end product; they are the scaffolding for a much
larger study.

\medskip\noindent
What the pilot proved above all is that the equity hypothesis is real and measurable.
The students who walk to Reservoir High School along Plenty Road every morning are
doing so on a street that scores 2.1 out of 10 for cycling safety, 3.8 out of 10 for
shade and shelter, and 4.2 out of 10 for footpath accessibility --- while living in a
catchment that ABS ranks in the bottom 40\,\% of the country for socioeconomic
advantage.  That is not a coincidence.  It is a pattern, and this pipeline exists to
make that pattern visible, measurable, and actionable.

\medskip\noindent
The 300,000 streets are waiting.

\bigskip
\begin{flushright}
\textit{300,000 Streets of Melbourne --- Regen Melbourne $\times$ RMIT University}\\
\textit{City of Darebin Pilot, 2024}
\end{flushright}

\end{document}
